% Paper 32: The GeoVac Spectral Triple — A Synthesis
% GeoVac Project, May 2026
% v1.0 -- Synthesis paper. Constructs the spectral triple (A, H, D)
% explicitly on the Fock-projected S^3 graph, checks axioms, places
% the project inside the Marcolli-van Suijlekom gauge-network lineage,
% and reads Papers 25, 28, 30, 31 as structural sub-sectors.

\documentclass[11pt]{article}
\usepackage{amsmath,amssymb,amsthm}
\usepackage{booktabs}
\usepackage{array}
\usepackage[margin=1in]{geometry}

% Shorthand macros used throughout the paper.
% Added 2026-05-16 (Sprint L2-B/C/D paper-edits round): restoring
% definitions that were referenced but undefined in pre-existing content
% (§VIII.D frontier-of-field paragraphs).  Definitions are conservative
% (calligraphic / mathematical-italic) and consistent with the usage
% contexts.
\providecommand{\Tcal}{\mathcal{T}}
\providecommand{\Op}{\mathcal{O}}
\providecommand{\sthree}{S^3}
\providecommand{\Q}{\mathbb{Q}}
\providecommand{\Z}{\mathbb{Z}}
\providecommand{\nmaxa}{n_{\max, a}}
\providecommand{\nmaxb}{n_{\max, b}}
\providecommand{\Tr}{\mathrm{Tr}}

\newtheorem{theorem}{Theorem}
\newtheorem{proposition}{Proposition}
\newtheorem{corollary}{Corollary}
\newtheorem{observation}{Observation}
\newtheorem{definition}{Definition}
\newtheorem{remark}{Remark}

\title{The GeoVac Spectral Triple:\\
A Synthesis Paper Locking the Framework into the\\
Marcolli--van~Suijlekom Gauge-Network Lineage}
\author{GeoVac Project}
\date{May 2026}

\begin{document}
\maketitle

\begin{abstract}
The GeoVac framework has, over Papers~0--31, accumulated a body of
structural results---the $S^3$ Fock-projection equivalence
(Paper~7~\cite{paper7}), the angular sparsity theorem
(Paper~22~\cite{paper22}), the Hopf-graph lattice-gauge structure
(Paper~25~\cite{paper25}), QED on $S^3$ (Paper~28~\cite{paper28}),
the Bargmann--Segal $S^5$ HO lattice (Paper~24~\cite{paper24}), the
$\mathrm{SU}(2)$ Wilson gauge (Paper~30~\cite{paper30}), and the
universal/Coulomb partition (Paper~31~\cite{paper31})---each
self-contained but together suggesting one parent structure.  This
paper names and constructs that structure: a finite spectral triple
$(\mathcal{A}_{\mathrm{GV}}, \mathcal{H}_{\mathrm{GV}},
D_{\mathrm{GV}})$ on the Fock-projected $S^3$ graph, with
$\mathcal{A}_{\mathrm{GV}}$ the algebra of complex functions on graph
nodes, $\mathcal{H}_{\mathrm{GV}}$ the Camporesi--Higuchi spinor space
restricted to the truncation, and $D_{\mathrm{GV}}$ a graph Dirac
operator whose continuum limit is the Camporesi--Higuchi
operator~\cite{camporesi_higuchi1996} on unit~$S^3$.  This is not a
new derivation; it is a structural map.  Three results follow.
(1)~\emph{Lineage.}~The construction sits inside the published
Marcolli--van~Suijlekom gauge-network framework~\cite{marcolli_vs2014}
(with the Perez-Sanchez correction~\cite{perez_sanchez2025} that the
continuum limit is Yang--Mills without Higgs), and inside the
Connes--van~Suijlekom spectral-truncation
program~\cite{connes_vs2021,hekkelman2022,hekkelman_mcdonald2024,vansuijlekom2024_ksystems}.  Every
structural ingredient of the GeoVac construction has published
precedent; what is not duplicated is the $\alpha$ prediction at
$8.8 \times 10^{-8}$ (Paper~2~\cite{paper2}) from a discrete spectral
construction.
(2)~\emph{Sub-sector identification.}  Paper~25's Hodge-1 Laplacian is
the Cartan-subalgebra (maximal-torus) projection of an
almost-commutative extension of the present triple by~$M_2(\mathbb{C})$;
Paper~30's $\mathrm{SU}(2)$ Wilson action is the full non-abelian
slice of the same extension.  Paper~28's QED on~$S^3$ is the
spectral-action computation on $D_{\mathrm{GV}}$; Paper~31's
universal/Coulomb partition is the $\mathcal{A}/D$ split of the
present triple.  Papers~25, 28, 30, 31 are not four parallel
results---they are four projections of the single triple constructed
here.
(3)~\emph{Axiom audit.}  $\mathcal{A}_{\mathrm{GV}}$ is finite,
involutive, faithfully represented on $\mathcal{H}_{\mathrm{GV}}$;
$D_{\mathrm{GV}}$ is self-adjoint with bounded $[D_{\mathrm{GV}}, a]$
for $a \in \mathcal{A}_{\mathrm{GV}}$; the order-one condition is
trivially satisfied by the abelian core; a real structure
$J_{\mathrm{GV}}$ exists on the Dirac sector with
$J_{\mathrm{GV}}^2 = -\mathbb{1}$ (the quaternionic/pseudoreal sign
for the 2-dimensional Dirac spinor representation of
$\mathrm{Spin}(3) = \mathrm{SU}(2)$) and $J_{\mathrm{GV}} D_{\mathrm{GV}}
= +D_{\mathrm{GV}} J_{\mathrm{GV}}$, matching the standard signs
$(\epsilon, \epsilon') = (-1, +1)$ for $\mathrm{KO}$-dimension~$3$
mod~$8$ of~$S^3$~\cite{connes1995}; a $\mathbb{Z}_2$-grading
$\gamma_{\mathrm{GV}}$ is absent at the Dirac level (correctly:
$S^3$ is odd-dimensional, the spectral triple is \emph{odd}).  The four-way $S^3$ coincidence---Fock projection image,
Hopf bundle base, Camporesi--Higuchi spin carrier,
$\mathrm{SU}(2)$ gauge manifold---is the statement that the manifold
of $\mathcal{A}_{\mathrm{GV}}$ in the continuum limit, the spinor
bundle of $D_{\mathrm{GV}}$, and the gauge group of natural
fluctuations of $D_{\mathrm{GV}}$ all coincide on the unique rank-1
non-abelian compact simply-connected Lie group.

\noindent\textbf{Scope.}  This paper does not derive new physics.  It
synthesizes existing GeoVac results into the axiomatic spectral-triple
language.  The synthesis lifts each result from a free-standing
observation to a sub-sector of a single named structure, and makes
explicit which axioms hold rigorously, which hold up to the standard
finite-truncation caveats, and which remain open---specifically
whether the GeoVac Dirac graph operator carries a real structure
matching the topological $\mathrm{KO}$-dimension of $S^3$ exactly at
finite $n_{\max}$ or only in the continuum limit.
\end{abstract}

% ===================================================================
\section{Introduction}
\label{sec:intro}
% ===================================================================

\subsection{Why a spectral triple, and why now}

Across thirty-one papers the GeoVac framework has carried four
distinct strands.  Strand~1 is the Fock projection
$p_0 = \sqrt{-2E_n}$ that maps the hydrogenic Hamiltonian onto the
free Laplacian on the unit three-sphere, with eigenvalues
$\lambda_n = -(n^2 - 1)$ on graph nodes labeled by quantum-number
triples $(n, l, m_l)$ (Paper~7).  Strand~2 is the Hopf bundle
$S^1 \to S^3 \to S^2$ that organizes the discrete edge structure into
ladder operators with $U(1)$ phase content (Papers~2, 25).  Strand~3
is the Camporesi--Higuchi Dirac spectrum
$|\lambda_n^{\mathrm{D}}| = n + 3/2$ with degeneracy
$g_n^{\mathrm{D}} = 2(n+1)(n+2)$ that controls QED on $S^3$
(Paper~28).  Strand~4 is the Wilson lattice gauge structure with
$\mathrm{SU}(2)$ link variables on the same graph (Paper~30).

Each strand uses the unit three-sphere in a different role:
\begin{itemize}\setlength{\itemsep}{0pt}
\item[(a)] $S^3$ as the \emph{conformal image} of the hydrogenic
spectrum (Paper~7);
\item[(b)] $S^3$ as the \emph{base} of the bundle carrying the
fine-structure-constant construction (Paper~2);
\item[(c)] $S^3$ as the \emph{spin carrier} of the Camporesi--Higuchi
Dirac operator (Paper~28);
\item[(d)] $S^3$ as the \emph{gauge manifold} of $\mathrm{SU}(2)$
Wilson links (Paper~30).
\end{itemize}
Working hypothesis~WH4 in the project register flags this as the
strangest and most consequential coincidence in the framework: $S^3 =
\mathrm{SU}(2)$ is the unique rank-1 non-abelian compact simply
connected Lie group where (a)--(d) all coincide; in any other
dimension or any other rank, the four roles separate.  WH4 takes the
position that the four strands are not four coincidences but a single
spectral-triple structure viewed in four projections.

This paper makes WH4 explicit.  We give a concrete construction of a
finite spectral triple
$(\mathcal{A}_{\mathrm{GV}}, \mathcal{H}_{\mathrm{GV}},
D_{\mathrm{GV}})$ (Sec.~\ref{sec:construction}); audit the standard
Connes axioms one by one (Sec.~\ref{sec:axiom_audit}); identify
Papers~25, 28, 30, and~31 as structural sub-sectors of the resulting
object (Sec.~\ref{sec:subsectors}); read the Paper~2 combination
formula $K = \pi(B + F - \Delta)$ inside the spectral-action language
(Sec.~\ref{sec:K_decomposition}); and place the Coulomb/HO asymmetry
of Paper~31 in spectral-triple terms (Sec.~\ref{sec:coulomb_ho}).

\subsection{What this paper is and is not}

\paragraph{What it is.}  A synthesis paper.  The construction is
explicit, finite, and reproducible from existing GeoVac modules
(\texttt{geovac/lattice.py}, \texttt{geovac/dirac\_matrix\_elements.py},
\texttt{geovac/fock\_graph\_hodge.py}); the lineage citations are to
published mathematical work; the audit of the Connes axioms uses
standard textbook definitions~\cite{connes1995,connes_marcolli2008,%
suijlekom_book2015}.

\paragraph{What it is not.}  A derivation of $\alpha$.  Paper~2 stays
where it is: a structural observation with three independent spectral
homes for the three pieces $B$, $F$, $\Delta$ but no first-principles
derivation of the combination rule.  This paper does not change that
status; it gives the spectral-action language in which the open
question is phrased without changing the conjectural label.

\paragraph{What it is not, second.}  A claim of the
Connes--Chamseddine Standard Model.  The GeoVac construction is a
\emph{rank-1} non-abelian almost-commutative extension; the SM
spectral triple of \cite{chamseddine_connes2010} is rank~$\ge 3$ with
specific finite-algebra content $\mathbb{C} \oplus \mathbb{H} \oplus
M_3(\mathbb{C})$.  GeoVac is not the Standard Model.  It is a finite
spectral triple in the same lineage, with the empirically interesting
property that one of its spectral-action coefficient combinations
$K = \pi(B + F - \Delta)$ satisfies, after subtraction of the leading
Sprint-A correction~$\alpha^2$, a residual
$K - 1/\alpha - \alpha^2 \approx 1.2 \times 10^{-5}$ corresponding to
a post-correction relative residual of $8.8\times 10^{-8}$ that
matches~$\pi^3\alpha^3$ to within 0.25\% (the raw match
$|K - 1/\alpha|/(1/\alpha)$ is $4.77 \times 10^{-7}$).

\paragraph{Internal register.}  This paper makes explicit the
spectral-triple framing already operating informally in the working
hypothesis register (\texttt{CLAUDE.md}~\S 1.7, WH1 ``MODERATE--STRONG'').
Papers continue under the rhetoric rule of \S 1.5: dual-description
framing, no ontological priority claims, lead with computational
results.

% ===================================================================
\section{The Connes Data Recalled}
\label{sec:connes_data}
% ===================================================================

A spectral triple is a triple $(\mathcal{A}, \mathcal{H}, D)$
where~\cite{connes1995}
\begin{itemize}\setlength{\itemsep}{2pt}
\item $\mathcal{A}$ is an involutive (unital) algebra,
\item $\mathcal{H}$ is a Hilbert space carrying a faithful
$*$-representation $\pi: \mathcal{A} \to \mathcal{B}(\mathcal{H})$,
\item $D : \mathcal{H} \to \mathcal{H}$ is an unbounded self-adjoint
operator with compact resolvent, such that the commutators
$[D, \pi(a)]$ are bounded for every $a \in \mathcal{A}$.
\end{itemize}
The triple is called \emph{even} if there is a self-adjoint involution
$\gamma$ commuting with $\pi(\mathcal{A})$ and anticommuting with $D$;
otherwise \emph{odd}.  A \emph{real structure} on the triple is an
antilinear isometry $J: \mathcal{H} \to \mathcal{H}$ with
$J^2 = \epsilon\,\mathbb{1}$, $JD = \epsilon'\,DJ$,
$J\gamma = \epsilon''\,\gamma J$ (the signs $(\epsilon,\epsilon',
\epsilon'')$ are determined by the $\mathrm{KO}$-dimension
mod~$8$~\cite{connes1995}); for the canonical Riemannian spin
manifold of dimension $n$, the $\mathrm{KO}$-dimension is $n$ mod~$8$
and the signs are tabulated.

The two structural axioms of central interest below are:
\begin{itemize}
\item[(\textbf{O1})] \emph{Order one:} $[[D, \pi(a)], \pi(b)^{\circ}]
= 0$ for all $a, b \in \mathcal{A}$, where $b^{\circ} = JbJ^{-1}$ is
the right action induced by the real structure;
\item[(\textbf{R})] \emph{Reality:} $J$ exists with the signs above.
\end{itemize}

For an \emph{almost-commutative} spectral triple
$(\mathcal{A}_M \otimes \mathcal{A}_F, \mathcal{H}_M \otimes
\mathcal{H}_F, D_M \otimes \mathbb{1} + \gamma_M \otimes
D_F)$~\cite{vansuijlekom_book2015,chamseddine_connes2010}, the
manifold sector $\mathcal{A}_M = C^{\infty}(M)$ is replaced in our
finite setting by the algebra of complex-valued functions on a finite
graph that converges to $C^{\infty}(M)$ in the continuum limit.  This
is exactly the move made by Marcolli and van~Suijlekom in their gauge
networks construction~\cite{marcolli_vs2014} and continued by
Perez-Sanchez~\cite{perez_sanchez2024,perez_sanchez2025}: replace the
manifold algebra by an algebra on a discrete vertex set, with edges
carrying connection data.  The GeoVac construction sits in this
lineage by design.

% ===================================================================
\section{The GeoVac Triple: Explicit Construction}
\label{sec:construction}
% ===================================================================

Fix a truncation $n_{\max} \ge 1$.  All objects below are finite.

\subsection{The algebra $\mathcal{A}_{\mathrm{GV}}$}

Let $V_{\mathrm{Fock}}$ denote the set of Fock-graph nodes at
truncation $n_{\max}$, i.e. quantum-number triples $(n, l, m_l)$ with
$1 \le n \le n_{\max}$, $0 \le l \le n-1$, $-l \le m_l \le l$.  The
total node count is $N_{\mathrm{Fock}}(n_{\max}) = \sum_{n=1}^{n_{\max}}
n^2$.

\begin{definition}[Function algebra on the Fock graph]
\label{def:A_GV}
$\mathcal{A}_{\mathrm{GV}}^{(n_{\max})} := \mathbb{C}^{V_{\mathrm{Fock}}}$,
the commutative $*$-algebra of complex-valued functions on
$V_{\mathrm{Fock}}$ under pointwise multiplication, with involution
$f^* (v) = \overline{f(v)}$.
\end{definition}

This is precisely the manifold algebra of the gauge networks of
\cite{marcolli_vs2014} (their $C(V)$, with $V$ the vertex set), with
the GeoVac-specific choice that $V$ is the Fock-projected $S^3$ graph
of Paper~7 rather than a generic graph.  It is finite-dimensional
($\dim_{\mathbb{C}} \mathcal{A}_{\mathrm{GV}} = N_{\mathrm{Fock}}$),
unital ($1$~is the constant function~$1$), and abelian.

\begin{remark}[$\mathcal{A}_{\mathrm{GV}}$ as $C^*$-envelope of an operator system]
\label{rem:operator_system}
In the parallel reading of GeoVac through the spectral-truncation
framework of Connes and van~Suijlekom~\cite{connes_vs2021},
Definition~\ref{def:A_GV} describes
$\mathcal{A}_{\mathrm{GV}}^{(n_{\max})} = \mathbb{C}^{V_{\mathrm{Fock}}}$
as the $C^*$-envelope $C^*(\mathcal{O}_{n_{\max}})$ of the operator
system
\[
   \mathcal{O}_{n_{\max}} := P_{n_{\max}} \, C^\infty(S^3) \, P_{n_{\max}},
\]
where $P_{n_{\max}}$ is the index-truncation projecting onto the
span of the first $n_{\max}$ Fock shells (in the spinor-harmonic
basis of $L^2(S^3)$ adapted to the SO(4) decomposition).  This is
structurally identical to the Connes--van~Suijlekom worked example
on $S^1$ (their $P_n$ onto $\mathrm{span}_{\mathbb{C}}\{e_1, e_2,
\ldots, e_n\}$, \cite[\S 3.1]{connes_vs2021}): a basis-index cutoff,
not a Borel functional-calculus projection
$\chi_{[-\Lambda, \Lambda]}(D)$.  The genuine truncated algebra-image
$\mathcal{O}_{n_{\max}}$ is $*$-closed but \emph{not} multiplicatively
closed; it carries nontrivial off-diagonal matrix elements
\[
   \langle (n,l,m_l) | \, f \, | (n', l', m_l') \rangle
   = \int_{S^3} \overline{Y_{n l m_l}}\, f \, Y_{n' l' m_l'}\, dV
\]
between distinct $n$-shells.  These overlap integrals are exactly
the Gaunt~/~Wigner-3j angular machinery and the hypergeometric
Slater radial integrals already implemented throughout the GeoVac
pipeline (Paper~22 angular sparsity; Paper~14 ERI structure).  The
off-diagonal data therefore already exists in the framework;
$\mathcal{A}_{\mathrm{GV}}$ as written in Definition~\ref{def:A_GV}
is its diagonal restriction.

The two readings agree on diagonal elements and on the spectral
data of $D_{\mathrm{GV}}$; they differ in what is taken as the
algebraic content surviving the truncation.  Connes and
van~Suijlekom note that taking the $C^*$-envelope of $\mathcal{O}_\Lambda$
typically collapses to a ``non-informative full matrix algebra'' --
the algebraic-structure information lives in the operator-system
data thrown away by the envelope construction.  Constructing
$\mathcal{O}_{n_{\max}}$ explicitly in $\mathcal{B}(\mathcal{H}_{\mathrm{GV}})$
and computing its propagation number and Connes distance on
node-evaluation states is the deliverable of the WH1 round-2 sprint
(see \texttt{debug/wh1\_connes\_vs\_mapping\_memo.md}, 2026-05-03).
The present subsection retains the
$\mathcal{A}_{\mathrm{GV}} = \mathbb{C}^{V_{\mathrm{Fock}}}$ definition
because (a) it matches the Marcolli--van~Suijlekom gauge-network
convention used throughout Papers~25 and~30, and (b) the spectral
triple's structural verifications (Theorem~\ref{thm:GV_triple}
below; KO-dimension audit; Hodge structure) are unaffected by
which representative of the equivalence class is taken.
\end{remark}

\begin{proposition}[Propagation number of $\mathcal{O}_{n_{\max}}$]
\label{prop:propagation_number}
Let $\mathcal{O}_{n_{\max}} = P_{n_{\max}} C^\infty(S^3) P_{n_{\max}}$
be the truncated operator system of
Remark~\ref{rem:operator_system}, viewed as a $*$-closed subspace of
$M_N(\mathbb{C})$ with $N = N_{\mathrm{Fock}}(n_{\max})$.  The
propagation number in the sense of Connes \& van~Suijlekom
(\cite{connes_vs2021}, Definition~2.39) satisfies
\begin{equation}
   \mathrm{prop}(\mathcal{O}_{n_{\max}}) = 2
   \qquad \text{for every } n_{\max} \in \{2, 3, 4\},
   \label{eq:prop_two}
\end{equation}
matching the Toeplitz operator system $C(S^1)^{(n)}$
(\cite{connes_vs2021}, Proposition~4.2: $\mathrm{prop}(C(S^1)^{(n)}) =
2$ independently of $n$) verbatim.  Equivalently,
$\dim_{\mathbb{C}}(\mathcal{O}_{n_{\max}}^2) = N^2$ while
$\dim_{\mathbb{C}}(\mathcal{O}_{n_{\max}}) < N^2$ strictly:
\[
\renewcommand{\arraystretch}{1.15}
\begin{array}{c|c|c|c|c}
n_{\max} & N & N^2 & \dim(\mathcal{O}_{n_{\max}}) &
   \dim(\mathcal{O}_{n_{\max}}^2) \\ \hline
2 & 5  & 25  & 14  & 25  \\
3 & 14 & 196 & 55  & 196 \\
4 & 30 & 900 & 140 & 900 \\
\end{array}
\]
The dimension fraction $\dim(\mathcal{O}_{n_{\max}}) / N^2$ falls
$0.560 \to 0.281 \to 0.156$ across the three truncations: the
operator system is a vanishing fraction of its $C^*$-envelope as
$n_{\max}$ grows, but its second power fills the envelope
completely.
\end{proposition}

\begin{proof}[Computational verification]
The full construction and verification is documented in
\texttt{geovac/operator\_system.py} (the \texttt{TruncatedOperatorSystem}
class) and \texttt{tests/test\_operator\_system.py} (24 tests, all
passing; CPU time ${\sim}20$~s).  In summary: a basis of
$\mathcal{O}_{n_{\max}}$ is built from the multiplier matrices
$M_{NLM}$ corresponding to the hyperspherical harmonics
$Y^{(3)}_{NLM}$, with selection rules $|n - n'| + 1 \le N \le n + n' -
1$ (SO(4) triangle), $|l - l'| \le L \le l + l'$ with $l + l' + L$
even (S$^2$ Gaunt parity), $M = m - m'$, and the angular S$^2$ part
evaluated symbolically via \texttt{sympy.physics.wigner.wigner\_3j}.
Linear independence of $\{M_{NLM}\}$ is verified at machine precision;
the propagation number is read off as the smallest $k$ with
$\dim(\mathcal{O}_{n_{\max}}^k) = N^2$.  The result
$\mathrm{prop} = 2$ is robust: it is invariant under the choice of
SO(4) radial-overlap placeholder (verified across three placeholders
including a pure indicator that retains only the support pattern), so
\eqref{eq:prop_two} is a structural invariant of the SO(4) selection
rules at finite cutoff and does not depend on the radial value
convention.  See \texttt{debug/wh1\_round2\_propagation\_memo.md}
(2026-05-03) for the full computational record, including the explicit
witness pair $a = M_{2,1,0}$ at $n_{\max} = 2$ for which $a^2 \notin
\mathcal{O}_{n_{\max}}$ (Frobenius residual $\approx 14.9\%$), with
the deficit concentrated on the top-shell diagonal entries.
\end{proof}

\begin{remark}[Significance of $\mathrm{prop} = 2$]
Connes \& van~Suijlekom prove $\mathrm{prop}(C(S^1)^{(n)}) = 2$ for
the Toeplitz operator system on the circle by showing that products
of two Toeplitz shift generators reach every matrix unit
$e_{ij}$~\cite{connes_vs2021}.  The companion dual operator system
$C(\mathbb{Z})^{(n)}$ has $\mathrm{prop} = \infty$ (their
Proposition~4.3) — so the value $\mathrm{prop} = 2$ is non-trivial
even at this level: it is specific to the Toeplitz construction, not
to operator systems in general.  Eq.~\eqref{eq:prop_two} therefore
states that the Fock-projected $S^3$ truncation is in the same
propagation class as the Toeplitz $S^1$ — one of the simplest worked
examples in the spectral-truncation literature — at the level of a
quantitative Connes--van~Suijlekom invariant, not merely at the
level of the underlying construction.  Together with the
$\mathrm{SO}(4)$ isometry preservation discussed in
Sec.~\ref{sec:axiom_audit} below, this is the strongest
quantitative alignment between GeoVac and the published Connes-vS
framework currently known.

The mechanism producing
$\mathcal{O}_{n_{\max}} \subsetneq M_N(\mathbb{C})$ is the
``raise-then-lower at the top shell escapes the truncation'' boundary
correction: a multiplier $M_{N \ge 2, L, M}$ acting at the top shell
$n = n_{\max}$ would, in the continuum, be balanced by intermediate
sums over $n' > n_{\max}$, but the truncation $P_{n_{\max}}$ kills
these contributions.  This is structurally the same boundary effect
as the pendant-edge theorem of Paper~28~\cite{paper28},
$\Sigma_{\mathrm{graph}}(\mathrm{GS}) = 2(n_{\max} - 1)/n_{\max}$ for
the graph-native ground-state self-energy, viewed from the
operator-system side rather than the QED side.
\end{remark}

\begin{corollary}[Structural specificity of $\mathrm{prop} = 2$ on $S^3$]
\label{cor:structural_specificity}
The propagation value
$\mathrm{prop}(\mathcal{O}_{n_{\max}}) = 2$ in
Proposition~\ref{prop:propagation_number} is structurally specific to
the Connes--van~Suijlekom spectral truncation of $S^3$ and \emph{not}
a generic feature of arbitrary finite-dimensional truncations of $S^3$.
The natural circulant-style comparator on $S^3$ at $N$ points
\[
   A_{\mathrm{circ}, N} := \mathrm{span}\{E_{i, i} : i = 0, \ldots, N - 1\}
   \subset M_N(\mathbb{C}),
\]
i.e., the commutative $C^*$-subalgebra of diagonal matrices in the
``point-evaluation'' basis of $\mathbb{C}^N$, satisfies
\begin{equation}
   \mathrm{prop}(A_{\mathrm{circ}, N}) = \begin{cases}
      1 & \text{(intrinsic envelope:\ $A_{\mathrm{circ}, N}$ is its own
                  $C^*$-algebra)}, \\
      \infty & \text{(ambient envelope:\ $M_N(\mathbb{C})$ is non-abelian;
                       no power of an abelian subalgebra reaches it)}.
   \end{cases}
   \label{eq:prop_circulant}
\end{equation}
Verified at $N \in \{2, 5, 14, 30, 55, 140\}$, matching both
dimension-matching conventions (envelope-match $N(\mathrm{circ}) =
N(\mathrm{GV})$ and operator-system-match $N(\mathrm{circ}) =
\dim(\mathcal{O}_{n_{\max}})$).  Either reading of
\eqref{eq:prop_circulant} is categorically different from the
spectral-truncation value~$2$.
\end{corollary}

\begin{proof}[Computational verification]
The construction lives in \texttt{geovac/circulant\_s3.py} (the
\texttt{CirculantS3Truncation} class) with verification in
\texttt{tests/test\_circulant\_s3.py} (35 tests, all pass; runtime
${\sim}1.4$~s).  Algorithmically, $\mathrm{prop}(A_{\mathrm{circ}, N})$
is computed via the same vec-stack rank algorithm used in the
proof of Proposition~\ref{prop:propagation_number}, applied to the
diagonal-matrix basis of $A_{\mathrm{circ}, N}$.  The intrinsic-envelope
reading exits at $k = 1$ since $\dim(A_{\mathrm{circ}, N}^1) = N =
\dim_{\mathbb{C}}(A_{\mathrm{circ}, N})$ already saturates the
intrinsic $C^*$-envelope.  In the ambient-envelope reading, the
$M_N(\mathbb{C})$ envelope has dimension $N^2$ for $N \ge 2$ but
$\dim(A_{\mathrm{circ}, N}^k) = N$ for every $k$ (commutativity is
preserved under self-products), so no finite $k$ achieves
saturation; $\mathrm{prop} = \infty$ by Connes--van~Suijlekom
Definition~2.39.  See
\texttt{debug/wh1\_round3\_falsification\_memo.md} (2026-05-03) for
the full record, including the verbatim Connes--van~Suijlekom Outlook~\S6
contrast (lines 2611--2618 of arXiv:2004.14115v2) between
spectral and circulant truncations on~$S^1$, of which
Eq.~\eqref{eq:prop_circulant} is the verbatim $S^3$ analog.
\end{proof}

\begin{remark}[Connes distance on $\mathcal{O}_{n_{\max}}$: SDP framework
and structural findings]
\label{rem:connes_distance}
The Connes distance on operator systems
(\cite{connes_vs2021}, Eq.~(2))
\[
   d_D(\phi, \psi) = \sup_{x \in \mathcal{O}} \big\{ |\phi(x) - \psi(x)|
      : \|[D, x]\|_{\mathrm{op}} \le 1 \big\}
\]
restricted to pure node-evaluation states $\phi_v(x) := \langle v | x | v
\rangle$ on the truncated operator system has been computed by
semidefinite programming on $\mathcal{O}_{n_{\max}}$ at $n_{\max} = 2, 3$.
The implementation lives in \texttt{geovac/connes\_distance.py} (a
\texttt{cvxpy} SDP with the operator-norm constraint lifted to the LMI
$\bigl[\begin{smallmatrix} I & [D, x] \\ [D, x]^* & I \end{smallmatrix}\bigr]
\succeq 0$, and a gauge-fixing constraint orthogonalizing $x$ to
$\ker([D, \cdot]) \cap \mathcal{O}_h$); 17 tests in
\texttt{tests/test\_connes\_distance.py} (all pass; runtime ${\sim}9$~s).
Three structural findings emerged at the placeholder convention:

\begin{enumerate}
\item \textbf{SO(3) $m$-reflection forces zero distance.}  For every
   $|n, l, -m\rangle$ and $|n, l, +m\rangle$ pair, $d_D(\phi_{|n,l,-m\rangle},
   \phi_{|n,l,+m\rangle}) = 0$ exactly (one such pair at $n_{\max} = 2$,
   four at $n_{\max} = 3$).  This is the SO(3) $m$-reflection symmetry
   of the operator system showing through, not a pathology: the symmetry
   identifies the two states as members of the same orbit, and any
   point-discriminating function must be constant on the orbit.

\item \textbf{Clean rational structure at $n_{\max} = 2$.}  The five
   non-zero entries of the $5 \times 5$ distance matrix at $n_{\max} = 2$
   group into ratios $1 : 2 : 3$ (specifically $28.87 : 57.73 : 86.60$),
   with $86.6020 = 50\sqrt{3}$ to five decimals.  This rational/algebraic
   structure is lost at $n_{\max} = 3$ as the SDP gains degrees of
   freedom.  We log this as an algebraic fingerprint of the truncated
   metric at small cutoff; whether it persists in any natural
   generalization is open.

\item \textbf{Non-monotonicity in the Fock-graph distance.}  Pearson
   correlation between $d_D(\phi_v, \phi_w)$ and the combinatorial
   Fock-graph distance $d_{\mathrm{graph}}(v, w)$ (counting $\Delta n =
   \pm 1$ and $\Delta l = \pm 1$ ladder steps) is $\rho \approx -0.04$
   at $n_{\max} = 2$, growing to $+0.14$ at $n_{\max} = 3$ (Spearman
   $0.09 \to 0.18$).  The trend is weakly toward correlation but the
   absolute level is far from a discretized-geodesic metric at our
   reachable cutoffs.
\end{enumerate}

These findings should be read with three explicit caveats, none of
which are addressed in the present paper.

\begin{itemize}
\item The SO(4) radial overlap integral $I_{n N n'}^{l L l'}$ in the
   construction of $\mathcal{O}_{n_{\max}}$
   (Remark~\ref{rem:operator_system}) is currently a placeholder rather
   than the genuine Avery--Wen--Avery $3$-$Y$ integral on $S^3$.
   Round~2's propagation-number result
   (Proposition~\ref{prop:propagation_number}) is invariant under this
   choice; the Connes distance is not.  The absolute distance values
   reported above are convention-dependent until the placeholder is
   replaced.

\item The construction is on the scalar Fock basis $\mathcal{H}_{\mathrm{GV}}$
   only; the spinor lift to $\mathcal{H}_{\mathrm{GV}} \otimes
   \mathbb{C}^4$ with the native Camporesi--Higuchi Dirac in the
   $(n_{\mathrm{D}}, \kappa, m_j)$ basis is not yet incorporated.  The
   diagonal scalar Dirac on $\mathcal{H}_{\mathrm{GV}}$ has a large
   $\ker([D, \cdot]) \cap \mathcal{O}$ that makes the SDP unbounded on
   most state pairs; the metric values reported here use an
   off-diagonal Dirac proxy as a finite-distance regulator, which is a
   second convention choice.

\item No comparison to the round-$S^3$ continuum geodesic distance is
   attempted.  Such a comparison is the actual content of Connes--vS
   Gromov--Hausdorff convergence
   (Sec.~\ref{sec:axiom_audit} discussion of Row~9 of the round-1
   mapping), which has not been proved on $S^3$ in the published
   literature.
\end{itemize}

The structural findings (SO(3) $m$-reflection forced zeros; the $1 : 2
: 3$ rational ratios at $n_{\max} = 2$) are robust across the
placeholder choice; the absolute distance values and the
non-monotonicity diagnostic are not.  See
\texttt{debug/wh1\_round3\_connes\_distance\_memo.md} for the full
record.  Caveats (1) and (2) above are subsequently closed in
Remark~\ref{rem:r31_r32_update}; one of the three structural
findings is retracted there as a placeholder artifact.
\end{remark}

\begin{remark}[R3.1+R3.2 closure of round-3 caveats]
\label{rem:r31_r32_update}
Two of the three caveats above have been closed in subsequent
sub-sprints (continuation of the same dispatch on 2026-05-03), and
one of the three structural findings has been retracted as a
placeholder artifact.

\paragraph{R3.1 (Avery--Wen--Avery integral, placeholder $\to$
physical).}  The placeholder for $I_{n N n'}^{l L l'}$ is replaced
with the closed-form Avery--Wen--Avery $3$-$Y$ integral on $S^3$,
factorized as $I_{\mathrm{rad}} \cdot G_{\mathrm{Gaunt}}$ with
$I_{\mathrm{rad}}$ the Gegenbauer triple integral on $\chi$ and
$G_{\mathrm{Gaunt}}$ the standard $S^2$ Gaunt integral.
Implemented in \texttt{geovac/so4\_three\_y\_integral.py}
(sympy-exact, \texttt{lru\_cache}), verified by hyperspherical
orthonormality at $n_{\max} \le 3$ in exact arithmetic.
$\mathrm{prop}(\mathcal{O}_{n_{\max}}) = 2$ holds at $n_{\max} =
2, 3, 4$ under the physical default with dim sequences $14 \to
25,\ 55 \to 196,\ 140 \to 900$ \emph{bit-identical} to the
placeholder values---confirming the round-2 invariance claim by
replacement, not by sampling.

The Connes distance recomputed under physical values retains
the $m$-reflection forced zeros and the $50\sqrt{3}$ fingerprint,
but \emph{retracts} the $1 : 2 : 3$ rational structure as a
placeholder artifact: under Avery the distinct nonzero values at
$n_{\max} = 2$ are $\{1.27, 86.60, 87.18, 174.93, 175.53,
262.11\}$ with no clean integer grouping.  The $M_{2,0,0}$
multiplier rescales by $\sim 3\times$ between placeholder and
physical values, propagating as $\sim 9\times$ on pairs that use
it twice.  More importantly, the non-monotonicity \emph{flips
sign} at $n_{\max} = 3$: Pearson goes from $+0.14$ (placeholder)
to $-0.36$ (physical), Spearman from $+0.18$ to $-0.33$.  The
optimistic ``slow trend toward correlation'' reading is
retracted; under physical values the metric is
\emph{anti-correlated} with the Fock-graph distance.  Memo:
\texttt{debug/wh1\_r31\_avery\_wen\_avery\_memo.md}.

\paragraph{R3.2 (spinor lift, scalar Fock $\to$ Weyl spinor on
$S^3$).}  The construction extends to the Weyl sector of the
Camporesi--Higuchi spinor bundle on $S^3$ via Clebsch--Gordan
decomposition $|\psi_{l, j = l + 1/2, m_j}\rangle = \alpha_+
|Y_{l, m_j - 1/2}\rangle \otimes |\!\!\uparrow\rangle + \alpha_-
|Y_{l, m_j + 1/2}\rangle \otimes |\!\!\downarrow\rangle$ with
$\alpha_\pm = \sqrt{(l \pm m_j + 1/2)/(2l + 1)}$.  Each scalar
Avery multiplier $M_{NLM}$ lifts to a CG-weighted sum of two
scalar matrix elements; the resulting
$\widetilde{\mathcal{O}}_{n_{\max}}$ is API-compatible with the
existing SDP machinery.  Implemented in
\texttt{geovac/spinor\_operator\_system.py}.  $\dim
\widetilde{\mathcal{O}} = \dim \mathcal{O}$ at every $n_{\max}$
(the CG decomposition introduces no new linear independencies);
identity, $\ast$-closure, and $m_j$-reflection forced zeros all
generalize naturally.

Two clean R3.2 findings.  \textbf{(i)} \emph{The truthful
Camporesi--Higuchi Dirac on the spinor bundle is too
$n$-degenerate} for the Connes distance to be well-defined on
cross-shell pure-state pairs: at $n_{\max} = 2$, $24$ of $28$
cross-pairs give $+\infty$ because the $n$-shell-diagonal
multipliers $M_{N, L = \mathrm{even}, M = 0}$ commute with $D =
\mathrm{diag}(n_{\mathrm{Fock}} + 1/2)$, making the SDP
unbounded.  This is the spinor analog of the round-3
\texttt{shell\_scalar} diagnostic and confirms the
$n$-degeneracy obstruction is structural, not scalar-only.
\textbf{(ii)} \emph{Under a CH+offdiag perturbation that lifts
the $n$-degeneracy, the metric remains anti-correlated} with the
Fock-graph distance: Pearson nz $-0.36$ at $n_{\max} = 2$,
$-0.26$ at $n_{\max} = 3$.  The non-monotonicity is therefore
robust under both the placeholder $\to$ Avery upgrade (R3.1)
\emph{and} the scalar $\to$ spinor lift (R3.2).  Memo:
\texttt{debug/wh1\_r32\_spinor\_lift\_memo.md}.

\paragraph{Disambiguation settled.}  The round-3 \S 5.4
two-reading question (``placeholder artifact'' vs
``pure-state-distance pathology'') lands on the latter.  R3.1
rules out the placeholder reading (anti-correlation
\emph{strengthens} under physical values).  R3.2 rules out the
scalar-Dirac-proxy reading (anti-correlation persists under the
spinor lift).  The Connes distance on $\mathcal{O}_{n_{\max}}$ or
its spinor lift, with pure node-evaluation states, is structurally
not a discretization of the round-$S^3$ geodesic distance at any
cutoff we can reach.  The third caveat above (no GH comparison
on the full state space) remains the natural next move; whether
the metric becomes monotone in the GH limit via a Kantorovich--
Wasserstein lift to mixed states requires Peter--Weyl on
$\mathrm{SU}(2)$ (genuinely new NCG mathematics; deferred by
Connes \& van~Suijlekom 2021 three times, no published treatment
on $S^3$).
\end{remark}

\begin{remark}[Two-sided structural alignment]
\label{rem:two_sided_alignment}
Connes \& van~Suijlekom's worked $S^1$ example exhibits the
spectral-vs-circulant distinction along two axes simultaneously:
the Toeplitz operator system $C(S^1)^{(n)}$ has $\mathrm{prop} = 2$
(their Proposition~4.2), while the circulant matrix algebra
$C(C_m)$ has $\mathrm{prop} = 1$ trivially (it is its own
$C^*$-algebra).  Eqs.~\eqref{eq:prop_two}
(Proposition~\ref{prop:propagation_number}) and~\eqref{eq:prop_circulant}
(Corollary~\ref{cor:structural_specificity}) reproduce both axes
verbatim on $S^3$:
\[
\renewcommand{\arraystretch}{1.15}
\begin{array}{c|c|c}
 & \text{spectral truncation} & \text{circulant analog} \\ \hline
S^1 \text{ (Connes--vS)} & C(S^1)^{(n)},\ \mathrm{prop} = 2 &
   C(C_m),\ \mathrm{prop} = 1 \\
S^3 \text{ (this work)}  & \mathcal{O}_{n_{\max}},\ \mathrm{prop} = 2 &
   A_{\mathrm{circ}, N},\ \mathrm{prop} = 1 \text{ or } \infty \\
\end{array}
\]
The two-sided match foreclosing the natural objection that
$\mathrm{prop} = 2$ might be a generic feature of any finite
truncation of $S^3$.  Together with the SO(4)-isometry-preservation
property of Sec.~\ref{sec:axiom_audit} below (the property Connes
\& van~Suijlekom themselves identify as distinguishing spectral
truncation from circulant discretization), this is the strongest
quantitative alignment between GeoVac and the published
spectral-truncation framework currently known.
\end{remark}

\subsection{The Hilbert space $\mathcal{H}_{\mathrm{GV}}$}

Let $V_{\mathrm{Dirac}}$ denote the set of Dirac-graph nodes at the
matching truncation, labeled by Camporesi--Higuchi triples
$(n_{\mathrm{D}}, \kappa, m_j)$ with $0 \le n_{\mathrm{D}} \le
n_{\max} - 1$ (the standard convention $n_{\mathrm{Fock}} = n_{\mathrm{CH}}
+ 1$ from \texttt{geovac/dirac\_matrix\_elements.py}), Dirac--Coulomb
relativistic angular momentum $\kappa \in \{-l-1, +l\}$ for each
$0 \le l \le n_{\mathrm{D}}$, and total angular momentum projection
$m_j \in \{-j, \ldots, j\}$ where $j = |\kappa| - 1/2$.  The total
node count is
\[
   N_{\mathrm{Dirac}}(n_{\max}) = \sum_{n=0}^{n_{\max}-1}
     g_n^{\mathrm{D}} = \sum_{n=0}^{n_{\max}-1} 2(n+1)(n+2)
   = \tfrac{2}{3} n_{\max}(n_{\max}+1)(n_{\max}+2).
\]

\begin{definition}[Spinor space on the Dirac graph]
\label{def:H_GV}
$\mathcal{H}_{\mathrm{GV}}^{(n_{\max})} := \mathbb{C}^{V_{\mathrm{Dirac}}}$,
with the standard Hermitian inner product on $\mathbb{C}^{N_{\mathrm{Dirac}}}$.
\end{definition}

The faithful $*$-representation of $\mathcal{A}_{\mathrm{GV}}$ on
$\mathcal{H}_{\mathrm{GV}}$ is by the natural pullback: for
$f \in \mathcal{A}_{\mathrm{GV}}$ and $\psi \in \mathcal{H}_{\mathrm{GV}}$,
\begin{equation}
   (\pi(f) \psi)(n_{\mathrm{D}}, \kappa, m_j)
   := f(n_{\mathrm{Fock}}(n_{\mathrm{D}}), l(\kappa), m_l(\kappa, m_j))
   \cdot \psi(n_{\mathrm{D}}, \kappa, m_j),
   \label{eq:representation}
\end{equation}
where the maps $n_{\mathrm{Fock}}, l, m_l$ are the standard Dirac-to-Fock
quantum-number bridges:
$n_{\mathrm{Fock}}(n_{\mathrm{D}}) = n_{\mathrm{D}} + 1$,
$l(\kappa) = -1 - \kappa$ if $\kappa < 0$ and $l(\kappa) = \kappa$ if
$\kappa > 0$, and $m_l$ via the standard
$|j, m_j\rangle = \langle l, m_l; \tfrac{1}{2}, m_s | j, m_j \rangle
|l, m_l\rangle |s, m_s\rangle$ Clebsch--Gordan decomposition.
Each $\pi(f)$ acts diagonally; in the spinor basis it is the diagonal
matrix $\mathrm{diag}(f(v_{\mathrm{Fock}}(\,\cdot\,)))$.  The
representation is faithful: $\pi(f) = 0 \Leftrightarrow f \equiv 0$
(every Fock node has at least one Dirac node above it for $n_{\max}
\ge 1$).

\subsection{The Dirac operator $D_{\mathrm{GV}}$}

The continuum Camporesi--Higuchi Dirac operator on the unit
three-sphere has spectrum $|\lambda_n^{\mathrm{D}}| = n + 3/2$ with
degeneracy $g_n^{\mathrm{D}} = 2(n+1)(n+2)$~\cite{camporesi_higuchi1996}.
We promote this to a finite operator on $\mathcal{H}_{\mathrm{GV}}$
in two equivalent ways and verify their agreement in the truncation.

\begin{definition}[Spectral form of $D_{\mathrm{GV}}$]
\label{def:D_GV_spectral}
$D_{\mathrm{GV}}^{\mathrm{spec}}$ is the diagonal operator on
$\mathcal{H}_{\mathrm{GV}}$ with eigenvalues
$\lambda_n^{\mathrm{D},\pm} = \pm (n_{\mathrm{D}} + 3/2)$
on the eigenstate $|n_{\mathrm{D}}, \kappa, m_j\rangle$ with sign
$+$ for $\kappa < 0$ and $-$ for $\kappa > 0$ (the standard
``large-component'' / ``small-component'' sign choice of the
Camporesi--Higuchi convention~\cite{camporesi_higuchi1996}).
\end{definition}

\begin{definition}[Graph form of $D_{\mathrm{GV}}$]
\label{def:D_GV_graph}
$D_{\mathrm{GV}}^{\mathrm{graph}}$ is the discrete first-order
hopping operator on the Dirac graph defined by the Camporesi--Higuchi
nearest-neighbor edge set with weights chosen so that the spectrum
matches Definition~\ref{def:D_GV_spectral} in the truncation.  The
edge set and weights are documented in
\texttt{geovac/dirac\_matrix\_elements.py}; the explicit form, written
in the $(n_{\mathrm{D}}, \kappa, m_j)$ basis, is the
Lichnerowicz-derived first-order operator with $\kappa$-coupled
adjacent-shell amplitudes.
\end{definition}

\begin{proposition}[Equivalence of spectral and graph forms]
\label{prop:D_equiv}
For all $n_{\max} \ge 1$, $D_{\mathrm{GV}}^{\mathrm{spec}}
= U^{\dagger} D_{\mathrm{GV}}^{\mathrm{graph}} U$ for the unitary
$U$ that diagonalizes $D_{\mathrm{GV}}^{\mathrm{graph}}$.  The
agreement is exact in finite arithmetic at every truncation; in
particular the spectrum, multiplicities, and trace invariants
$\mathrm{Tr}\,(D_{\mathrm{GV}})^{2k}$ match the truncated
Camporesi--Higuchi continuum values.
\end{proposition}

\begin{proof}[Proof sketch.]
Both operators are constructed to have the truncated
Camporesi--Higuchi spectrum by definition.  The equivalence is the
statement that there exists an orthonormal basis of
$\mathcal{H}_{\mathrm{GV}}$ in which both are diagonal, which holds
because $D_{\mathrm{GV}}^{\mathrm{graph}}$ is self-adjoint by
construction (Hermitian conjugate edge weights) and has the prescribed
spectrum.  Numerical agreement at $n_{\max} \le 6$ is verified by
existing tests in \texttt{tests/test\_dirac\_matrix\_elements.py} (108
tests) and \texttt{tests/test\_qed\_self\_energy.py}.
\end{proof}

We use the spectral form $D_{\mathrm{GV}}^{\mathrm{spec}}$ (henceforth
$D_{\mathrm{GV}}$) for axiom verification because it is closed-form;
the graph form is what one uses for finite-truncation computations of
the spectral action.  The $\pi$-free certificate of the Dirac sector
(Paper~28~\S graph\_native\_qed.6) applies to both: every eigenvalue
is a half-integer rational, every degeneracy is a positive integer.

\subsection{Bounded commutators}

\begin{proposition}[Bounded commutators]
\label{prop:bounded_commutators}
For every $a \in \mathcal{A}_{\mathrm{GV}}$, the commutator
$[D_{\mathrm{GV}}, \pi(a)]$ is a bounded operator on
$\mathcal{H}_{\mathrm{GV}}$.
\end{proposition}

\begin{proof}
$D_{\mathrm{GV}}$ has finite-dimensional domain
$\mathcal{H}_{\mathrm{GV}}$ at every $n_{\max}$, hence bounded.
$\pi(a)$ is a diagonal multiplication operator, also bounded.  Any
commutator of bounded operators is bounded.
\end{proof}

This is the trivial but necessary check: in finite truncation, every
operator is bounded.  The non-trivial content---that
$\|[D_{\mathrm{GV}}, \pi(a)]\|$ is controlled uniformly in
$n_{\max}$---is what makes the continuum limit a bona-fide spectral
triple in the sense of Connes; for a Lipschitz function $a$, the
commutator norm equals the Lipschitz constant on the graph metric, in
exact analogy with the smooth case.

\subsection{The triple, summarized}

Putting Definitions~\ref{def:A_GV},~\ref{def:H_GV},~\ref{def:D_GV_spectral}
and Proposition~\ref{prop:bounded_commutators} together:

\begin{theorem}[The GeoVac spectral triple]
\label{thm:GV_triple}
For every $n_{\max} \ge 1$, the data
$(\mathcal{A}_{\mathrm{GV}}^{(n_{\max})},
\mathcal{H}_{\mathrm{GV}}^{(n_{\max})},
D_{\mathrm{GV}}^{(n_{\max})})$ is a finite spectral triple in the
sense of~\cite{connes1995}: $\mathcal{A}_{\mathrm{GV}}$ is unital and
involutive, $\pi: \mathcal{A}_{\mathrm{GV}} \to
\mathcal{B}(\mathcal{H}_{\mathrm{GV}})$ is a faithful
$*$-representation, $D_{\mathrm{GV}}$ is self-adjoint with bounded
commutators.  The continuum limit $n_{\max} \to \infty$ recovers the
canonical spectral triple of the unit Riemannian spin three-sphere
$S^3$ with the Camporesi--Higuchi Dirac operator.
\end{theorem}

The continuum-limit clause is the spectral-truncation statement of
Connes--van~Suijlekom~\cite{connes_vs2021,hekkelman2022,hekkelman_mcdonald2024}:
the GeoVac triples form a Cauchy sequence in the truncation order whose
limit is the standard $S^3$ spectral triple.

% ===================================================================
\section{Axiom Audit}
\label{sec:axiom_audit}
% ===================================================================

We check each Connes axiom in turn, distinguishing what holds
rigorously at finite $n_{\max}$ from what holds only in the continuum
limit.

\subsection{Order zero (involutive faithful representation)}

Holds rigorously by Definitions~\ref{def:A_GV}--\ref{def:H_GV} and
Eq.~\eqref{eq:representation}.  Verified in
\texttt{tests/test\_dirac\_matrix\_elements.py}.

\subsection{Bounded commutators (regularity)}

Holds rigorously at every $n_{\max}$ by
Proposition~\ref{prop:bounded_commutators}.

\subsection{Compact resolvent / finite summability}

In the finite-$n_{\max}$ truncation $\mathcal{H}_{\mathrm{GV}}$ is
finite-dimensional so this is automatic.  The continuum statement is
that $D_{\mathrm{GV}}$ has compact resolvent and is $3^+$-summable
(Dixmier-trace summable in dimension $3+\epsilon$); this is the
Camporesi--Higuchi summability of the continuum operator on $S^3$ and
is a standard result~\cite{camporesi_higuchi1996,connes1995}.

\subsection{Order one}

\begin{proposition}[Order one]
\label{prop:order_one}
For every $a, b \in \mathcal{A}_{\mathrm{GV}}$,
$[[D_{\mathrm{GV}}, \pi(a)], \pi(b)^{\circ}] = 0$ holds
trivially because the right action $\pi(b)^{\circ}$ commutes with the
left action $\pi(a)$ (the algebra is commutative) and with the
diagonal-in-spectral-basis form of $D_{\mathrm{GV}}$.
\end{proposition}

This is the standard observation that the order-one condition is
automatic for any commutative spectral triple.  Order-one becomes
non-trivial only when one extends to an almost-commutative triple
$\mathcal{A}_{\mathrm{GV}} \otimes M_n(\mathbb{C})$ with non-abelian
fiber, at which point it constrains the form of fluctuations of $D$.
This is exactly the setting where Papers~25 ($n=1$, $U(1)$ fiber) and
30 ($n=2$, $\mathrm{SU}(2)$ fiber) live; we treat them in
Sec.~\ref{sec:subsectors}.

\subsection{Real structure}

$S^3$ is odd-dimensional; in dimension~$n$ the canonical
$\mathrm{KO}$-dimension is~$n$ mod~$8$ for a Riemannian spin manifold,
and the sign table~\cite{connes1995,vansuijlekom_book2015} for $n=3$
gives
\[
   \epsilon = -1, \quad \epsilon' = +1, \quad \epsilon'' = ~\text{(none)}~
   ~~\text{(no $\gamma$ in odd dimension).}
\]
The sign $\epsilon = -1$ reflects the fact that the 2-dimensional
spinor representation of $\mathrm{Spin}(3) = \mathrm{SU}(2)$ is
quaternionic (pseudoreal): the standard charge-conjugation operator
on this representation squares to $-\mathbb{1}$, equivalently
$\mathbb{H} \otimes_{\mathbb{R}} \mathbb{C} \cong M_2(\mathbb{C})$
acts on a complex 2-dimensional space with antiunitary $J$ satisfying
$J^2 = -\mathbb{1}$.  This is the same sign that appears in the
canonical Lorentzian Dirac formalism in $3+1$ dimensions for the
charge-conjugation matrix~$C$.

\begin{proposition}[Real structure on $\mathcal{H}_{\mathrm{GV}}$]
\label{prop:reality}
There exists an antilinear isometry $J_{\mathrm{GV}}:
\mathcal{H}_{\mathrm{GV}} \to \mathcal{H}_{\mathrm{GV}}$ with
$J_{\mathrm{GV}}^2 = -\mathbb{1}$ and $J_{\mathrm{GV}} D_{\mathrm{GV}}
= +D_{\mathrm{GV}} J_{\mathrm{GV}}$, namely the Camporesi--Higuchi
charge-conjugation operator restricted to the truncation.
\end{proposition}

\begin{proof}[Construction.]
The continuum charge conjugation $J_{S^3}$ on the
Camporesi--Higuchi spinor bundle satisfies the listed signs by
construction~\cite{camporesi_higuchi1996,connes1995}.  Restriction to
the truncation $\mathcal{H}_{\mathrm{GV}}^{(n_{\max})}$ commutes with
the spectral truncation because $D_{\mathrm{GV}}$ acts diagonally in
the spectral basis and $J_{S^3}$ permutes the spectral basis (sending
$|n_{\mathrm{D}}, \kappa, m_j\rangle \mapsto
|n_{\mathrm{D}}, -\kappa, -m_j\rangle$ up to a phase, whose
antilinear-squaring gives $-\mathbb{1}$ on the quaternionic 2-spinor
representation).  The restriction therefore inherits the signs.
\end{proof}

\paragraph{Finite-$n_{\max}$ audit (Sprint WH1-Connes Step~2,
2026-05).}  The qualitative argument above is now backed by a
quantitative computational audit at $n_{\max} \in \{1, 2, 3\}$ on
both the Weyl sector
(\texttt{geovac/spinor\_operator\_system.py}, $\dim_{\mathcal{H}} = 2,
8, 20$) and the full-Dirac sector
(\texttt{geovac/full\_dirac\_operator\_system.py},
$\dim_{\mathcal{H}} = 4, 16, 40$), with
$J_{\mathrm{GV}}$ realized as an antilinear permutation-phase operator
$J = U \mathcal{K}$ (with $\mathcal{K}$ complex conjugation) and phase
convention $\sigma(l, m_j) = i^{2 m_j} (-1)^l$ giving
$\sigma_a \overline{\sigma_{J(a)}} = -1$ on every basis label.
Implementation in \texttt{geovac/real\_structure.py};
verification in \texttt{tests/test\_real\_structure.py}
(43 tests passing); raw audit data
\texttt{debug/data/wh1\_real\_structure\_audit.json}.
The audit verdict for the four Connes-axiom conditions involving $J$
at finite $n_{\max}$ on the truthful Camporesi--Higuchi Dirac is:

\begin{enumerate}
\item \textbf{$J^2 = -\mathbb{1}$ HOLDS EXACTLY} at every
$n_{\max} \in \{1, 2, 3\}$ on both sectors (bit-exact zero
residual).  The phase formula
$\sigma_a \overline{\sigma_{\pi(a)}} = -1$ is verified pair-by-pair
across the spinor and full-Dirac bases.
\item \textbf{$J_{\mathrm{GV}} D_{\mathrm{GV}}^{\mathrm{truthful}}
= +D_{\mathrm{GV}}^{\mathrm{truthful}} J_{\mathrm{GV}}$ HOLDS
EXACTLY} at every $n_{\max}$ (bit-exact zero residual).  Mechanism:
the truthful CH Dirac is diagonal in
$(n_{\mathrm{Fock}}, l, m_j, \mathrm{chirality})$ with eigenvalue
$\mathrm{chi} \cdot (n_{\mathrm{Fock}} + 1/2)$ depending only on $n$
and~$\mathrm{chi}$; $J_{\mathrm{GV}}$ permutes only $m_j$, leaving
$n$ and $\mathrm{chi}$ fixed, so the matrix relation
$U \overline{D} = D U$ is diagonal-by-diagonal.
\item \textbf{$J_{\mathrm{GV}} \mathcal{O}_{n_{\max}}
J_{\mathrm{GV}}^{-1} = \mathcal{O}_{n_{\max}}$ HOLDS EXACTLY} at
every $n_{\max}$ (bit-exact zero residual on each generator $a$, with
$0$ of $1, 14, 55$ generators failing the $\mathcal{O}$-membership
test).  Mechanism: each multiplier matrix
$\widetilde{M}_{NLM}$ has $J$-conjugate
$\widetilde{M}_{N L (-M)}$ via the $m_l \to -m_l$ symmetry of the
Wigner-3j coupling coefficients (Edmonds Eq.~3.7.5 / Avery--Wen--Avery
3-Y integral); the operator system is closed under this involution.
\item \textbf{$J_{\mathrm{GV}} D_{\mathrm{GV}}^{\mathrm{offdiag}}
= +D_{\mathrm{GV}}^{\mathrm{offdiag}} J_{\mathrm{GV}}$ FAILS}
structurally (max residual $2 \times 10^0$ at $n_{\max} \in \{2, 3\}$;
$10^{-2}$ at $n_{\max} = 1$ from the m-linear lifter alone).  This is
the natural counter-example: the round-3 \texttt{offdiag} Dirac is a
hand-crafted Hermitian perturbation with E1-type hopping
$D[a, b] = \alpha$ on the rule
$|\Delta n| = |\Delta l| = 1, |\Delta m_j| \le 1$, and the
permutation-phase $\sigma$ with $(-1)^l$ parity does not commute with
this perturbation because the $l$-parity sign flips under bra/ket
$l \to l \pm 1$.  The genuine 4-component Camporesi--Higuchi spinor
Dirac (which couples large/small components via the proper
$\mathrm{SU}(2)_L \times \mathrm{SU}(2)_R$ angular gradient with
Clifford-algebra $\gamma$ matrices) would resolve this, but lies
outside the present construction; the offdiag CH is best read as an
SDP-bounding device (round-3 R3.5 Connes distance), not a
spectral-action-foundational object.
\item \textbf{Order-zero $[a, J b J^{-1}] = 0$ and order-one
$[[D, a], J b J^{-1}] = 0$ FAIL at finite $n_{\max} \ge 2$} with
residuals $O(0.05$--$0.2)$ in operator norm.  At $n_{\max} = 2$:
$124$ of $196$ pairs fail order-zero (max residual $5.5\times 10^{-2}$);
$43$ of $196$ pairs fail order-one (max residual $1.0 \times 10^{-1}$);
at $n_{\max} = 3$: $2{,}594$ of $3{,}025$ and $1{,}337$ of $3{,}025$
respectively (max residuals $7.9 \times 10^{-2}$ and
$2.0 \times 10^{-1}$).  The conditions hold trivially at
$n_{\max} = 1$ (single-generator operator system).
\end{enumerate}

\paragraph{Reading of the order-condition failures.}  The
order-zero condition $[a, J b J^{-1}] = 0$ is the
abelian-classical statement \emph{multiplications by classical
functions commute}, automatic for $f, g \in C^\infty(S^3)$ in the
continuum.  At finite $n_{\max}$, the truncated objects
$P_{n_{\max}} f P_{n_{\max}}$ and
$J P_{n_{\max}} g P_{n_{\max}} J^{-1}$ are no longer multiplications
by classical functions; they are operator-system elements
(\cite[Sec.~2]{connes_vs2021}), and their commutator is
generically nonzero.  This is structurally the same finite-resolution
mechanism as the failure of multiplicative closure that defines the
Connes--van Suijlekom truncated operator system itself
($a \cdot b \notin \mathcal{O}$ in general; round-2 propagation
result).  By the conjectured GH convergence on $S^3$
(Track-TS-A R2.5 keystone) the residuals should decay in the
$n_{\max} \to \infty$ limit, with quantitative scaling at $n_{\max}
\in \{4, 5\}$ flagged as deferred future work.

\paragraph{Net status of the real structure at finite
$n_{\max}$.}  The introductory abstract flagged as open
``whether the GeoVac Dirac graph operator carries a real structure
matching the topological $\mathrm{KO}$-dimension of $S^3$ exactly
at finite $n_{\max}$ or only in the continuum limit.''
The Sprint WH1-Connes Step~2 audit closes this question
\emph{positively for the truthful CH Dirac}: the three load-bearing
conditions ($J^2 = -\mathbb{1}$, $JD = +DJ$, and
$J \mathcal{O} J^{-1} = \mathcal{O}$) all hold with bit-exact zero
residual at every finite $n_{\max} \in \{1, 2, 3\}$ on both Weyl and
full-Dirac sectors.  The order-zero and order-one conditions
(involving the abelian-multiplication identity that defines the
classical limit of the algebra) fail at finite $n_{\max} \ge 2$ as
finite-resolution artifacts of the same class as the operator-system
non-multiplicative-closure already named by Connes--van Suijlekom.
The offdiag CH Dirac fails $JD = +DJ$ structurally and is therefore
not a spectral-action-foundational object; the truthful CH is the
natural domain for both the present audit and the inner-fluctuation
calculus relevant to the Track~3 almost-commutative extension
(Sec.~\ref{sec:K_decomposition} discussion).

\subsection{$\mathbb{Z}_2$-grading}

\begin{proposition}[Odd triple]
\label{prop:odd}
The GeoVac spectral triple is \emph{odd}: there is no chirality
operator $\gamma$ on $\mathcal{H}_{\mathrm{GV}}$ commuting with
$\pi(\mathcal{A}_{\mathrm{GV}})$ and anticommuting with
$D_{\mathrm{GV}}$.
\end{proposition}

\begin{proof}
$S^3$ is odd-dimensional and the Camporesi--Higuchi Dirac operator on
$S^3$ has no chirality operator in the standard sense.  The truncation
inherits the same statement.
\end{proof}

This is correct, not a defect.  Odd-dimensional spectral triples do
not carry a $\mathbb{Z}_2$-grading; their spectral action is
constructed from the absolute value $|D|$ rather than from $D^2$
projected onto chirality eigenspaces, and this is precisely the
structure that Paper~28's QED-on-$S^3$ computations use.

\subsection{Summary}

\begin{table}[h]
\centering
\begin{tabular}{lll}
\toprule
Axiom & Status at finite $n_{\max}$ & Status in continuum limit \\
\midrule
Order zero (faithful $*$-rep) & Holds rigorously & Holds rigorously \\
Bounded commutators           & Holds rigorously & Holds rigorously \\
Compact resolvent             & Auto (finite dim) & Holds (CH summability) \\
Finite summability            & Auto              & $3^+$-summable \\
Order one                     & Holds trivially   & Holds trivially \\
Real structure $J$            & Holds (CH conjugation, audited) & Holds \\
$J^2 = -\mathbb{1}$           & Holds exactly $\forall n_{\max}$ & Holds \\
$JD = +DJ$ (truthful CH)      & Holds exactly $\forall n_{\max}$ & Holds \\
$J \mathcal{O} J^{-1} = \mathcal{O}$  & Holds exactly $\forall n_{\max}$ & Holds \\
Order zero $[a, JbJ^{-1}]=0$  & Fails at $n_{\max}\ge 2$ (artifact) & Holds \\
Order one                     & Fails at $n_{\max}\ge 2$ (artifact) & Holds \\
$\mathbb{Z}_2$-grading        & Absent (odd)      & Absent (odd) \\
\bottomrule
\end{tabular}
\caption{Axiom audit for the GeoVac triple
$(\mathcal{A}_{\mathrm{GV}}, \mathcal{H}_{\mathrm{GV}}, D_{\mathrm{GV}})$
at signature $(3, 0)$ (Riemannian, KO-dimension~$3$).
All standard axioms are satisfied; the triple is real, odd, and of
$\mathrm{KO}$-dimension~$3$ in the continuum limit.  The
Lorentzian-extension audit at $(m, n) = (4, 6)$ via Sprint L2-D
extends this picture; see Table~\ref{tab:axiom_audit_lorentzian}
below.}
\label{tab:axiom_audit}
\end{table}

\paragraph{Lorentzian-extension audit at $(m, n) = (4, 6)$
(Sprint L2-B/C/D, 2026-05-16).}
The Riemannian audit above sits at signature $(s, t) = (3, 0)$
(KO-dim~$3$ in the standard $\mathrm{Spin}^c$ convention).  The
Lorentzian extension to $(s, t) = (3, 1)$ via the
Bizi--Brouder--Besnard~2018~\cite{bizi_brouder_besnard2018} Krein-space
spectral-triple classification translates to the BBB index pair
$(m, n) = (4, 6)$ (via $m \equiv t + s$, $n \equiv t - s \pmod 8$).
Sprints L2-B (Krein space, \texttt{geovac/krein\_space\_construction.py}),
L2-C (Lorentzian Dirac, \texttt{geovac/lorentzian\_dirac.py}), and L2-D
(Connes axiom audit, \texttt{geovac/connes\_axiom\_audit\_31.py})
verify four BBB-predicted-sign axioms bit-exactly at
$n_{\max} \in \{1, 2, 3\}$ and $N_t \in \{1, 11, 21\}$ on the truthful
Camporesi--Higuchi spatial Dirac.  Per BBB Table~1 at $(4, 6)$:
$\varepsilon = +1$, $\varepsilon'' = -1$, $\kappa = -1$, $\kappa'' = +1$
(verified directly against \cite{bizi_brouder_besnard2018} v2 PDF).
Sub-deliverables and verdicts are summarised in
Sec.~\ref{sec:lorentzian_audit_paper32} and the subsections
\S~\ref{sec:l2b_krein}, \S~\ref{sec:l2c_lorentzian_dirac},
\S~\ref{sec:l2d_axiom_audit_31}.

\begin{table}[h]
\centering
\begin{tabular}{lll}
\toprule
Axiom & KO-dim 3 (Riemannian) & $(m, n) = (4, 6)$ Lorentzian \\
\midrule
$J^2$
  & $-\mathbb{1}$ (bit-exact, $\forall n_{\max}$)
  & $+\mathbb{1}$ (bit-exact, BBB $\varepsilon = +1$) \\
$JD$
  & $+DJ$ (bit-exact, $\forall n_{\max}$)
  & $+DJ$ (bit-exact, BBB universal Sec~5(v)) \\
$J\chi$
  & n/a (no $\chi$ in odd dim)
  & $\{J, \chi\} = 0$ (bit-exact, BBB $\varepsilon'' = -1$) \\
$J\eta$
  & n/a (no $\eta$ in Riemannian)
  & $\{J, \eta\} = 0$ (bit-exact, BBB $\varepsilon\kappa = -1$) \\
$\eta\chi$
  & n/a
  & $\{\eta, \chi\} = 0$ (bit-exact, BBB $\varepsilon''\kappa'' = -1$) \\
$\chi D$
  & n/a (no $\chi$)
  & $-D\chi$ (BBB universal); FAILS on truthful $D_{\mathrm{GV}}$
    (structural finding, see below) \\
Order zero $[a, JbJ^{-1}]$
  & 5--20\% finite-resolution
  & $\le 7\%$ finite-resolution (sample-of-3) \\
Order one $[[D,a], JbJ^{-1}]$
  & 5--20\% finite-resolution
  & $\le 11\%$ finite-resolution (sample-of-3) \\
\bottomrule
\end{tabular}
\caption{Lorentzian-extension axiom audit at $(m, n) = (4, 6)$
via Sprint L2-B/C/D, 2026-05-16
(\texttt{debug/l2\_\{b,c,d\}\_*\_memo.md},
\texttt{debug/data/l2\_d\_connes\_axiom\_audit\_31.json}).
Four BBB-predicted-sign axioms hold bit-exactly across the
$(n_{\max}, N_t) \in \{1, 2, 3\} \times \{1, 11, 21\}$ panel.  The
BBB universal $\chi D = -D\chi$ axiom fails structurally on the
truthful Camporesi--Higuchi $D_{\mathrm{GV}}$:\ this is a structural
feature of GeoVac, not a bug, and is the load-bearing scope finding
of Sprint L2-D.  Resolution path R1+R2 (truthful for axioms requiring
$JD = +DJ$, offdiag CH for axioms requiring $\{\chi, D\} = 0$) mirrors
the Riemannian-side R3.5/Paper~42 trade.  WH1 PROVEN
(Theorem~\ref{thm:gh_convergence}) is NOT re-opened by these results;
the Lorentzian audit extends the framework rather than re-tests its
Riemannian foundation.}
\label{tab:axiom_audit_lorentzian}
\end{table}

\paragraph{Structural finding on $\chi D$.}
The universal BBB Sec~5(v) axiom $\chi D = -D \chi$ fails bit-exactly
on the pair $(\gamma^5_K, D_L)$ built in Sprints L2-B/C, with Frobenius
residual $2\|D_{\mathrm{GV}}\|_F$ at $N_t = 1$ (numerical values
$6.00, 18.33, 38.88$ at $n_{\max} = 1, 2, 3$ respectively).  The
mechanism is structural and follows from two locked Riemannian-side
conventions:\ (i) the Sprint L2-B convention identifies
\verb|FullDiracLabel.chirality| with the $\gamma^5$ eigenvalue (the
chiral-basis Peskin--Schroeder convention which keeps the Riemannian
limit identification with $\mathcal{H}_{\mathrm{GV}}$ transparent);
(ii) the truthful Camporesi--Higuchi Dirac is diagonal in
$\mathrm{chirality}$ with eigenvalue $\mathrm{chi} \cdot
(n_{\mathrm{Fock}} + 1/2)$, so $D_{\mathrm{GV}}$ \emph{commutes} with
$\gamma^5$ rather than anticommuting with it.  These two facts together
force $\{\gamma^5, D_L\} = 2i\,(\gamma^5 D_{\mathrm{GV}}) \otimes
I_{N_t} \neq 0$ in general.  This is a structural property of
GeoVac's chirality-diagonal $D_{\mathrm{GV}}$ in the Peskin--Schroeder
chiral basis, not a basis-convention error and not a defect of the
Lorentzian lift.  Three resolutions are documented:
\textbf{(R1) accept the structural finding} (truthful $D_{\mathrm{GV}}$
is Krein-self-adjoint with the BBB-predicted (4,~6) signs on the
$J$-relations and is the right object for axioms requiring
$JD = +DJ$, even though it does not form a full BBB indefinite
spectral triple);
\textbf{(R2) use the offdiag Camporesi--Higuchi Dirac} for axioms that
require $\{\chi, D\} = 0$ (the same offdiag CH used as an SDP-bounding
device in the Riemannian Connes-distance round 3 R3.5);
\textbf{(R3) redefine $\gamma^5$ on $\mathcal{H}_{\mathrm{GV}}$} as an
off-diagonal grading anticommuting with truthful $D_{\mathrm{GV}}$,
which breaks Sprint L2-B's chirality-as-$\gamma^5$ identification and
is essentially a basis-convention rebuild.  The recommended path for
Sprint L2-E (Krein-level Paper~42 redo) is R1 + R2 in parallel,
matching the Riemannian-side practice of using truthful CH for
spectral-action and reality-condition tests while using offdiag CH for
SDP-bounding Connes-distance computations.

\begin{theorem}[Rigorous identification]
\label{thm:axiom_complete}
The continuum limit of the GeoVac spectral triple is the canonical
spectral triple of the unit Riemannian spin three-sphere with the
Camporesi--Higuchi Dirac operator.  In particular it is real, odd, and
of $\mathrm{KO}$-dimension~$3$.
\end{theorem}

This is a structural theorem; no fitting parameters or empirical
inputs.  The audit in Table~\ref{tab:axiom_audit} establishes that the
spectral-triple framing is not metaphorical.

% ===================================================================
\section{Sub-sector Identification: Papers 25, 28, 30, 31}
\label{sec:subsectors}
% ===================================================================

We now read the major framework papers as projections of the GeoVac
triple.

\subsection{Paper 28 (QED on $S^3$): the spectral action}

Paper~28~\cite{paper28} computes Tr~$f(D^2/\Lambda^2)$ and related
spectral sums on the Dirac operator $D_{\mathrm{GV}}$ in the continuum
limit, with $\mathrm{SO}(4)$ vertex selection rules from the
Hopf-bundle decomposition.  In the language of Sec.~\ref{sec:connes_data}:
\begin{itemize}\setlength{\itemsep}{0pt}
\item Theorem~T9 (Paper~28 Sec.~T9): the spectral zeta of $D^2$ is a
two-term polynomial in $\pi^2$ with rational coefficients at every
integer $s$.  In spectral-triple language: the heat-kernel
coefficients $a_k(D^2)$ on $S^3$ saturate at $k=2$ (Paper~28's
``SD two-term exactness theorem'').
\item Self-energy structural zero (Paper~28 Theorem~4):
$\Sigma(n_{\mathrm{ext}}=0) = 0$ on the continuum spectral sum,
identically by vertex parity.  This is the standard heat-kernel
bosonic sector cancellation specialized to the Camporesi--Higuchi
spectrum.
\item Vector-photon QED on $S^3$ (Paper~28 \S vector\_photon\_qed):
adding photon $(q, m_q)$ labels with Wigner 3$j$ vertex coupling
recovers 7/8 selection rules for scalar electrons and 8/8 for Dirac
electrons.  The 1/(4$\pi$) per-loop calibration is the $S^2$ Weyl
exchange constant of the Hopf base.
\end{itemize}

In the spectral-triple language Paper~28 is the spectral-action
computation $\mathrm{Tr}\,f(D_{\mathrm{GV}}^2/\Lambda^2)$ at one loop;
the Camporesi--Higuchi spectrum is the input data of $D_{\mathrm{GV}}$;
the SO(4) selection rules are the harmonic decomposition of the spinor
bundle.  Paper~28 is not a free-standing observation about $S^3$; it is
the spectral action of the present triple.

\subsection{Paper 25 (Hopf graph as lattice gauge): the Cartan torus}

Paper~25~\cite{paper25} introduces the signed node-edge incidence
matrix $B$ of the Fock graph and identifies $L_0 = BB^{\top}$ as the
matter propagator and $L_1 = B^{\top}B$ as the discrete
Hodge-1 Laplacian.  In the spectral-triple language: extend
$\mathcal{A}_{\mathrm{GV}}$ to
$\mathcal{A}_{\mathrm{GV}} \otimes \mathbb{C}$ (trivial $M_1$ fiber)
and consider inner fluctuations of $D_{\mathrm{GV}}$ in the abelian
$U(1)$ inner-derivation algebra.  The fluctuated Dirac operator
$D_{\mathrm{GV}} + A$, with $A \in \Omega_D^1$ a self-adjoint 1-form
in the $D_{\mathrm{GV}}$-calculus, has $L_1$ as its quadratic kinetic
term.  This is the Cartan-subalgebra (maximal-torus) sector.

Paper~25's $L_1$ is the $U(1)$-fluctuated kinetic term of the GeoVac
spectral action; its assignment as ``photon propagator'' in Paper~25
is exactly the spectral-action reading.

\subsection{Paper 30 (SU(2) Wilson gauge): the full non-abelian sector}

Paper~30~\cite{paper30} builds the $\mathrm{SU}(2)$ Wilson lattice
gauge theory on the same graph and proves: (a) the maximal-torus
reduction reproduces Paper~25's $U(1)$ action exactly; (b) the
weak-coupling kinetic term equals $L_1$ acting on
$\mathfrak{su}(2)$-valued 1-cochains.  In the spectral-triple
language: extend $\mathcal{A}_{\mathrm{GV}}$ to
$\mathcal{A}_{\mathrm{GV}} \otimes M_2(\mathbb{C})$ (a non-abelian
$M_2$ fiber, equivalently a rank-1 non-abelian almost-commutative
extension of the GeoVac triple).  The order-one condition for this
extension is non-trivial and constrains the form of fluctuations to be
$\mathfrak{su}(2)$-valued connections, exactly Paper~30's link
variables.  The Wilson plaquette action is the lattice-gauge
implementation of the spectral-action curvature term in this
extension.

Paper~30 is the rank-1 non-abelian almost-commutative extension of the
GeoVac triple; Paper~25 is the abelian projection of Paper~30;
together they exhibit the standard Cartan torus / full group structure
of the gauge sector of the present triple.

\subsection{Paper 31 (universal/Coulomb partition): the
$\mathcal{A}/D$ split}

Paper~31~\cite{paper31} partitions GeoVac results into a universal
sector that depends only on the algebra $\mathcal{A}$ (angular
sparsity, composed scaling, selection rules) and a potential-specific
sector that depends on the Dirac operator $D$ (spectrum, calibration
$\kappa$, transcendental content, rigidity, $\alpha$ conjecture).  In
the spectral-triple language this is exactly the $\mathcal{A}/D$ split:
properties of $\mathcal{A}_{\mathrm{GV}}$ alone (rank, dimension,
representation content) transfer to any spectral triple with the same
algebra; properties of $D_{\mathrm{GV}}$ specifically depend on the
choice of Dirac operator.

Paper~31 is the structural reading of Paper~22 (angular sparsity
theorem, an $\mathcal{A}$-statement) and Paper~23 (nuclear shell
model, a different $D$-choice on the same $\mathcal{A}$): when one
keeps $\mathcal{A}$ fixed and varies $D$, the universal results
transfer; when one varies $\mathcal{A}$ (say, replace $S^3$ with
$S^5$, Paper~24), the entire structural picture reorganizes.

\subsection{Summary of sub-sector identification}

\begin{table}[h]
\centering
\begin{tabular}{lll}
\toprule
Paper & Object in this paper's language & Sector \\
\midrule
Paper~7 & $\mathcal{A}_{\mathrm{GV}}$ continuum limit & Manifold sector \\
Paper~22 & Selection rules from $\mathcal{A}_{\mathrm{GV}}$ & $\mathcal{A}$ sector (universal) \\
Paper~23 & $D$-replacement: nuclear shell model & Different $D$, same $\mathcal{A}$ \\
Paper~24 & Different $\mathcal{A}$: $S^5$ holomorphic & Distinct triple (HO) \\
Paper~25 & $U(1)$-fluctuated $D_{\mathrm{GV}} + A$ & Cartan-torus inner fluctuation \\
Paper~28 & Tr~$f(D_{\mathrm{GV}}^2/\Lambda^2)$ & Spectral action (one loop) \\
Paper~30 & $\mathrm{SU}(2)$-fluctuated, $M_2$ fiber & Non-abelian extension \\
Paper~31 & $\mathcal{A}/D$ split & Structural partition \\
\bottomrule
\end{tabular}
\caption{Sub-sector identification of the major framework papers as
projections of the GeoVac spectral triple.}
\label{tab:subsectors}
\end{table}

\paragraph{Track NI as an explicit composed real-space spectral triple.}
The ``Different $D$, same $\mathcal{A}$'' row of
Table~\ref{tab:subsectors} (Paper~23~\cite{paper23}, nuclear shell
model) covers a second class of construction in Paper~23~\S VI: the
composed nuclear--electronic deuterium Hamiltonian, a 26-qubit
Hilbert-space tensor product of a 16-qubit deuteron register
(harmonic-oscillator basis) with a 10-qubit hydrogenic electronic
register (Fock-projected $S^3$ lattice), coupled through a
cross-register hyperfine $\mathbf{I}\cdot\mathbf{S}$ operator.  This
is structurally distinct from both the SM almost-commutative geometry
(which couples a Riemannian factor with a finite \emph{internal}
factor) and from the standard spectral-truncation convergence
machinery (which handles single triples or one-infinite-plus-one-finite
products);\ it is an explicit Connes-style real-space multi-particle
spectral triple, validated against the analytic $21\,$cm
singlet--triplet hyperfine gap at $1.62 \times 10^{-7}$ Ha.  The
positioning observation (that the canonical NCG / spectral-triple
literature does not contain a directly comparable construction) is
formalized in the standalone Track NI Zenodo memo
(\verb|papers/group4_quantum_computing/track_ni_spectral_triple_zenodo.md|).
The W2b cross-manifold structural blocker of
Sec.~\ref{sec:frontier_framing} is exactly the convergence-theorem
gap that surrounds this construction.

\paragraph{Paper 39 closes the same-manifold tensor-product case.}
The same-manifold sub-case of W2b --- the tensor product
$\Tcal_{\sthree}^{\lambda_a} \otimes \Tcal_{\sthree}^{\lambda_b}$ of
two truncated Camporesi--Higuchi spectral triples on $\sthree$ at
distinct focal lengths $\lambda_a \neq \lambda_b$ --- is closed in a
companion standalone paper (Paper~39~\cite{paper39}). The result
extends Paper~38's five-lemma machinery factor-by-factor to the
KO-dim-$6$ tensor product, with the joint Lipschitz comparison
constant and the joint Lipschitz-distortion height term both
controlled by the Connes--Marcolli graded anticommutation
$\{D_a \otimes I, \gamma_a \otimes D_b\} = 0$ via a Pythagorean
refinement of the triangle inequality on the graded module. This is
the mathematical foundation for the GeoVac multi-focal-composition
program, where two electrons in distinct hydrogenic shells live on
$\sthree$ at distinct focal lengths $\lambda = n/Z$;\ the convergence
theorem of Paper~39 says that the cross-register joint
operator system on $\Op_{\nmaxa} \otimes \Op_{\nmaxb}$ converges
qualitatively to the continuum tensor algebra $C^\infty(\sthree
\times \sthree)$ in the Latr\'emoli\`ere propinquity. The two
remaining cross-manifold variants
($\Tcal_{\sthree} \otimes \Tcal_{S^5}^{\mathrm{Bargmann}}$ and
$\Tcal_{\sthree} \otimes \Tcal_{S^5}^{\mathrm{Riem}}$) remain the W2b
content of Sec.~\ref{sec:frontier_framing}.

% ===================================================================
\section{The $K = \pi(B + F - \Delta)$ Formula in Spectral-Action
Language}
\label{sec:K_decomposition}
% ===================================================================

Paper~2~\cite{paper2} observed the numerical coincidence
$K = \pi(B + F - \Delta) = 1/\alpha$ at $8.8 \times 10^{-8}$, with
\[
   B = 42, \quad F = \tfrac{\pi^2}{6}, \quad \Delta = \tfrac{1}{40},
\]
and the Phase~4B--4I sprint series identified three independent
spectral homes for the three pieces (CLAUDE.md~\S 2):
\begin{itemize}\setlength{\itemsep}{0pt}
\item $B$ is the finite Casimir trace at $m=3$:
$B = \sum_{n=1}^{3} (2l+1)l(l+1)|_{n,l} = 42$ (Paper~2 Eq.~17).
\item $F = D_{n^2}(d_{\max} = 4) = \zeta_R(2)$, the Fock-degeneracy
Dirichlet series at the packing exponent (Phase~4F $\alpha$-J).
\item $\Delta^{-1} = g_3^{\mathrm{D}} = 40$, the third single-chirality
Camporesi--Higuchi Dirac mode degeneracy (Phase~4H SM-D).
\end{itemize}

In spectral-triple language, $B$, $F$, $\Delta$ live in three different
sectors of the same triple:

\begin{itemize}\setlength{\itemsep}{2pt}
\item $B$ is a \emph{finite scalar Casimir trace} on the truncated
algebra: in the language of Sec.~\ref{sec:axiom_audit}, it is
$\mathrm{tr}_{\mathcal{A}_{\mathrm{GV}}}(C \cdot \mathbb{1})$ for $C$
the second Casimir of $\mathrm{SO}(4)$ at $n_{\max} = 3$.
\item $F$ is an \emph{infinite scalar spectral zeta}: it equals
$\zeta_{\Delta_{S^3}}(2)$ in the continuum limit, where
$\Delta_{S^3}$ is the scalar Laplace--Beltrami operator on $S^3$
(equivalently, the eigenvalue-weighted Dirichlet of the Fock
degeneracy).
\item $\Delta^{-1}$ is a \emph{single-level Dirac degeneracy}: it
equals $g_3^{\mathrm{D}}$, the multiplicity of the third
Camporesi--Higuchi single-chirality eigenvalue
$|\lambda_3^{\mathrm{D}}| = 9/2$.
\end{itemize}

These three are categorically different objects: a finite
combinatorial trace, an infinite analytic Dirichlet series, and a
single-level multiplicity.  They sit in different sectors of the
spectral triple (scalar/spinor; finite/infinite; trace/zeta/dimension)
and Phases 4B--4I exhausted nine attempts to find a common spectral
generator that produces all three simultaneously, without success.

\begin{observation}[Three-sector composition]
\label{obs:three_sector}
The Paper~2 combination rule $K = \pi(B + F - \Delta)$ is a sum across
three structurally distinct sectors of the GeoVac spectral triple.
Each piece has been individually identified as a clean spectral
invariant; the combination rule itself remains a numerical observation.
\end{observation}

This rephrasing is what working hypothesis WH5 in the project register
is built on: $\alpha$ is the conversion factor between three projection
regimes that do not share a generator.  The spectral-triple framing
does not derive $\alpha$; it makes precise where the obstruction is.

\begin{remark}[Sprint A, $\pi^3 \alpha^3$ residual]
\label{rem:sprint_A}
Sprint~A (April~2026, CLAUDE.md~\S 1.7) verified at 80~dps that the
post-cubic residual $K - 1/\alpha - \alpha^2 = 1.2079 \times 10^{-5}$
matches $\pi^3 \alpha^3 = 1.2049 \times 10^{-5}$ to within 0.25\%.
The structural reading is $\pi^3 = \mathrm{Vol}(S^5)$, hinting at a
contribution from the Bargmann--Segal $S^5$ HO sector
(Paper~24~\cite{paper24}) at higher order.  This is a structural
hint, not a derivation; the cross-check (Paper~24's HO triple is a
distinct spectral triple, not a sub-sector of the GeoVac triple)
shows that the $\pi^3 \alpha^3$ identification, if real, is a
\emph{cross-triple} relation and not internal to the GeoVac
spectral-action expansion.
\end{remark}

% ===================================================================
\section{The Coulomb/HO Asymmetry in Spectral-Triple Language}
\label{sec:coulomb_ho}
% ===================================================================

Paper~31~\cite{paper31} formalized a three-layer asymmetry between the
Coulomb spectral triple of the present paper and the harmonic
oscillator spectral triple of Paper~24~\cite{paper24}, summarized
here:

\begin{table}[h]
\centering
\small
\begin{tabular}{lll}
\toprule
Layer & Coulomb (this paper) & Harmonic oscillator (Paper~24) \\
\midrule
Manifold algebra $\mathcal{A}$ &
  Functions on $S^3$ Fock graph &
  Functions on $S^5$ holomorphic Hardy graph \\
Operator order &
  Second-order (Laplace--Beltrami) &
  First-order (complex Euler) \\
Spectrum &
  Quadratic $|\lambda_n| = n^2 - 1$ &
  Linear $E_N = N + 3/2$ \\
Calibration constant &
  $\pi$ (Hopf measure $\mathrm{Vol}(S^2)/4$) &
  None ($\pi$-free in exact arithmetic) \\
Wilson gauge content &
  Full $\mathrm{SU}(2)$ (Paper~30) &
  $U(1)$ only (Sprint~5 negative) \\
$L_0$ matter propagator role &
  Yes (eigenvalues = spectrum) &
  No (only adjacency, not spectrum) \\
\bottomrule
\end{tabular}
\caption{The Coulomb/HO asymmetry from Paper~31, in spectral-triple
language.  The two systems are distinct spectral triples sharing only
the algebra type (functions on a graph) and the trace structure of the
spectral action.}
\label{tab:coulomb_ho}
\end{table}

In spectral-triple language the asymmetry is structural: the Coulomb
case is a Riemannian-spin spectral triple of $\mathrm{KO}$-dimension~3,
the HO case is a spectral triple of complex Euler type on the
holomorphic Hardy space of $S^5$.  These are different objects even
though both are ``almost-commutative'' in the broad sense.  Calibration
$\pi$ enters in the Coulomb case through the Hopf-bundle measure
$\mathrm{Vol}(S^2) / 4$ in the spectral-action coefficients; it does
not enter the HO case because the HO spectral action is built from a
linear (first-order) operator whose heat kernel does not produce $\pi$
at the rational level of the truncation~\cite{paper24}.

\begin{observation}[Coulomb-specificity of $\pi$]
\label{obs:pi_coulomb}
The presence of $\pi$ in the spectral-action coefficients of QED is
specific to the Coulomb-projected $S^3$ triple of this paper, not a
universal feature of GeoVac-style discretizations.  The HO triple is
$\pi$-free at the discrete level.  This is the structural answer to
``why does QED need $\pi$''---it needs $\pi$ because it lives on the
Coulomb-projection $S^3$ triple, where the calibration is the
$S^2$ Hopf-base Weyl constant.
\end{observation}

\begin{theorem}[G3 closure: Bertrand-rigidity category obstruction]%
\label{thm:g3_closure}
There is no tensor-product spectral triple
$\Tcal_{\sthree} \otimes \Tcal^{\mathrm{BL}}_{S^5}$
that simultaneously preserves:
\begin{enumerate}
\item[(a)] the Riemannian-Dirac character of $\Tcal_{\sthree}$
  (second-order spectral action, spinor bundle,
  Camporesi--Higuchi spectrum $|\lambda_n| = n + 3/2$); and
\item[(b)] the Bargmann character of $\Tcal^{\mathrm{BL}}_{S^5}$
  ($\pi$-free certificate, Hardy-sector tower structure,
  first-order Euler spectrum $E_N = N + 5/2$).
\end{enumerate}
The obstruction is categorical and rests on three premises:
\begin{itemize}
\item \emph{Bertrand's theorem (1873):}  the only central
  potentials with all bounded orbits closed are $V = -Z/r$
  (Coulomb) and $V = \omega^2 r^2$ (harmonic oscillator).
\item \emph{Fock rigidity (Paper~23~\cite{paper23}, Thm.~4):}
  the $\sthree$ conformal projection is unique to $-Z/r$;\ the
  projection is second-order (stereographic / conformal), yielding
  a Riemannian spectral triple with the Laplace--Beltrami
  operator.
\item \emph{HO rigidity (Paper~24~\cite{paper24}, Thm.~3):}
  the $S^5$ Bargmann--Segal projection is unique to $\omega^2 r^2$;\
  the projection is first-order (complex-analytic / holomorphic),
  yielding a Berezin--Toeplitz construction with the Euler number
  operator on the Hardy space.
\end{itemize}
\end{theorem}

\begin{proof}[Proof sketch]
\emph{Step 1 (category identification).}
Fock rigidity and HO rigidity place the two potentials in different
mathematical categories: Coulomb $\to$ second-order conformal
$\to$ Riemannian spectral triple; HO $\to$ first-order
holomorphic $\to$ Berezin--Toeplitz quantization.

\emph{Step 2 (exhaustion of tensor-product options).}
Three constructions are available:
(i)~the standard graded tensor product with the round-$S^5$
Riemannian Dirac---this preserves~(a) but violates~(b)
($\pi$ re-enters via the Seeley--DeWitt coefficients on $S^5$);
(ii)~the tensor product with the Bargmann--Euler operator---this
preserves~(b) but violates~(a)
($E_{\mathrm{BL}}$ is not a Dirac-type operator on $S^5$,
being the radial part of the Laplacian restricted to the
holomorphic sector, not an elliptic differential operator
on the total space);
(iii)~a hypothetical framework-extending construction beyond
standard Connes spectral triples and Berezin--Toeplitz
quantization---no such framework exists in published
mathematics.

\emph{Step 3 (irreducibility).}
Bertrand exhausts the potentials;\ the two rigidity theorems fix
the projection type for each.  The category mismatch is a
structural feature of central-force quantum mechanics, not an
artifact of the GeoVac construction.  The five asymmetry layers
of Paper~24~\S~V (spectrum-computing role of $L_0$; calibration
$\pi$; Wilson gauge matter coupling; modular-Hamiltonian
Pythagorean orthogonality; spectral-action gravity termination)
are all downstream consequences of this single categorical
obstruction.
\end{proof}

\begin{remark}[G3 scope]%
\label{rem:g3_scope}
Theorem~\ref{thm:g3_closure} does not assert that $S^3$ and $S^5$
are unrelated --- they are the two Bertrand geometries, deeply connected
through the angular sparsity theorem (Paper~22).  It asserts that their
\emph{spectral-triple-type structures} live in different mathematical
categories, and no published framework bridges the gap while preserving
both structural contents.  Cross-manifold unification, if achievable,
requires framework-extending mathematics beyond the scope of this paper
and the published NCG literature.
\end{remark}

% ===================================================================
\section{The $\pi$-Source Case-Exhaustion Theorem}
\label{sec:pi_source_theorem}
% ===================================================================

The empirical principle of Paper~35~\cite{paper35} (Prediction~1,
graduated to load-bearing status on a cumulative 208/208 confirmation
across sprints KG and TX-B, May~2026) states that a GeoVac observable
contains $\pi$ if and only if its evaluation includes a continuous
integration over a temporal or spectral parameter promoted from the
discrete graph spectrum.  Read in the spectral-triple language of
\S\ref{sec:construction}--\S\ref{sec:axiom_audit}, this principle has
a structural underpinning: the only places in the GeoVac framework
where such continuous integrations occur are anchored to three named
mechanisms of the present triple.  We state and prove the
case-exhaustion.

\begin{theorem}[$\pi$-Source Case-Exhaustion]
\label{thm:pi_source_case_exhaustion}
Let $\mathcal{C} = (P_{i_k} \circ \cdots \circ P_{i_1})$ be any finite
composition of projections drawn from the Paper~34 list~\cite{paper34}
\S\,III.1--\S\,III.15, applied to a Layer-1 datum on
$(\mathcal{A}_{\mathrm{GV}}, \mathcal{H}_{\mathrm{GV}},
D_{\mathrm{GV}})$ or its Bargmann--Segal sibling triple.  The
following are equivalent:
\begin{enumerate}
\setlength\itemsep{0pt}
\item[(i)] The converged value $\mathcal{C}(\mathrm{datum})$ contains a
transcendental of the form $\pi$, $\sqrt{\pi}$, $\pi^k$ for $k \in
\mathbb{Q}$, or $\pi$ in a denominator.
\item[(ii)] The chain $\mathcal{C}$ contains at least one projection
$P_{i_j}$ whose evaluation pipeline includes a continuous integration
over a parameter promoted from the discrete graph spectrum
(heat-kernel time, Matsubara frequency, Mellin contour, Hopf $S^2$-base
angular measure, $\mathrm{SU}(N)$ Haar measure, or
analytic-continuation parameter).
\item[(iii)] The chain $\mathcal{C}$ engages at least one of three
spectral-triple mechanisms:
\begin{itemize}\setlength\itemsep{0pt}
\item[\textbf{M1}] (Hopf-base measure):\ introduces $\pi =
\mathrm{Vol}(S^2)/4$ via the Hopf bundle $S^3 \to S^2 \times S^1$,
or its higher-rank analogues~\cite{paper25,paper30}.
\item[\textbf{M2}] (Seeley--DeWitt heat-kernel coefficient):\
introduces $\sqrt{\pi}$ via heat-kernel coefficients on the
Camporesi--Higuchi Dirac spectrum on unit~$S^3$ (Paper~28
Theorem~T9~\cite{paper28}: $\zeta_{D^2}(s)$ is a two-term polynomial
in $\pi^2$ with rational coefficients at every integer $s \ge 1$).
The sub-mechanism for the spinor sector is the half-integer Hurwitz
$\zeta(s, 3/2)$.
\item[\textbf{M3}] (vertex-topology Dirichlet $L$-content):\
introduces Catalan~$G = \beta(2)$ and Dirichlet $\beta(s)$ via
quarter-integer Hurwitz shifts $\zeta(s, 3/4)$, $\zeta(s, 5/4)$ at
two-loop sunset vertices (Paper~28 Theorem~3:\
$D_{\mathrm{even}}(s) - D_{\mathrm{odd}}(s) = 2^{s-1}(\beta(s) -
\beta(s-2))$).
\end{itemize}
\end{enumerate}
\end{theorem}

\begin{proof}[Proof sketch]
(i)~$\Leftrightarrow$~(ii) is Paper~35 Prediction~1, verified across
208 individual checks (sprints KG, TX-B, May 2026~\cite{paper35}).

(ii)~$\Rightarrow$~(iii):\ the proof body that follows enumerates the
fifteen Paper~34 projections in §III.1--§III.15 (the corpus state at
this theorem's drafting; Paper~34 has since grown to twenty-eight
projections~\cite{paper34}, and the theorem's claim is here restricted
to those first fifteen pending an extension of the case-by-case
enumeration to projections~16--28).  Each of the fifteen Paper~34
projections is
either $\pi$-free at the projection step (projections~1,~3,~4,~5,~8,
9,~12,~13,~14: Fock conformal, Bargmann--Segal, Stereographic, Sturmian,
Wigner~$3j$, Wigner~$D$, Molecule-frame, Drake--Swainson, Rest-mass) or
$\pi$-bearing through a continuous integration that reduces to one of
M1,~M2,~M3 by direct identification (Hopf~bundle, Spectral~action,
Camporesi--Higuchi spinor lift, Wilson plaquette, Vector-photon promotion,
Observation/temporal-window).  The seven $\pi$-free projections live in
either Paper~31's universal sector (algebra-only, transfers to any
spherical-fermion system) or Paper~22's potential-independent angular
sparsity tier; neither sector contains $\pi$.  The eight $\pi$-bearing
projections are tagged in
\texttt{debug/track\_ts\_e1\_theorem\_promotion\_memo.md}~\S2 to one or
more of M1/M2/M3, with closed-form identification for each (Paper~25
for M1, Paper~28 T9 theorem for M2, Paper~28 Theorem~3 for M3).

(iii)~$\Rightarrow$~(i) is closed-form for each mechanism:\ M1's
$\mathrm{Vol}(S^2)/4 = \pi$ identity, M2's $a_0(S^3) = \sqrt{\pi}$
heat-kernel identity, M3's $\beta(s) - \beta(s-2)$ closed form via
the $\chi_{-4}$ Dirichlet character.  Each forces $\pi$ to appear in
any chain that engages it.
\end{proof}

\begin{corollary}[M1 sub-mechanism lies in the localised pure-Tate
sub-ring generated by $\pi$]
\label{cor:m1_pure_tate}
The M1 sub-mechanism (the Hopf-base measure factor
$\mathrm{Vol}(S^2)/4 = \pi$ on the $S^3 \to S^2$ Hopf bundle)
produces period values in the localised pure-Tate sub-ring
$\mathbb{Q}[\pi, \pi^{-1}]$ of mixed Tate periods over $\mathbb{Q}$,
at depth $0$ and Tate weight $\in \mathbb{Z}$.  Every M1 closed form
witnessed across the corpus (Paper 25 Hopf base
$\mathrm{Vol}(S^2)/4 = \pi$; Paper 38 §L2 asymptotic propinquity
rate $4/\pi = \mathrm{Vol}(S^2)/\pi^2$; Paper 40 Thm 1 universal
$4/\pi$ across compact Lie groups; Paper 43 §10.2 Pythagorean
orthogonality $1/\pi^2$; Paper 50 §3 F-theorem $-3\zeta(3)/(16\pi^2)$
as a M1$\times$M3 cross-product) is a $\mathbb{Q}$-linear combination
of $\pi^n$ for $n \in \mathbb{Z}$, with $|n| \le 2$ on the empirical
witnesses logged so far.  Higher integer powers of $\pi$ enter only
through M2 (Seeley--DeWitt) or M3 (vertex-parity sums).  The M1 ring
is convention-free (unlike M2's $\sqrt\pi$ vs $\pi^2$
volume-normalisation choice), because the Hopf-base measure is itself
a volume integral and does not engage the heat-kernel asymptotic.
Sprint M1 Pure-Tate (June~2026, memo
\texttt{debug/sprint\_m1\_pure\_tate\_memo.md}).
\end{corollary}

\begin{corollary}[M2 sub-mechanism inherits Fathizadeh--Marcolli
mixed-Tate classification, with pure-Tate refinement]
\label{cor:m2_mixed_tate}
At the standard volume-normalized Connes--Chamseddine
convention, the M2 sub-mechanism Seeley--DeWitt coefficients on
$\sthree$ are mixed-Tate periods over $\mathbb{Q}$ by direct
inheritance from
Fathizadeh--Marcolli~\cite{fathizadeh_marcolli2016} (the static-R--W
sub-case $a(t) \equiv 1$ on $\mathbb{R} \times \sthree$). The GeoVac
$\sthree$ M2 sector is in fact the strictly smaller pure-Tate
sub-ring
$\bigoplus_k \pi^{2k}\cdot\mathbb{Q}$
of even-weight Tate periods, with no $\zeta(3), \zeta(5)$, or MZV
content appearing at the SD level. This refines F--M's continuum
classification in two structural ways:\ (i) the Dirac SD expansion
is two-term exact ($a_k^{D^2} = 0$ for $k \ge 2$,
Paper~51~\cite{paper51} Cor.~2.1), not generically infinite; (ii)
the scalar Laplacian SD coefficients have the closed-form
geometric-series collapse $a_k^\Delta = 2\pi^2/k!$. The M3
sub-mechanism (vertex-parity Hurwitz, Catalan $G$, $\beta(s)$)
lives in cyclotomic mixed Tate at level $4$ (mixed Tate over
$\mathbb{Z}[i]$, not over $\mathbb{Q}$) and does \emph{not}
contaminate the SD coefficients, by the operator-order partition
of Paper~18~\cite{paper18} \S~III.7 ($k=2$ for SD vs $k=1$ for
vertex parity). Sprint Mixed-Tate Test (June~2026, memo
\texttt{debug/sprint\_mixed\_tate\_test\_memo.md}).
\end{corollary}

\begin{corollary}[M3 sub-mechanism inherits Deligne--Glanois
cyclotomic mixed-Tate classification at level $\le 4$]
\label{cor:m3_cyclotomic_mixed_tate}
The M3 sub-mechanism (vertex-parity Hurwitz / Dirichlet $L$ /
depth-$k$ loop tower) on the discrete Camporesi--Higuchi $\sthree$
spectral triple lies in the cyclotomic mixed-Tate period ring at
level $\le 4$ over $\mathcal{O}_4[1/4] = \mathbb{Z}[i, 1/2]$ at depth
equal to the loop order, by direct application of
Deligne~\cite{deligne2010} (proven for $N \in \{2,3,4,8\}$) and the
explicit basis of Glanois~\cite{glanois2015}.  Three sub-rings are
realised:
\begin{equation*}
\mathcal{MT}(\mathbb{Z})
\;\subset\;\mathcal{MT}(\mathbb{Z}[1/2])
\;\subset\;\mathcal{MT}(\mathbb{Z}[i, 1/2]),
\end{equation*}
selected by (i) un-restricted vs vertex-restricted summation, and
(ii) loop order $=$ motivic depth.  Specifically:
\begin{itemize}
\item Un-restricted Dirac Dirichlet $D(s)$ at integer $s \ge 2$
      (Paper~28~\cite{loutey_paper28} Thm~2) lies in
      $\mathcal{MT}(\mathbb{Z})$ at depth $1$;
\item Vertex-restricted $D_{\mathrm{even}} - D_{\mathrm{odd}} =
      2^{s-1}(\beta(s) - \beta(s-2))$ at integer $s \ge 2$
      (Paper~28~\cite{loutey_paper28} Thm~3) lies in
      $\mathcal{MT}(\mathbb{Z}[i, 1/2])$ at level $4$, depth $1$, via
      Hurwitz at quarter-integer shifts;\ Catalan $G = \beta(2)$ and
      $\beta(4)$ are level-$4$ Dirichlet $L$-value periods in this
      ring;
\item Two-loop irreducible $S_{\min}$ (Paper~28 §$S_{\min}$) lies in
      $\mathcal{MT}(\mathbb{Z}[1/2])$ at depth $2$; the depth-$k$
      tower (Paper~28 Prop.~depth\_k) generates new irreducible
      periods at depth $k = $ loop order in $\mathcal{MT}$ at level
      $\le 4$.
\end{itemize}
The vertex-parity-as-$\chi_{-4}$ Dirichlet character mechanism IS the
Deligne--Glanois Galois descent from level $4$ to level $2$
(Glanois~\cite{glanois2015} Cor.~1.2) realised on the framework's
natural QED observables.  Combined with
Corollary~\ref{cor:m2_mixed_tate}, the master Mellin engine of
Theorem~\ref{thm:pi_source_case_exhaustion} now has a complete
period-theoretic stratification:\ M1 in $\mathbb{Q}[\pi, 1/\pi]$, M2
in the pure-Tate sub-ring $\bigoplus_k \pi^{2k}\cdot\mathbb{Q}$, and
M3 in $\mathcal{MT}(\mathbb{Z}[i, 1/2])$ at level $\le 4$.  Specific
motivic identification of irreducible constants at depth $\ge 2$
($S_{\min}$, $S^{(3)}$, $\ldots$) within the Deligne--Glanois basis
is an open question at each depth.  Sprint M3 Cyclotomic Mixed-Tate
(June~2026, memo
\texttt{debug/sprint\_m3\_cyclotomic\_mixed\_tate\_memo.md}).
\end{corollary}

\begin{remark}[Joint engagement]
M1, M2, M3 are not set-theoretically disjoint as constructions:\ the
spectral action invokes both M1 (Hopf-base measure on the spinor
bundle) and M2 (Seeley--DeWitt coefficients) simultaneously, and the
vacuum-polarization coefficient $\Pi = 1/(48\pi^2)$ (Paper~28~\S\,V)
contains one factor of $\pi$ from each mechanism.  The theorem requires
only that \emph{at least one} of the three be engaged; multiple
engagement is the typical case at higher loop order.  The mechanisms
are role-disjoint in the Paper~18~\S\,III.7 Mellin-engine sense
(M1~$=$ packing-tier $\omega_d$, M2~$=$ heat-kernel via Mellin/Gaussian,
M3~$=$ vertex-parity quarter-integer Hurwitz character mechanism).
The Track~TS-E3 sprint (May~2026) further sharpened this by showing
that M1, M2, M3 are all sub-cases of a single master Mellin-engine
mechanism $\mathrm{Tr}(D^k \cdot e^{-tD^2})$ on a compact Riemannian
manifold, with $k \in \{0, 1, 2\}$ selecting the sub-mechanism.
\end{remark}

\begin{remark}[Status of $K = \pi(B + F - \Delta)$]
\label{rem:K_under_theorem}
Paper~2's combination $K = \pi(B+F-\Delta) \approx 1/\alpha$ matches
its target at $8.8 \times 10^{-8}$.  Under
Theorem~\ref{thm:pi_source_case_exhaustion}, $K$'s chain
(Fock~$\circ$~Hopf$\circ$~spinor) engages M1 (the explicit prefactor $\pi$,
identified by Sprint~A as $\mathrm{Vol}(S^2)/4$~\cite{paper2}) and M2
(in $F = \pi^2/6 = D_{n^2}(d_{\max})$ via the Mellin transform of
the scalar Fock-degeneracy).  The $\pi$-content is therefore consistent
with case-exhaustion.  The remaining anomaly is at the
\emph{depth-prediction} layer of Paper~34~\S\,VIII.3 (depth predicts
$\sim 3\%$ residual at three projections; $K$ matches at
$10^{-8}$): a quantitative-error puzzle, distinct from the
qualitative case-exhaustion of the present theorem, and
unaffected by it.
\end{remark}

\begin{remark}[Falsification target]
A counter-example to Theorem~\ref{thm:pi_source_case_exhaustion}
would be a sixteenth projection (or a finite chain of existing
projections) introducing $\pi$ via a fourth mechanism not in
\{M1, M2, M3\}.  Three natural candidates were investigated in
Track~TS-E3 (May~2026) and each REDUCED to combinations of the
three existing mechanisms:\ continuum-gravity coupling (M1+M2),
anomaly-class projections (M1+M2+M3 via the eta-invariant's
first-order Mellin), and non-trivial principal-bundle projections
(M1 normalisations against integer integrands).  The strongest
remaining computational falsification target is the discrete
first-Chern-class invariant on the Hopf $S^3$ graph at $n_{\max}=3$,
predicted integer-valued by the case-exhaustion; if a $\pi$ appears
unexpectedly, a fourth mechanism M4 would need to be named.
\end{remark}

\begin{remark}[Mechanism-as-domain sharpening (Sprint MR-A/B/C, May~2026)]
\label{rem:master_mellin_domain}
A three-track follow-up sprint (MR-A: Dirac propinquity rate;
MR-B: spectral-action modular residual; MR-C: L2 next-order
constant) tested whether the master Mellin engine of
Paper~18~\S~III.7 is also \emph{predictive} for convergence rates,
not only classifying for $\pi$-sources. The combined verdict
sharpens the case-exhaustion theorem into a mechanism-as-domain
partition.
\begin{itemize}
\item M1 ($k = 0$) is the natural domain of state-space
propinquity rates. The L2 sprint (2026-05-06) measured the
asymptotic constant of the SU(2) scalar central Fej{\'e}r kernel as
$4/\pi = \mathrm{Vol}(S^2)/\pi^2$.
\item M2 ($k = 2$) is the natural domain of heat-kernel and
spectral-action convergence. Sprint MR-B derived a closed form for
the Camporesi--Higuchi modular residual:
$$
\varepsilon(t) \;=\; \sum_{m \ge 1} (-1)^m \sqrt{\pi}\,
e^{-m^2 \pi^2 / t}\left[t^{-3/2} - 2m^2 \pi^2\, t^{-5/2}
   - \tfrac{1}{2}\, t^{-1/2}\right].
$$
Verified at $> 100$ digits at every tested $t$; every coefficient in
$\sqrt{\pi}\,\mathbb{Q} \oplus \pi^2\,\mathbb{Q}$; modular exponent
exactly $\pi^2$ via Jacobi $\vartheta_2$.
\item M3 ($k = 1$) is the natural domain of vertex-restricted
parity-character observables (Paper~28~\S\,QED-vertex:
Catalan~$G$, $\beta(4)$). Sprint MR-A demonstrated that asking
for M3 from a half-integer-only Peter--Weyl propinquity is a
\emph{category error}: the half-integer-spin truncation cannot span
the constant function $\chi_0$ on $\mathrm{SU}(2)$, so the kernel
does not localize and the moment is identically $\gamma_n^{\rm Dirac}
= \pi$ for every $n_{\max} \ge 1$. The L2 scalar rate's
$(4/\pi)\log(n)/n$ decay is an integer$\times$half-integer
interference effect; removing the integer sector kills it. This is
exactly what the M1 mechanism predicts when the sub-bundle lacks
the constant function.
\end{itemize}
The reading is therefore: the index $k \in \{0, 1, 2\}$ classifies
\emph{both} the sub-mechanism and the natural domain of observable
from which its transcendental signature can be extracted. The
case-exhaustion theorem~\ref{thm:pi_source_case_exhaustion}
classifies \emph{which} transcendentals appear; the domain partition
refines to \emph{which observables} can produce them. Sprint MR-C
extracted the L2 next-order constant $c = 4.1093214674877940927\ldots$
to $\geq 80$ dps via Richardson extrapolation but found no PSLQ
identification in $M_1 \cup M_2 \cup M_3 \cup \{\gamma_E, \log 2,
G, \zeta(3)\}$ at coefficient ceiling $10^4$. This is consistent
with $c$ being a downstream Stein--Weiss IBP derivative of the M1
mechanism rather than a direct Mellin output, and is a partial
negative against extending the predictive engine to all
asymptotic-expansion coefficients along subsequent analytical
pipelines. The engine is predictive \emph{within} each domain.
Source: \texttt{debug/mr\_a\_dirac\_propinquity\_*},
\texttt{debug/mr\_b\_spectral\_action\_rate\_*},
\texttt{debug/mr\_c\_l2\_subleading\_*}.
\end{remark}

\begin{remark}[CH-substrate-specificity of the chirality-parity
factorization (Sprint Q5'-S5-falsifier, 2026-06-05)]
\label{rem:q5p_ch_specificity}
The bit-exact $\mathbb{Z}/2$ chirality-parity factorization of the master
Mellin source on the truncated Camporesi--Higuchi spectral triple
(Sprint~Q5'-CH-1) is CH-specific:\ it requires both spectral
$\pm$-symmetry and chirality alignment $\gamma_{ii} = \chi_i$,
$\Lambda_{ii} = \chi_i |\lambda_i|$ of the CH Dirac.  The
Bargmann--Segal Hardy-on-$S^5$ triple (Paper~24) has neither at any
finite $N_{\max}$:\ the Hardy spectrum is strictly positive
(no $\pm$-partner) and the $(N, 0)$ $\mathrm{SU}(3)$ tower carries
no canonical chirality grading (the two natural candidates
$\chi_N = (-1)^N$ and $\chi_l = (-1)^l$ coincide because $(N, 0)$
branching forces $l$-parity $\equiv N$-parity, and neither
reproduces the CH-1 even-$j$ zeros).  The $\omega^{\mathrm{tri}}$
enrichment program of Paper~55~\S~subsec:open\_m2\_m3 correctly
targets the CH triple as substrate; the case-exhaustion
classification~\ref{thm:pi_source_case_exhaustion} remains
universal in scope but its bit-exact \emph{finite-cutoff} witness
form is CH-substrate-specific.  This places a sixth layer on the
Coulomb/HO asymmetry of Paper~24~\S~V.  Source:
\texttt{debug/sprint\_q5p\_s5\_falsifier\_memo.md}.
\end{remark}

\begin{remark}[Stage-1 cochain-level witness of the master-Mellin
partition (Sprint Q5'-Stage1-JLO-nmax2, 2026-06-05)]
\label{rem:q5p_stage1_jlo_witness}
On the truncated Camporesi--Higuchi spectral triple at $n_{\max}
= 2$, the master-Mellin partition $(M_1, M_2, M_3)$ is bit-exactly
visible at the cyclic-cochain level as the symbol
$\omega^{\mathrm{tri}}(\mathcal{T}_2) = (16, 84, 36) \in \mathbb{Z}^3$.
The first two slots live inside a single JLO degree-0 cochain:\
$M_1$ at $t^0$ of $\phi_0^{\mathrm{odd}}(1; t) = \mathrm{Tr}\,
e^{-tD^2}$, $M_2$ at $-t^1$ of the same.  $M_3$ lives outside the
JLO framework on commutative $\mathcal{A} = \mathbb{C}^5$ and instead in
the Connes--Moscovici residue framework as the leading
coefficient of $\mathrm{Tr}(\gamma\,D\,e^{-tD^2})$.  The first
non-trivial JLO cochain degree on commutative $\mathcal{A}$ is
$n = 2$;\ structural vanishing $\phi_1(a_0, a_1; t) \equiv 0$ is
proved bit-exactly for all idempotent pairs $(a_0, a_1)$.  The
cutoff sweep (Sub-Sprint~2a) extends the symbol bit-exactly to
$n_{\max} \in \{1, 2, 3, 4\}$ via the closed forms
$M_1(n) = 2n(n{+}1)(n{+}2)/3$,
$M_2(n) = n(n{+}1)(n{+}2)(2n{+}1)(2n{+}3)/10$,
$M_3(n) = n(n{+}1)^2(n{+}2)/2$, all re-derived from the
chirality-balanced shell degeneracy $2n(n{+}1)$ and the CH
eigenvalue $n{+}1/2$.  See
Paper~55~\S~subsec:open\_m2\_m3 for the
Stage-1~Q5'~framing.  Source:
\texttt{debug/sprint\_q5p\_jlo\_nmax2\_memo.md},
\texttt{debug/sprint\_q5p\_2a\_jlo\_nmax\_sweep\_memo.md}.
\end{remark}

\begin{remark}[Mellin extraction points of the master-Mellin
partition (Sprint Q5'-Stage1-2b, 2026-06-05)]
\label{rem:q5p_mellin_extraction_points}
The master Mellin partition $(M_1, M_2, M_3)$ corresponds to three
categorically different extraction points of the formal Mellin
transform of the JLO degree-0 cochain on a compact 3-dimensional
spectral triple:
\begin{itemize}
\item $M_1$ at the spectral-dimension pole $s = d/2 = 3/2$ of
$\Gamma(s)\,\zeta_{D^2}(s)$:\ the Weyl-law short-time leading
$\mathrm{Tr}(e^{-tD^2}) \sim \sqrt{\pi}/2 \cdot t^{-3/2}$ supplies
the Hopf-base measure $\pi$ residue.
\item $M_2$ at integer-$s$ regular points of the same
$\Gamma(s)\,\zeta_{D^2}(s)$:\ Seeley--DeWitt pure-Tate continuum
values (e.g., $\zeta_{D^2}(2) = \pi^2 - \pi^4/12$ on unit~CH)
reproduce the Sprint~Q5'-CH-2 panel bit-exactly.
\item $M_3$ at the analogous extraction on a structurally
different Mellin complex, the $\eta$-style pairing
$\Gamma(s)\,\mathcal{M}[\mathrm{Tr}(\gamma\,D\,e^{-tD^2})](s)$:\
vertex-parity Hurwitz content at quarter-integer shifts.
\end{itemize}
At finite $n_{\max}$ the spectral zeta is bit-exact rational
(e.g.,
$\zeta_{D^2}^{(2)}(1) = 1\,978\,465\,196\,515\,232/532\,976\,619\,280\,625$
for the full Dirac at $n_{\max} = 2,\ \kappa = -1/16$); the
$s \to 0$ residue is exactly $\dim \mathcal{H} \in \mathbb{Z}$ (no~$\pi$).
Master-Mellin transcendentals enter only in the $n_{\max} \to
\infty$ continuum limit; the cosmic-Galois Stage~2 motivic-Galois
action acts at continuum, not at finite cutoff.  Source:
\texttt{debug/sprint\_q5p\_2b\_phi0\_mellin\_memo.md}.
\end{remark}

\begin{remark}[Continuum residue identification of the master-Mellin
$M_1$ Hopf-base coefficient (Sprint Q5'-Stage1-2b-continuum, 2026-06-05)]
\label{rem:q5p_continuum_residue}
The Mellin extraction-points reading of
Remark~\ref{rem:q5p_mellin_extraction_points} is closed at theorem
grade for the $M_1$ component on the truncated Camporesi--Higuchi
spectral triple of dimension $d = 3$.  The continuum spectral zeta
$\zeta_{D^2}^{\mathrm{cont}}(s)
   = 2\,\zeta(2s - 2, 3/2) - \tfrac{1}{2}\,\zeta(2s, 3/2)$
has a simple pole at the spectral-dimension value $s = d/2 = 3/2$
with bit-exact meromorphic residue $1$ (verified two independent
ways: from the Hurwitz pole at $u = 2s - 2 = 1$ with residue $1$,
and from direct Laurent expansion via the digamma identity).  The
Mellin residue of $\Gamma(s)\,\zeta_{D^2}^{\mathrm{cont}}(s)$ at
$s = 3/2$ is therefore $\Gamma(3/2) \cdot 1 = \sqrt{\pi}/2$, which
matches the standard heat-kernel Weyl-law leading coefficient
$r\,\mathrm{Vol}(S^3)/(4\pi)^{3/2}\big|_{r=2}$
(\cite{gilkey1995, vassilevich2003}) bit-exactly.  The three
published instances of the $M_1$ Hopf-base measure
($\mathrm{Vol}(S^2)/4 = \pi$ in Paper~18~\S~III.2;
$\mathrm{Vol}(S^2)/\pi^2 = 4/\pi$ in Paper~38~\S~VIII as the L2
asymptote rate constant; $\Gamma(d/2) = \sqrt{\pi}/2$ here at the
Mellin pole) are sibling normalizations of the same spectral object,
related by exact factors of $\pi^2/4$ and $\pi^{3/2}/8$ with no
spectral-data input.  The M2 continuum panel at integer
$s \in \{1, 2, 3, 4, 5\}$ reproduces the Sprint~Q5'-CH-2
Seeley--DeWitt panel bit-exactly (5/5 matches).  Tauberian rate
uniformity in a neighborhood of $s = 3/2$ is the Stage-2-relevant
open question, reachable from Karamata~\cite{karamata1962} and
Korevaar~\cite{korevaar2004}.  Source:
\texttt{debug/sprint\_q5p\_continuum\_mellin\_memo.md}.
\end{remark}

\begin{remark}[$\mathrm{HP}^{\mathrm{even}}$ class identification on
the bicomplex of the truncated CH triple (Sprint Q5'-Stage1-2c-Bicomplex, 2026-06-05)]
\label{rem:q5p_hp_even_class}
The cyclic $(b, B)$ bicomplex of the truncated Camporesi--Higuchi
spectral triple at $n_{\max} = 2$ has been constructed
bit-exactly.  The load-bearing degree-1 entire-cyclic cocycle
condition $b\phi_0 + B\phi_2 = 0$ holds bit-exactly on the full
panel of $36$ idempotent-pair inputs at three $t$-orders
$t^0, t^1, t^2$, both even and odd flavors ($216$ bit-exact zero
residuals).  The explicit $\mathrm{HP}^{\mathrm{even}}$ class on
the Morita-trivial baseline $\mathbb{Q}^5 = \mathrm{HP}_0(\mathbb{C}^5)$ is
\[
   [\phi_{\mathrm{JLO}}]^{\mathrm{HP}^{\mathrm{even}}}
   \;=\; (+2, -2, +2, +2, -4) \in \mathbb{Q}^5,
\]
the sector-resolved McKean--Singer index summing to
$\mathrm{Tr}(\gamma) = 0 = \chi(S^3)$.  The Sub-Sprint~1 symbol
$(16, 84, 36)$ is therefore a JLO cocycle representative, not a
coboundary.  A structural finding (honest scope):\ at degree~3,
$(b\phi_2 + B\phi_4)$ has bit-exact residual
$\pm 1/(3 \cdot 2^{16}) = \pm 1/196\,608$ on specific palindromic
$4$-tuples, a JLO-tower truncation artifact tied to the structural
vanishing of odd-degree cochains $\phi_1 \equiv \phi_3 \equiv 0$
on commutative $\mathcal{A}$.  The CM-residue cocycle for $M_3$ is
single-cochain and immune to this artifact;\ the JLO/CM-residue
distinction is Stage-2-relevant.  Source:
\texttt{debug/sprint\_q5p\_2c\_bicomplex\_memo.md}.
\end{remark}

\begin{remark}[CM-residue $\eta$-class identification, Sprint Q5'-Stage1-CM-Bicomplex,
June 2026]
\label{rem:q5p_cm_residue_class}
The second span of the Stage-1 cosmic-Galois $U^*$ bridge --- the
Connes--Moscovici $\eta$-pairing residue cocycle on the truncated
Camporesi--Higuchi spectral triple at $n_{\max} = 2$ --- has been
constructed bit-exactly in $\mathsf{sympy.Rational}$
(\texttt{debug/sprint\_q5p\_cm\_bicomplex\_memo.md},
\texttt{debug/data/sprint\_q5p\_cm\_bicomplex.json}). The CM-$\eta$
cochain $\psi_n^{\eta}(a_0, \dots, a_n; t) := \int_{\Delta_n}
\mathrm{Tr}(\gamma D\,a_0\,e^{-s_0 t D^2}\,[D, a_1]\cdots[D, a_n]\,
e^{-s_n t D^2})\,ds$ replaces the JLO leftmost insertion $\gamma$ by
$\gamma D$, and shares the bicomplex structure $(b + B)\psi = 0$ with
JLO via the same formulas. The degree-1 cocycle condition
$b\psi_0^{\eta} + B\psi_2^{\eta} = 0$ holds bit-exactly on the full
panel of 36 idempotent-pair inputs at three $t$-orders both flavors
(216 bit-exact zero residuals), with NO truncation artifact of the
JLO $\pm 1/196608$ type. The CM-$\eta$ class on the Morita-trivial
baseline $\mathbb{Q}^5 = \mathrm{HP}_0(\mathbb{C}^5)$ is explicitly
\[
[\psi_{\mathrm{CM}}^{\eta, \mathrm{even}}]^{\mathrm{HP}_{\eta}}
= (3, 3, 5, 15, 10) \in \mathbb{Z}^5,
\]
summing to $36 = M_3(n_{\max} = 2)$ (the Sub-Sprint 1 master-Mellin
M3 ground truth). The sector-resolved structure is the
chirality-symmetrized eigenvalue magnitude per sector
$\eta_s = \dim_s \cdot (n_s + 1/2)$, in clean dual to the JLO
$\chi_s = \dim(\gamma_+ e_s \mathcal{H}) - \dim(\gamma_- e_s \mathcal{H})$
chirality balance per sector. The two classes share the same $K_0$
sector-decomposition $\{[e_s]\}_{s=0..4}$ but encode distinct
cohomological content: JLO captures the zeroth heat-kernel
supertrace coefficient (McKean--Singer index $\chi(S^3) = 0$),
CM-$\eta$ captures the leading $t^0$ $\eta$-density supertrace
coefficient ($M_3$ mechanism leading). Cross-check at $n_{\max} = 3$:
$M_3(3) = 120$ bit-exact, CM-$\eta$ degree-1 cocycle bit-exact zero
on all 9 diagonal sectors. See Paper~55~\S~subsec:open\_m2\_m3 for
the Stage-1 Q5' construction context.
\end{remark}

\begin{remark}[Q5' Stage 1 pro-system functoriality of the JLO and CM-$\eta$
classes (Sprint Q5'-Stage1-Prosystem, June 2026)]
\label{rem:q5p_prosystem_functoriality}
The Stage-1 cocycle classes form a strictly compatible inverse system
across the cutoff pro-system $\{\mathcal{T}_{n_{\max}}\}_{n_{\max} \ge 1}$
(\texttt{debug/sprint\_q5p\_prosystem\_memo.md},
\texttt{debug/data/sprint\_q5p\_prosystem.json}). Sector idempotents are
indexed by $(n, l)$ with $1 \le n \le n_{\max}$, $0 \le l \le n$,
giving $N(n_{\max}) = n_{\max}(n_{\max} + 3)/2$ sectors at each cutoff
(2, 5, 9, 14 at $n_{\max} = 1, 2, 3, 4$, with strict nesting); the
truncation $P_{n+1 \to n}: \mathcal{O}_{n+1} \to \mathcal{O}_n$ is
sector projection (algebra homomorphism), the algebraic Berezin
$B_{n \to n+1}$ is sector inclusion. Both the JLO $\mathrm{HP}^{\mathrm{even}}$
class $\chi_{(n, l)}$ and the CM-$\eta$ class $\eta_{(n, l)}$ are
\emph{sector-local} --- the per-sector value depends only on the
sector label, not on the cutoff:
\[
   \chi_{(n, l)} = \begin{cases} +2 & l < n \\ -2n & l = n \end{cases},
   \qquad
   \eta_{(n, l)} = \begin{cases} (2l + 1)(2n + 1) & l < n \\ n(2n + 1) & l = n \end{cases}.
\]
Bit-exact pull-back identity $P^*_{n+1 \to n}[\psi^{(n+1)}] = [\psi^{(n)}]$
holds across all three consecutive cutoff pairs for both classes (six
inverse-image identities, each with bit-zero residual); algebraic
Berezin compatibility holds with the same bit-exactness on the common
sectors. Total sums match the Sub-Sprint~2a polynomial closed forms
$M_1(n_{\max}) = 2 n_{\max}(n_{\max}+1)(n_{\max}+2)/3$,
$\sum_s \chi_s = 0$ (McKean--Singer index $\chi(S^3) = 0$ preserved
cutoff-uniformly), $\sum_s \eta_s = M_3(n_{\max})
= n_{\max}(n_{\max}+1)^2(n_{\max}+2)/2$ bit-exactly at all four
cutoffs. The pro-system is the cleanest possible inverse-system lift
of the cohomological Stage-1 data: strict at the class level
(no twist, no coboundary modification), consistent with the
sector-locality of the Camporesi--Higuchi shell decomposition.
The strict-strong-form question --- pro-system functoriality at the
level of \emph{cochain morphisms} in the entire-cyclic bicomplex (vs
the class level closed here) --- is sharpened to a bit-exact dichotomy
in Remark~\ref{rem:q5p_strict_strong_drift};\ the multi-year NCG-mathematics
sub-question on the published Connes 1994 Ch.~IV / Loday \S~1.4 gap
remains open, independent of the Stage-1 closure here.
See Paper~55~\S~subsec:open\_m2\_m3 for the Q5' Stage 1 construction
context.
\end{remark}

\begin{remark}[Q5' Stage 2 candidate Hopf algebra substrate, Sprint Q5'-Stage2-Hopf, June 2026]
\label{rem:q5p_stage2_hopf_substrate}
The first scoping step of Stage 2 of the cosmic-Galois $U^*$ program produces
a concrete bit-exact finite-cutoff candidate Hopf algebra
$\mathcal{H}_{\mathrm{GV}}^{(n_{\max})} = \mathrm{Sym}_{\mathbb{Q}}(V_{n_{\max}})$
on the pro-system substrate of Remark~\ref{rem:q5p_prosystem_functoriality},
where $V_{n_{\max}}$ is the $\mathbb{Q}$-vector space of primitive generators
$x_{(n, l), k}$ indexed by $(n, l) \in \mathcal{S}_{n_{\max}}$ and Mellin slot
$k \in \{0, 1, 2\}$ (with $k = 0, 1, 2$ corresponding to M1, M3, M2 of
\S\ref{sec:case_exhaustion}, respectively). With primitive coproduct
$\Delta(x_g) = x_g \otimes 1 + 1 \otimes x_g$, counit $\varepsilon(x_g) = 0$,
and antipode $S(x_g) = -x_g$ on generators (extended as algebra homomorphisms),
$\mathcal{H}_{\mathrm{GV}}$ satisfies all five Hopf axioms bit-exactly at
$n_{\max} \in \{2, 3\}$ (269 bit-exact zero residuals on the axiom panel:\
52 coassociativity, 104 counit-left/right, 104 antipode-left/right, 9
bialgebra-compatibility), and the $k$-grading is preserved by $\Delta$ so that
$\mathcal{H}_{\mathrm{GV}} = \mathcal{H}_{\mathrm{GV}}^{[0]} \otimes
\mathcal{H}_{\mathrm{GV}}^{[1]} \otimes \mathcal{H}_{\mathrm{GV}}^{[2]}$
factors as a tensor product of three Hopf sub-algebras (42 bit-exact
generator-level $k$-preservation checks). The pro-system truncation
$P_{n+1 \to n}$ of Remark~\ref{rem:q5p_prosystem_functoriality} lifts
bit-exactly to a Hopf-algebra homomorphism
$\mathcal{H}_{\mathrm{GV}}^{(n+1)} \to \mathcal{H}_{\mathrm{GV}}^{(n)}$ (126
bit-exact zero residuals on the truncation panel at
$(n_{\max}, n_{\max} - 1) \in \{(2, 1), (3, 2)\}$). The candidate motivic
Galois group at finite cutoff is the additive affine group
$U^{*(n_{\max})}_{\mathrm{GeoVac}} = \mathbb{G}_a^{3 N(n_{\max})}$
(dimensions 15 and 27 at $n_{\max} = 2, 3$), with the Mellin-slot
factorisation $U^{*(n_{\max})}_{\mathrm{GeoVac}} = \prod_{k = 0}^{2}
\mathbb{G}_a^{N(n_{\max})}$ structurally enforcing the M1/M2/M3 partition
of Theorem~\ref{thm:pi_source_case_exhaustion} at the Stage-2 cosmic-Galois
level. The candidate is the \emph{abelian primitive-cospan} of the
Connes--Kreimer Hopf algebra of Feynman graphs (Connes--Kreimer 1998,
arXiv:hep-th/9808042):\ same generator-and-grading shape, same pro-system
truncation structure, but the sub-graph coproduct of CK is replaced by
the primitive coproduct here because the CH Fock sector idempotents are
sector-disjoint (no proper sub-sector nesting). The natural multi-year
continuation toward Stage-2's full Tannakian construction
(Connes--Marcolli 2007, arXiv:math/0409306) is to enrich the coproduct
with a sub-sector or nested-Hopf-tower structure that introduces
non-abelian pro-unipotent content. See Paper~55~\S~subsec:open\_m2\_m3
for the cosmic-Galois narrative.
\end{remark}

\begin{remark}[Q5' Stage 1 strict-strong-form cochain-morphism drift,
Sprint Q5'-Stage1-StrictStrong, June 2026]
\label{rem:q5p_strict_strong_drift}
Sharpening the v3.60.0 honest-scope caveat of
Remark~\ref{rem:q5p_prosystem_functoriality}:\ while pro-system
functoriality at the level of cocycle CLASSES is bit-exactly strict
across the cutoff pro-system $\{\mathcal{T}_{n_{\max}}\}_{n_{\max} \ge 1}$
on commutative $\mathcal{A}$, pro-system functoriality at the
COCHAIN-MORPHISM level of the entire-cyclic bicomplex FAILS bit-exactly
on the same algebra (\texttt{debug/sprint\_q5p\_strict\_strong\_memo.md},
\texttt{debug/data/sprint\_q5p\_strict\_strong.json}). The degree-3
closure residual $(b\phi_2 + B\phi_4)$ on the load-bearing $(e_2, e_3)$
palindromic 4-tuples carries bit-exact values
\[
n_{\max} = 2: \pm 1/(3 \cdot 2^{16})
\quad\to\quad
n_{\max} \ge 3: \pm 1/2^{16}
\quad(\text{bit-exact fixed point at } n_{\max} = 3, 4),
\]
with the pull-back cochain increment $\Delta(B\phi_4)_{3 \to 2}
= -1/(3 \cdot 2^{14})$ closing the residual identity bit-exactly:\
$1/196608 + (-1/49152) = -1/65536$. The CM-$\eta$ cocycle tower
exhibits the same structural pattern at degree 3 with structurally
distinct quantitative values (the v3.59.0 reading that the CM-$\eta$
tower is the cochain-clean alternative is restricted to degree 1 and
does not extend to degree 3). The pro-system functoriality verdict is
therefore a sharp dichotomy:\ strict at the class level
(Remark~\ref{rem:q5p_prosystem_functoriality}), bit-exact non-strict
with $n_{\max} = 3$ fixed-point structural drift at the cochain-morphism
level (this remark). The drift is consistent with a boundary-vs-interior
shell partition:\ $n_{\max} = 2$ is the boundary cutoff (no shell above
$n = 2$ in the $\phi_4$ matrix-element channel), $n_{\max} \ge 3$ are
interior cutoffs (stabilized). Empirical confirmation at bit-exact
precision of the published-open Connes 1994 Ch.~IV / Loday \S 1.4 gap
on entire-cyclic functoriality under algebra truncations. See Paper~55
\S~subsec:open\_m2\_m3 for the Stage-1 Q5' construction context.
\end{remark}

\begin{remark}[Tauberian rate uniformity in a neighborhood of $s = 3/2$
(Sprint Q5'-Stage1-Tauberian, 2026-06-05)]
\label{rem:tauberian_uniformity}
The pointwise rate of approach in
Remark~\ref{rem:q5p_continuum_residue} is sharpened to a uniform rate
on compact above-pole neighborhoods of the spectral-dimension pole
$s = d/2 = 3/2$ via the standard uniform Tauberian theorem of
Karamata~\cite{karamata1962} (foundational), Korevaar~\cite{korevaar2004}
\S~III.4 (modern treatment with boundary-saturation Remark~3), and
Tenenbaum~\cite{tenenbaum2015} \S~II.7 Thm~II.7.7 (uniform rates of
approach with explicit constants). The four hypotheses transport
bit-exactly to the CH spectral zeta:\ (H1) positivity $a_n = 2n(n+1) > 0$;\
(H2) polynomial growth $a_n = O(n^2)$;\ (H3) meromorphic continuation
with unique simple pole at $\sigma_0 = 3/2$ of residue $1$ (the v3.59.0
Track 2 closure);\ (H4) vertical-strip boundedness from standard Hurwitz
convexity. Conclusion:\ on the compact above-pole neighborhood
$U_{>} = [3/2 + \epsilon,\, 3/2 + 1/4]$,
$\sup_{s \in U_{>}} |\zeta^{(N)}_{D^2}(s) - \zeta^{\mathrm{cont}}_{D^2}(s)|
= O(N^{-2\epsilon + \delta})$ uniformly in $U_{>}$, for any $\delta > 0$.
Boundary-saturation:\ the supremum is attained at the boundary point
closest to the pole, verified bit-exactly at $n_{\max} \in \{2, 3, 4, 5\}$
on the six-grid panel $\{5/4, 11/8, 23/16, 25/16, 13/8, 7/4\}$ (sup
attained at $s = 25/16$ at every cutoff;\ empirical slope $-0.092$ at
$n_{\max} \in \{2,\ldots,5\}$ monotonically approaches the predicted
$-1/8$, hitting $-0.123$ at $n_{\max} = 100$). Source:\
\texttt{debug/sprint\_q5p\_tauberian\_memo.md}.
\end{remark}

\begin{remark}[Continuum residue identification of the master-Mellin $M_3$
$\eta$-tower (Sprint Q5'-Stage1-M3-Continuum, 2026-06-05)]
\label{rem:q5p_m3_continuum_residue}
The continuum $M_3$ $\eta$-Mellin object on the truncated
Camporesi--Higuchi spectral triple of dimension $d = 3$ is closed at
theorem grade with parity-of-$s$ stratification. The continuum
$\eta$-Dirichlet series
$\eta_D(s) := \sum_{n \ge 0} g_n |\lambda_n|^{1 - s}$ admits the
Hurwitz reduction $\eta_D(s) = 2 \zeta(s - 3, 3/2) - \tfrac{1}{2} \zeta(s - 1, 3/2)$
(equivalently $\eta_D(s) = D(s - 1)$ via Paper~28's Dirac Dirichlet
series), with TWO simple poles at $s = 4$ and $s = 2$ of bit-exact
meromorphic residues $2$ and $-1/2$ respectively (verified two
independent routes:\ Hurwitz pole rule at $u = 1$ AND direct Laurent
via digamma). The corresponding Mellin residues
$\mathrm{Res}_{s = 4} \Gamma(s) \eta_D(s) = 12$ and
$\mathrm{Res}_{s = 2} \Gamma(s) \eta_D(s) = -1/2$ are rational integer
values, NOT in the M3 cyclotomic ring;\ the genuine M3 cyclotomic
content lives in the CHIRALITY-RESOLVED $\eta$-Mellin
$(D_{\mathrm{even}} - D_{\mathrm{odd}})(s - 1)$, which inherits the
Sprint Q5'-CH-3 even-$s$ M3 stratification via the
$\Q(i) = \Q(\zeta_4)$ Galois descent. Specifically at $s = 3$ the
chirality discriminant equals $-1 + 2G$ (Catalan $G$;\ genuine M3
depth 1);\ at $s = 5$ it equals $8 \beta_4 - 8G$;\ at $s = 7$ it equals
$-32 \beta_4 + 32 \beta_6$ (both genuine M3 depth $1+$). Odd-$(s - 1)$
values collapse to $M_1$ via $\beta(2k+1) \in \pi^{\mathrm{odd}} \Q$.
The discrete-side Karamata Tauberian signature at the $s = 4$ pole
holds exactly:\ partial sums on the v3.60.0 sector-local substrate
satisfy $S(n_{\max}) / \log(n_{\max}) \to 2$ at $n_{\max} \ge 200$
(meromorphic residue;\ no Jacobian since the natural $\eta$-Mellin
variable IS the spectral $|\lambda| = n + 1/2$). This closes the $M_3$
component of the continuum-residue identification triad started by
Remark~\ref{rem:q5p_continuum_residue}. Structural sharpening:\ unlike
$M_1$ (which admits three exact-factor sibling normalizations $\pi$,
$4/\pi$, $\sqrt{\pi}/2$ of a single generator), $M_3$'s ring
$\mathcal{MT}(\Z[i, 1/2], 4)$ is depth-graded with generators
$\{1, G, \beta_4, \beta_6, \ldots\}$ at successive depths with no
$\pi$-power reduction at any depth $\ge 1$ (Apéry-style irrationality
of Catalan $G$ vs $\pi^{2k} \Q$). \emph{This asymmetry IS the
operational content of the master-Mellin partition} $M_1$ at $k = 0$
depth $0$ vs $M_3$ at $k = 1$ depth $1+$. Source:\
\texttt{debug/sprint\_q5p\_m3\_continuum\_memo.md}.
\end{remark}

\begin{remark}[Q5' Stage 2 substrate enrichment via $J^*(S^3)$, Sprint Q5'-J-Star-S3, June 2026]
\label{rem:q5p_j_star_substrate}
A first scoping step of the multi-year nested-Hopf-tower substrate
enrichment of Remark~\ref{rem:q5p_stage2_hopf_substrate} replaces the
disjoint sector-idempotent substrate $\mathbb{C}^{N(n_{\max})}$ with
the SU(2) Peter--Weyl matrix-coefficient substrate
\[
J^*_{j_{\max}}(S^3) = \operatorname{span}_{\Q}\{\pi^j_{mn} : 0 \le j \le j_{\max},\; -j \le m, n \le j\},
\]
with the matrix-coefficient coproduct
$\Delta \pi^j_{mn} = \sum_p \pi^j_{mp} \otimes \pi^j_{pn}$.
This coproduct is bit-exactly \emph{non-primitive} on every $j > 0$
generator (verified at $j_{\max} \in \{1/2, 1, 3/2\}$ on 29 generators
with $j > 0$);\ coassociativity and counit hold bit-exactly in the free
polynomial algebra (147 zero residuals);\ the antipode axiom requires
passing to the standard $\mathcal{O}(SU(2)) \cong \mathcal{O}(SL_2)$
quotient by SU(2) unitarity relations (textbook construction;\
Klimyk--Schmüdgen 1997 §1.3.2);\ the pro-system truncation
$P_{j_{\max} + 1/2 \to j_{\max}}$ lifts bit-exactly to a Hopf-algebra
homomorphism (132 zero residuals). The candidate motivic Galois group
at the quotient is $U^{*(j_{\max})}_{\mathrm{GeoVac}, J^*} = SL_2$ for
every $j_{\max} \ge 1/2$ --- a non-abelian semisimple algebraic group,
achieving the qualitative goal of breaking the abelian primitive shape
of Remark~\ref{rem:q5p_stage2_hopf_substrate}. The trade is
axis-orthogonal:\ $J^*(S^3)$ has non-abelian content but sacrifices the
Mellin-slot $k$-partition;\ v3.61.0's $\mathcal{H}_{\mathrm{GV}}$ has
the Mellin partition but is abelian. The natural multi-year continuation
is a synthesis combining both --- either (a) tensor product
$\mathcal{H}_{\mathrm{v3.61}} \otimes \mathcal{H}^{J^*}$ giving
$U^* = \mathbb{G}_a^{3 N} \times SL_2$, or (b) Mellin-slot-decorated
Peter--Weyl $\pi^{j, k}_{mn}$ giving $U^* = SL_2^3$. The
$SL_2 \times \mathbb{G}_a^d$ form would be the Levi decomposition of a
general algebraic group (semisimple times pro-unipotent), suggesting
the synthesis substrate is the right Stage-2 target. See
Paper~55~\S~subsec:open\_m2\_m3.
\end{remark}

\begin{remark}[Q5' Stage 2 cross-shell off-diagonal Dirac enrichment scoping,
Sprint Q5'-OffDiag-Dirac, June 2026]
\label{rem:q5p_offdiag_dirac_enrichment}
The first multi-year scoping step of the cross-shell off-diagonal Dirac
perturbation enrichment (second of three structural ingredients flagged
in Remark~\ref{rem:q5p_stage2_hopf_substrate}) promotes the v3.61.0
pro-system substrate to track the off-diagonal Dirac perturbation
$\kappa A$ via transition generators
$T_{s' \to s} := e_s \cdot (\kappa A) \cdot e_{s'}$ (non-zero for
E1-Gaunt-selected sector pairs). At $n_{\max} = 2$ on the truncated CH
spectral triple, the enriched substrate breaks the sector-locality
structural condition that v3.61.0 Track A identified as forcing the
abelian primitive coproduct:\ bit-exact in $\mathsf{sympy.Rational}$,
12 of 12 single-step transitions carry non-vanishing $\eta$-class
values in $\kappa^2 \cdot \mathbb{Z}$ (the M3 cocycle slot non-vanishes
on inter-sector content), 12 of 12 chain compositions land on
idempotents $T_2 \cdot T_1 = c \cdot e_s$ with non-zero
$2 \eta(T_1) \eta(T_2)$ cross-term cocycle image, and an additional 18
chain compositions land on two-step transitions outside the basic
single-step generator basis (algebra-closure failure). The lowest-order
non-trivial commutators $[T_a, T_b] \ne 0$ appear at 24 of 66 = 36\%
of transition pairs --- explicit non-abelian Lie generators. The
non-primitivity content has denominator family $\kappa^2 \cdot \mathbb{Z}
= 2^{-8} \cdot \mathbb{Z}$ times inverse path counts, structurally
aligned with the bit-exact $\pm 1/65536$ closure drift residual of
Remark~\ref{rem:q5p_strict_strong_drift} modulo a $2^3$ JLO simplex
normalization factor;\ the exact bit-exact bridge identification
between the two phenomena is a feasible follow-on. The enriched
substrate's candidate motivic Galois group
$U^{*(n_{\max})}_{\mathrm{GeoVac}, \mathrm{enriched}}$ has both an
abelian (idempotent-generated) part and a unipotent (transition-generated)
part, structurally targeting the pro-unipotent factor of
Connes--Kreimer--Marcolli's cosmic-Galois $U^*$ (arXiv:hep-th/9808042;\
arXiv:math/0409306) at the Stage-2 level. See
Paper~55~\S~subsec:open\_m2\_m3 for the cosmic-Galois narrative.
\end{remark}

\begin{remark}[Q5' Stage 2 Cuntz-extension enrichment:\ structural incompatibility,
Sprint Q5'-Cuntz, June 2026]
\label{rem:q5p_stage2_cuntz_incompat}
The first scoping step of the multi-year Cuntz-extension enrichment
ingredient flagged in Remark~\ref{rem:q5p_stage2_hopf_substrate} returns
a structural negative. Replace the commutative algebra
$\mathcal{A}_{\mathrm{GV}}^{(n_{\max})} = \mathbb{C}^{N(n_{\max})}$ of
sector idempotents by the Cuntz algebra $\mathcal{O}_N$ generated by
$N = 3 \cdot N(n_{\max})$ isometries $S_{(n, l), k}$ (one per sector
$(n, l)$ and Mellin slot $k \in \{0, 1, 2\}$). The two natural candidate
coproducts on $\mathcal{O}_N$ are bit-exactly incompatible with the
Cuntz relations:\ (i) the primitive coproduct
$\Delta(S_i) = S_i \otimes 1 + 1 \otimes S_i$ extended as a
$*$-algebra homomorphism fails $S_i^* S_j = \delta_{ij}$ on $0/225$
pairs at $n_{\max} = 2$ with $N = 15$;\ (ii) the diagonal coproduct
$\Delta(S_i) = S_i \otimes S_i$ passes the first Cuntz relation on
$225/225$ pairs but fails $\sum_i S_i S_i^* = 1$ by exactly
$N (N - 1) = 210$ missing cross-terms. Net:\ $\mathcal{O}_N$ is
structurally a Hopf-\emph{comodule} (over the gauge $U(1)$ scaling
each $S_i$ by a phase) rather than a Hopf-algebra. The M1/M2/M3
partition survives only at the trivial comodule level via fixed-$k$
generator labels, structurally weaker than the clean Hopf tensor
product $\mathcal{H}_{\mathrm{GV}} = \otimes_k \mathcal{H}_{\mathrm{GV}}^{[k]}$
established for the commutative case in
Remark~\ref{rem:q5p_stage2_hopf_substrate}. The JLO degree-1 cochain
$\phi_1$ is structurally non-zero on $\mathcal{O}_N$ (the
commutative-case vanishing theorem requires sector orthogonality
$e_s e_t = \delta_{st} e_s$ which Cuntz isometries violate), but
extraction at theorem grade requires the full CH-adjacency-extended
Dirac representation through the Pimsner--Voiculescu sequence
(Pimsner--Voiculescu 1980, \emph{J. Operator Theory} 4;\ Pimsner 1997,
Fields Inst. Commun. 12), genuinely multi-year. The Cuntz route is
therefore the structurally most disruptive of the three enrichment
ingredients flagged in Remark~\ref{rem:q5p_stage2_hopf_substrate}:\
it does not enrich the Hopf substrate but shifts the framework into
the categorically different Hopf-comodule-over-gauge-group setting
(Cuntz--Pimsner extension). The remaining two ingredients
(nested-Hopf-tower $J^*(S^3)$ closed at finite cutoff in
Remark~\ref{rem:q5p_j_star_substrate};\ cross-shell off-diagonal Dirac
closed at finite cutoff in Remark~\ref{rem:q5p_offdiag_dirac_enrichment})
preserve the Hopf-algebra category at the substrate level and are
correspondingly more tractable. See Paper~55~\S~subsec:open\_m2\_m3 for
the cosmic-Galois program narrative.
\end{remark}

\begin{remark}[Q5' Stage 2 substrate HEADLINE:\ Levi-decomposition tensor synthesis,
Sprint Q5'-Levi-Synthesis, June 2026]
\label{rem:q5p_levi_synthesis_substrate}
The HEADLINE Stage-2 substrate construction combines the v3.61.0 abelian
primitive substrate of Remark~\ref{rem:q5p_stage2_hopf_substrate} with the
v3.62.0 Peter--Weyl substrate of Remark~\ref{rem:q5p_j_star_substrate} into
the tensor-product Hopf algebra
\[
   \mathcal{H}_{\mathrm{GV}}^{\mathrm{Levi}, (n_{\max}, j_{\max})}
   := \mathcal{H}_{\mathrm{GV}}^{(\mathrm{v3.61}), n_{\max}}
      \otimes_{\Q} \mathcal{H}_{\mathrm{GV}}^{J^*, j_{\max}}.
\]
At finite cutoff $(n_{\max}, j_{\max}) \in \{(2, 1/2), (2, 1), (3, 1/2)\}$,
all five Hopf axioms (coassociativity, counit, antipode), bialgebra
compatibility, the M1/M3/M2 Mellin-slot $k$-grading preservation on the
first factor, and the non-cocommutativity on the second factor (encoding
non-abelian $SL_2$ content) are verified bit-exactly (882 zero residuals
across the joint axiom + truncation panel:\ 636 axiom residuals + 246
truncation residuals). The Mellin partition lifts to
$\mathcal{H}^{\mathrm{Levi}} = \bigotimes_{k = 0}^{2}
\mathcal{H}_{\mathrm{v3.61}}^{[k]} \otimes \mathcal{H}^{J^*}$. The
candidate motivic Galois group at the standard $\mathcal{O}(SL_2)$
quotient is
\[
   \boxed{
   U^{*(n_{\max}, j_{\max})}_{\mathrm{GeoVac}, \mathrm{Levi}}
   = \mathbb{G}_a^{3 N(n_{\max})} \times SL_2
   }
\]
by the Tannakian theorem (tensor product of Hopf algebras = direct
product of affine group schemes;\ Waterhouse 1979 §1.4;\ Sweedler 1969
Ch.~IV;\ Deligne 1990 \emph{Cat\'egories tannakiennes}). The semidirect
product $\mathbb{G}_a^{3 N(n_{\max})} \rtimes SL_2$ reduces to a direct
product because the Levi action on the unipotent radical is trivial:\
the v3.61.0 generators $x_{(n, l), k}$ carry no $SU(2)$ representation
content (they are abelian primitives indexed by sector and Mellin slot),
so any $SL_2$ element acts on them as the identity (Hochschild 1981
Ch.~VII Levi-decomposition direct-product reduction). The pro-unipotent
$\mathbb{G}_a^{3 N(n_{\max})}$ factor carries the master-Mellin engine
M1/M3/M2 partition (Remark~\ref{rem:q5p_stage2_hopf_substrate});\ the
semisimple $SL_2$ factor carries the Peter--Weyl non-abelian content
(Remark~\ref{rem:q5p_j_star_substrate}). The Levi-decomposition shape is
the published structural target for the cosmic-Galois $U^*$ in the
Connes--Marcolli motivic Galois machinery (Connes--Marcolli 2007
arXiv:math/0409306;\ Connes--Marcolli 2008 book Ch.~4):\ pro-unipotent
times semisimple. The multi-year Stage-2 Tannakian construction reduces
to (a) full Tannakian closure of the pro-unipotent factor, and (b)
verification that the cosmic-Galois $U^*$ of Connes--Marcolli 2007
acts on $\mathcal{H}^{\mathrm{Levi}}_{\mathrm{GV}}$ in the expected way.
See Paper~55~\S~subsec:open\_m2\_m3.
\end{remark}

\begin{remark}[Q5' Stage 2 substrate enrichment via Mellin-slot-decorated Peter--Weyl,
Sprint Q5'-Decorated-PW, June 2026]
\label{rem:q5p_decorated_pw}
The alternative synthesis flagged in
Remark~\ref{rem:q5p_j_star_substrate} option (b) --- the
Mellin-slot-decorated Peter--Weyl substrate
\[
   \mathcal{H}_{\mathrm{dec}}^{(j_{\max})}
   = \mathrm{span}_{\Q}\{\pi^{j, k}_{mn} : 0 \le j \le j_{\max},\;
     -j \le m, n \le j,\; k \in \{0, 1, 2\}\}
\]
with $k$-preserving coproduct $\Delta \pi^{j, k}_{mn}
= \sum_p \pi^{j, k}_{mp} \otimes \pi^{j, k}_{pn}$ --- is constructed and
verified bit-exactly at $j_{\max} \in \{1/2, 1\}$ with $n_k = 3$ Mellin
slots. All Hopf axioms hold bit-exactly with the standard counit and
per-slot $SU(2) \to SL_2$ quotient (411 zero residuals:\ coassoc 57,
counit L+R 114, antipode at quotient 57, $k$-grading preservation 57,
pro-system Hopf-hom 126). The substrate factorises as
$\mathcal{H}_{\mathrm{dec}} = \mathcal{O}(SL_2)^{\otimes 3}$ with the
three Mellin mechanisms M1, M2, M3 carried by three independent
non-abelian semisimple $SL_2$ factors. The candidate motivic Galois
group at the $SU(2)^{\otimes 3}$ quotient is
$U^*_{\mathrm{dec}, (a)} = SL_2^3$ for every $j_{\max} \ge 1/2$. The
alternative $k$-summing coproduct
$\Delta \pi^{j, k}_{mn} = \sum_{k_1 + k_2 = k \pmod 3} \sum_p \pi^{j, k_1}_{mp} \otimes \pi^{j, k_2}_{pn}$
passes coassociativity and a $\mathbb{Z}/3$-augmented counit
$\varepsilon_{\mathrm{aug}}(\pi^{j, k}_{mn}) = \delta_{mn} \delta_{k, 0}$
but FAILS the antipode axiom at the per-slot $SU(2)$ quotient by a
structural factor of $n_k = 3$ --- fixing the normalisation requires
$1/n_k \in \Q(1/3)$ breaking the discrete-for-skeleton discipline. The
two synthesis routes (this remark's $SL_2^3$ vs
Remark~\ref{rem:q5p_levi_synthesis_substrate}'s tensor synthesis
$\mathbb{G}_a^{3 N} \times SL_2$) capture different aspects of the
Stage-2 substrate:\ the tensor synthesis is the Levi-decomposition
shape with explicit v3.61.0 abelian primitive content;\ this remark's
decorated-PW is the Mellin-decorated symmetry shape with the M1/M2/M3
partition realised as independent $SL_2$ factors. Both substrates are
valid candidates;\ the strategic choice between them is the next
sprint-scale decision in the multi-year Stage-2 program. See
Paper~55~\S~subsec:open\_m2\_m3.
\end{remark}

\begin{remark}[Q5' Stage 2 bridge identity closing the T3b umbrella,
Sprint Q5'-Bridge-Id, June 2026]
\label{rem:q5p_bridge_identity}
The bit-exact bridge identification flagged as a feasible follow-on in
Remark~\ref{rem:q5p_offdiag_dirac_enrichment} closes
(\texttt{debug/sprint\_q5p\_bridge\_id\_memo.md},
\texttt{debug/data/sprint\_q5p\_bridge\_id.json}). On the load-bearing
$(e_2, e_3, e_3, e_2)$ palindrome at the strict-strong-form
cochain-morphism interior cutoff $n_{\max} \ge 3$, the closure drift
residual is bit-exactly $-\kappa^4 = -1/2^{16}$, equivalent in four
closed forms:\
\[
   \mathrm{drift}_{n_{\max} \ge 3}
   = -\kappa^4
   = -\frac{1}{4!} \kappa^4 \, T_{\mathrm{path}}^{\mathrm{int}}
   = -\frac{2 \eta(T_{(2,0) \to (2,1)}) \eta(T_{(2,1) \to (2,0)})}{40}
   = -\frac{5/8192}{40},
\]
with $T_{\mathrm{path}}^{\mathrm{int}} = 24$ the chirality-weighted
two-step E1 path count at $n_{\max} \ge 3$ and the integer $40
= 2 \cdot 10 \cdot 2$ the palindromic two-step path-count product on
the OffDiag substrate $T_{(2, 0) \to (2, 1)}, T_{(2, 1) \to (2, 0)}$
of Remark~\ref{rem:q5p_offdiag_dirac_enrichment}. The boundary cutoff
$n_{\max} = 2$ value $+\kappa^4 / 3 = +1/196608$ uses
$T_{\mathrm{path}}^{\mathrm{bdy}} = 8$, with the bit-exact ratio
$T_{\mathrm{path}}^{\mathrm{int}} / T_{\mathrm{path}}^{\mathrm{bdy}}
= 24/8 = 3$ identifying the $-3$ boundary-vs-interior factor of
Remark~\ref{rem:q5p_strict_strong_drift} as an explicit integer
path-count ratio. The pullback increment closes the integer arithmetic:\
$32 = 24 + 8$. The denominator-scale mismatch $2^{13}$ (T3b
$\eta \otimes \eta$) vs $2^{16}$ (Track B drift) is bit-exactly the
$2^3$ piece of the JLO simplex factor $1/4! = 1/(3 \cdot 2^3)$;\ the
leftover numerator $5$ in $5/8192$ is half the path-count product
$40 = 5 \cdot 8$. The T3b umbrella conjecture (``the two findings are
aspects of one Stage-2 substrate-enrichment story with denominator
scales aligned modulo a $2^3$ JLO simplex factor'') is therefore
bit-exactly correct, and the bridge identity is constructive at finite
cutoff. Multi-year follow-ons (closed-form $T_{\mathrm{path}}(n_{\max})$,
Lie-algebra structure constants, continuum-limit Mellin lift) remain
open.
\end{remark}

\begin{remark}[Q5' Stage 2 OffDiag algebra-closure dimension at $n_{\max} = 3$,
Sprint Q5'-OffDiag-Closure-nmax3, June 2026]
\label{rem:q5p_offdiag_closure_nmax3}
The first multi-year scoping step of the algebra-closure dimension at
general $n_{\max}$ (named follow-on (1) of
Remark~\ref{rem:q5p_offdiag_dirac_enrichment}) identifies the closure of
the OffDiag-enriched substrate $\mathcal{A}_{\mathrm{OD}}^{(n_{\max})}
= \Q\langle e_s, A_{\mathrm{off}} \rangle$ under matrix multiplication as
the path algebra of the sector quiver realised in $M_{\dim H}(\Q)$ via
the canonical adjacency representation $A_{\mathrm{off}}
= D - \Lambda = \kappa A_{\mathrm{graph}}$. Bit-exact in
$\mathsf{sympy.Rational}$ at $n_{\max} \in \{2, 3\}$:\
$\dim \mathcal{A}_{\mathrm{OD}}^{(2)} = 84$ and
$\dim \mathcal{A}_{\mathrm{OD}}^{(3)} = 592$, with closure decomposing
as a sum of per-block ranks over all $N(n_{\max})^2$ sector pairs
(all 25 of 25 blocks at $n_{\max} = 2$, all 81 of 81 at $n_{\max} = 3$
realised non-trivially --- the sector quiver is fully connected via
multi-step E1 dipole paths). Fill-fractions $84/256 = 21/64$ and
$592/1600 = 37/100$ slowly increasing;\ asymptotic limit requires
$n_{\max} = 4$ for honest identification. Lie subalgebra of
single-step transitions does NOT close:\ 82 of 378 commutator pairs
lie outside the 28-dim transition span at $n_{\max} = 3$, generalising
the 24 of 66 at $n_{\max} = 2$. Named sprint-scale follow-on:\ extend
bit-exact closure to $n_{\max} = 4$ for third data point + asymptotic
fill-fraction identification. Multi-year:\ closed-form per-block rank
formula from quiver combinatorics, and embedding of
$\{\mathcal{A}_{\mathrm{OD}}^{(n_{\max})}\}_{n_{\max}}$ into the
inductive system whose limit hosts the conjectural pro-unipotent factor
of the cosmic-Galois $U^*$ Hopf algebra (Connes--Marcolli 2007,
arXiv:math/0409306). See Paper~55~\S~subsec:open\_m2\_m3.
\end{remark}

\begin{remark}[Q5' Stage-2 combined substrate L5 = L1 categorically,
Sprint Q5'-Combined-Substrate, June 2026]
\label{rem:q5p_combined_substrate_levi}
The first multi-year scoping step of combining the T3a Peter--Weyl
substrate (Remark~\ref{rem:q5p_j_star_substrate}) and T3b cross-shell
off-diagonal Dirac substrate
(Remark~\ref{rem:q5p_offdiag_dirac_enrichment}) on a SINGLE combined
substrate closes via categorical reduction. The combined substrate
$\mathcal{A}^{\mathrm{combined}} = \mathcal{A}^{J^*}_{j_{\max}} \otimes
\mathcal{A}^{\mathrm{OD}}_{n_{\max}}$ reduces to the categorical TENSOR
PRODUCT of Remark~\ref{rem:q5p_levi_synthesis_substrate} at the basic
substrate level, not a smash product. Two structural obstructions block
the smash-product upgrade:\ (i) Peter--Weyl matrix coefficients
$\pi^j_{mn}$ are polynomial \emph{functions} on $SL_2$, while OffDiag
transitions $T_{s' \to s}$ are operators on $\mathcal{H}$ with
\emph{scalar} cocycle data ($\eta$-class values in $\kappa^2 \cdot \Z$)
--- categorically different objects;\ no $SL_2$ group element evaluates
the 4 spin-$1/2$ matrix coefficients to the 12 OffDiag $\eta$-rationals
simultaneously. (ii) The natural $SU(2)$ z-rotation $U_\theta$ does NOT
preserve the basic 5-sector OffDiag basis:\ bit-exact symbolic
computation (\texttt{debug/compute\_q5p\_combined\_substrate\_part2.py})
gives three distinct phase factors $\{1, e^{i\theta}, e^{-i\theta}\}$ on
the conjugate $U_\theta T U_\theta^{-1}$ --- the result lives in the
span of $m_J$-resolved sub-transitions, not in the span of the basic
basis. The candidate motivic Galois group at the quotient is therefore
the Levi-decomposition product
\[
   U^{*(n_{\max}, j_{\max})}_{\mathrm{combined}}
   = SL_2 \times \mathbb{G}_a^{N_{\mathrm{OD}}(n_{\max})},
\]
exactly the multi-year synthesis target predicted by
Remark~\ref{rem:q5p_levi_synthesis_substrate}.
\textbf{L5 reduces to L1 at the substrate level};\ the genuinely new
content (a non-trivial smash product) lives at the $m_J$-resolved
refinement of the OffDiag basis, which loses the v3.61.0 sector
structure as load-bearing data and constitutes a multi-year continuation.
See Paper~55~\S~subsec:open\_m2\_m3.
\end{remark}

\begin{remark}[Q5' Stage-2 substrate decision: L1 is forced via the
primitive-space invariant, Sprint Q5'-L1-vs-L2-Diagnostic, June 2026]
\label{rem:q5p_l1_forced}
The strategic question flagged in
Remark~\ref{rem:q5p_combined_substrate_levi} and in Sprint Q5'-Levi-Arc
\S~5 -- whether the v3.61.0 abelian primitives $x_{(n, l), k}$ in
$\mathcal{H}_{\mathrm{GV}}^{(n_{\max})}$ are reducible to Peter--Weyl
matrix coefficients in disguise (so that the simpler $SL_2^3$ substrate
of Remark~\ref{rem:q5p_decorated_pw} would suffice) or carry independent
content (so that the Levi-decomposition substrate of
Remark~\ref{rem:q5p_levi_synthesis_substrate} is forced) -- is resolved
bit-exactly at finite cutoff by the primitive-space invariant
$\dim \mathrm{Prim}(\mathcal{H}) = \dim \mathrm{Lie}(\mathrm{Spec}\,\mathcal{H})$,
a Hopf-isomorphism invariant. For $\mathcal{H}_{\mathrm{GV}}^{(n_{\max})}$
the Milnor--Moore structure theorem for connected cocommutative Hopf
algebras over $\mathbb{Q}$ identifies $\mathrm{Prim}$ with the
generating vector space $V_{n_{\max}}$, giving
$\dim \mathrm{Prim}(\mathcal{H}_{\mathrm{GV}}) = 3 N(n_{\max}) \in \{6, 15, 27, 42, 60, 81\}$
at $n_{\max} \in \{1, \dots, 6\}$. For $\mathcal{H}_{\mathrm{dec}}^{(j_{\max})}$
the standard identification $\mathrm{Prim}\,\mathcal{O}(SL_2)^{\otimes 3} =
\mathfrak{sl}_2 \oplus \mathfrak{sl}_2 \oplus \mathfrak{sl}_2$ gives
$\dim \mathrm{Prim}(\mathcal{H}_{\mathrm{dec}}) = 9$ independent of
$j_{\max}$ (higher $j_{\max}$ adds symmetric-square Clebsch--Gordan
polynomials of degree $\ge 2$ in the $j = 1/2$ generators, contributing
zero new dimensions to $\mathfrak{m}/\mathfrak{m}^2$ at the augmentation
ideal). At every $n_{\max} \ge 2$ the dim defect is strictly positive
$(6, 18, 33, 51, 72)$, so no Hopf isomorphism exists between the two
substrates;\ L1 is forced. The substrate-dim coincidences flagged in
Remark~\ref{rem:q5p_decorated_pw} ($15 = 15$ at
$(n_{\max}, j_{\max}) = (2, 1/2)$ and $42 = 42$ at $(4, 1)$) are
small-cutoff artifacts breaking at $n_{\max} \in \{1, 3, 5, 6\}$.
An independent Hopf-isomorphism invariant (cocommutativity) separates
the substrates at every $n_{\max} \ge 1$ -- $\mathcal{H}_{\mathrm{GV}}$
is cocommutative (every generator primitive) while
$\mathcal{H}_{\mathrm{dec}}$ is not on $j > 0$ matrix coefficients.
Bit-exact panel:\ 258 zero residuals / consistent checks (driver
\texttt{debug/compute\_q5p\_l1\_vs\_l2\_diagnostic.py};\ data
\texttt{debug/data/sprint\_q5p\_l1\_vs\_l2\_diagnostic.json}). The
L2 $SL_2^3$ substrate of Remark~\ref{rem:q5p_decorated_pw} remains
valid as a complementary semisimple-only reading of the M1/M2/M3
partition;\ the Levi-decomposition
$U^*_{\mathrm{GeoVac}, \mathrm{Levi}} = \mathbb{G}_a^{3 N(n_{\max})} \times SL_2$
of Remark~\ref{rem:q5p_levi_synthesis_substrate} is the
structurally-correct Stage-2 target, matching the published
Connes--Marcolli motivic-Galois shape. The canonical multi-year next
step of the Q5' program is the Tannakian closure of the pro-unipotent
factor $\mathbb{G}_a^{3 N(n_{\max})}$ in the Connes--Marcolli machinery.
See Paper~55~\S~subsec:open\_m2\_m3.
\end{remark}

\begin{remark}[Q5' Stage-2 closed-form generating function for $T_{\mathrm{path}}^{\mathrm{abs}}(n_{\max})$ on the OffDiag substrate, with cutoff-independence + E1 selection + McKean-Singer-on-$A^2$ closures, Sprint Q5'-T-Path-Generating (T2), June 2026]
\label{rem:q5p_t_path_quartic}
The chirality-weighted two-step path count
\[
N^{(2)}_{s' \to s} \;:=\; \mathrm{Tr}(\gamma\, A\, e_s\, A\, e_{s'})/\kappa^2
\]
on the OffDiag substrate (Remark~\ref{rem:q5p_offdiag_dirac_enrichment})
is bit-exactly \emph{cutoff-independent}:\ every realised sector pair
$(s', s)$ has the same $N^{(2)}$ value at every $n_{\max}$ where both
sectors are present. Hypothesis verified bit-exactly across
$n_{\max} \in \{1, 2, 3, 4, 5\}$, including boundary pairs. The
E1 selection rule $|\Delta n|, |\Delta l| \le 1$ holds for all
realised transitions. The total signed sum vanishes identically:
\[
T_{\mathrm{path}}^{\mathrm{signed}}(n_{\max}) \;=\; \sum_{(s', s)} N^{(2)}_{s' \to s} \;=\; \mathrm{Tr}(\gamma A^2)/\kappa^2 \;\equiv\; 0,
\]
extending McKean-Singer ($\chi(S^3) = 0$, McKean-Singer on $A^0 = I$)
to the chirality-weighted $A^2$-trace. The total absolute path count
admits the bit-exact closed form
\[
\boxed{\;T_{\mathrm{path}}^{\mathrm{abs}}(n_{\max}) \;=\; \frac{(n_{\max}^2 + 9 n_{\max} - 4)(n_{\max}^2 + 13 n_{\max} - 6)}{6}\;}
\]
verified at panel $(8, 72, 224, 496, 924)$ for
$n_{\max} \in \{1, 2, 3, 4, 5\}$, predicting $T_{\mathrm{path}}^{\mathrm{abs}}(6) = 1548$.
Bit-exact verification (driver
\texttt{debug/compute\_q5p\_t\_path\_generating.py};\ data
\texttt{debug/data/sprint\_q5p\_t\_path\_generating.json};\ memo
\texttt{debug/sprint\_q5p\_t\_path\_generating\_memo.md}). Unlocks the
continuum Mellin lift of Remark~\ref{rem:q5p_mellin_lift_bridge}.
\end{remark}

\begin{remark}[Q5' Stage-2 continuum Mellin lift of the discrete bridge identity to the spectral zeta side with M2/M3 decomposition, Sprint Q5'-Mellin-Lift-Bridge (T1), June 2026]
\label{rem:q5p_mellin_lift_bridge}
Using the quartic closed form of
Remark~\ref{rem:q5p_t_path_quartic}, the continuum Mellin lift of
the v3.63.0 L3 discrete bridge identity $\mathrm{drift}_{n_{\max} \ge 3} = -\kappa^4$
(Remark~\ref{rem:q5p_bridge_identity}) is the generating function
\[
F(s) \;:=\; \sum_{n \ge 1} T_{\mathrm{path}}^{\mathrm{abs}}(n) \cdot n^{-s} \;=\; \tfrac{1}{6}\zeta(s-4) + \tfrac{11}{3}\zeta(s-3) + \tfrac{107}{6}\zeta(s-2) - \tfrac{53}{3}\zeta(s-1) + 4\,\zeta(s).
\]
$F(s)$ admits the bit-exact decomposition by shift parity
\[
F(s) \;=\; F_{\mathrm{M2}}(s) + F_{\mathrm{M3}}(s),
\]
\[
F_{\mathrm{M2}}(s) = \tfrac{1}{6}\zeta(s-4) + \tfrac{107}{6}\zeta(s-2) + 4\,\zeta(s) \quad (\text{even shifts},\ \text{M2 Seeley--DeWitt}),
\]
\[
F_{\mathrm{M3}}(s) = \tfrac{11}{3}\zeta(s-3) - \tfrac{53}{3}\zeta(s-1) \quad (\text{odd shifts},\ \text{M3 vertex-parity-Hurwitz}).
\]
Bit-exact verification:\ $F - F_{M2} - F_{M3} = 0$ symbolically;\ 5
simple poles at $s \in \{1, 2, 3, 4, 5\}$ with residues
$\{4, -53/3, 107/6, 11/3, 1/6\}$;\ integer-$s$ panel exhibits \emph{mixed
M2 ($\pi^{2k}$) + M3 ($\zeta(\text{odd})$) content} at every test point
(e.g.\ $F(6) = -53\zeta(5)/3 + \pi^2/36 + 4\pi^6/945 + 11\zeta(3)/3 + 107\pi^4/540$).
$F$ is structurally distinct from the v3.62.0 M3 $\eta_D(s)$ Hurwitz
Mellin (shares only the odd-shift structure with $\mathbb{Q}$
vs $2^s\,\mathbb{Q}$ coefficients). This is the \emph{first algebraic-side
realisation of the M1/M2/M3 master Mellin engine partition appearing
on a single Hopf-algebra-level invariant} — the Mellin engine extends
from the trace-evaluation level (case-exhaustion theorem
\ref{thm:pi_source_case_exhaustion}) to the substrate-side Mellin
generating function. Driver
\texttt{debug/compute\_q5p\_mellin\_lift\_bridge.py};\ data
\texttt{debug/data/sprint\_q5p\_mellin\_lift\_bridge.json};\ memo
\texttt{debug/sprint\_q5p\_mellin\_lift\_bridge\_memo.md}.
\end{remark}

\begin{remark}[Q5' Stage-2 $m_J$-resolved smash-product foundation:\ bit-exact $\Delta m_J \in \{-1, 0, +1\}$ grading and $U(1)$ z-rotation diagonalisation at $n_{\max} = 2$, Sprint Q5'-mJ-Smash-Product (T4), June 2026]
\label{rem:q5p_mJ_smash_foundation}
The $m_J$-resolved OffDiag substrate at $n_{\max} = 2$ admits a
bit-exact $\mathbb{Z}_3$-grading by
$\Delta m_J = m_{J,\mathrm{to}} - m_{J,\mathrm{from}} \in \{-1, 0, +1\}$
(Wigner-Eckart selection for the vector operator $A = \kappa A_{\mathrm{graph}}$)
with symmetric dimension count
\[
\dim \mathcal{A}_{-1} = 30, \quad \dim \mathcal{A}_0 = 40, \quad \dim \mathcal{A}_{+1} = 30, \quad \text{total} = 100,
\]
over all 100 non-zero off-diagonal Dirac entries.
The $U(1)$ z-rotation $U_\theta |n, \kappa, m_J\rangle = e^{i m_J \theta}|n, \kappa, m_J\rangle$
acts diagonally:\ $U_\theta T_{m_J' \to m_J} U_\theta^{-1} = e^{i \Delta m_J \theta} T_{m_J' \to m_J}$,
exactly reproducing the three phase factors
$\{1, e^{i\theta}, e^{-i\theta}\}$ that
Remark~\ref{rem:q5p_combined_substrate_levi} identified as broken
on the basic substrate — the $m_J$-resolved refinement
\emph{diagonalises} the $U(1)$ action. Critical scope finding:\
the $j$-content at $n_{\max} = 2$ includes $j = 3/2$ states (from
$\kappa = \pm 2$ in sectors $(2, 1)$ and $(2, 2)$), so the decorated-PW
substrate of Remark~\ref{rem:q5p_decorated_pw} at $j_{\max} = 1/2$ is
\emph{insufficient} for the full $SU(2)$ action on the $m_J$-resolved
OffDiag;\ \emph{L2 must be extended to $j_{\max} \ge 3/2$} as the
sprint-scale prerequisite (\textasciitilde 1 day) to the multi-year
smash-product construction
$\mathcal{A}^{m_J\text{-OD}} \rtimes \mathcal{O}(SL_2)$ — the only path
to a non-trivial Hopf extension beyond Levi-decomposition. Driver
\texttt{debug/compute\_q5p\_mj\_smash\_product.py};\ data
\texttt{debug/data/sprint\_q5p\_mj\_smash\_product.json};\ memo
\texttt{debug/sprint\_q5p\_mj\_smash\_product\_memo.md}.
\end{remark}

\begin{remark}[Q5' L2 substrate extended to $j_{\max} = 3/2$:\ sprint-scale prerequisite for the multi-year smash-product construction closed, Sprint Q5'-FO1 (v3.66.0), June 2026]
\label{rem:q5p_l2_extended_jmax32}
The sprint-scale prerequisite flagged by
Remark~\ref{rem:q5p_mJ_smash_foundation} (T4 of v3.65.0) for the
multi-year smash-product construction
$\mathcal{A}^{m_J\text{-OD}} \rtimes \mathcal{O}(SL_2)^{\otimes 3}$
at $n_{\max} = 2$ is closed bit-exactly.  The L2 decorated-PW substrate
of Remark~\ref{rem:q5p_decorated_pw} extends to $j_{\max} = 3/2$ with
all 7 Hopf axioms passing:\ dim $\mathcal{H}_{\mathrm{dec}}^{(3/2)} = 90$
(per slot:\ $\sum_{2j \le 3}(2j+1)^2 = 1 + 4 + 9 + 16 = 30$), 621 bit-exact
zero residuals across coassociativity ($90/90$), counit L+R ($180/180$),
antipode at SU(2)$^{\otimes 3}$ quotient ($90/90$ structural per slot),
$k$-grading preservation ($90/90$), and two pro-system Hopf-homomorphism
truncations $P_{3/2 \to 1}$ ($126/126$) and $P_{3/2 \to 1/2}$ ($45/45$).
The substrate factorisation $\mathcal{O}(SL_2)^{\otimes 3}$ and candidate
motivic Galois group $SL_2^3$ are unchanged from v3.63.0
$j_{\max} \in \{1/2, 1\}$ panel — the $j_{\max}$ parameter controls only
which irrep matrix coefficients are explicitly included;\ the underlying
group structure is $j_{\max}$-invariant.  This unblocks the multi-year
smash-product construction at the sprint-scale prerequisite level.
Driver \texttt{debug/compute\_q5p\_fo1\_l2\_extension\_jmax32.py};\
data \texttt{debug/data/sprint\_q5p\_fo1\_l2\_extension\_jmax32.json};\
memo \texttt{debug/sprint\_q5p\_fo1\_l2\_extension\_jmax32\_memo.md}.
\end{remark}

\begin{remark}[{Q5' Interpretation C closure:\ $F(s)$ integer-$s$ panel sits in MT period ring at level $\le 4$ over $\mathbb{Z}[i, 1/2]$, AND $U^*$ acts on $\chi, \eta, F(s)$ cocycle classes via period-pairing, Sprint Q5'-FO2+FO3 (v3.66.0), June 2026}]
\label{rem:q5p_interpretation_C_closure}
T5 of v3.65.0 identified three interpretations of "$U^*$ acts on
$\mathcal{H}_{\mathrm{Levi}}$":\ (A) Hopf-coaction, (B) Hopf-automorphism,
(C) period-pairing on the image
$\pi(\mathcal{H}_{\mathrm{Levi}}) \subset \mathbb{C}$ via motivic Galois
conjugation. Interpretation C is the cleanest sprint-scale closure.
This remark records its bit-exact closure for all named cocycle-class
families. T5 SQ1+SQ2+SQ5 combined:\ (i) every $\chi_{(n, l)}$ value
and every $\eta_{(n, l)}$ value (60 + 60 entries up to $n_{\max} = 4$)
is an integer in $\mathbb{Z}$ at depth 0 in the MT period ring;\
(ii) every term of $F(s)$ (Remark~\ref{rem:q5p_mellin_lift_bridge}) at
integer $s \in \{6, 7, 8, 9, 10\}$ (5 cells $\times$ 5 terms = 25 term
classifications) is either pure-Tate $\pi^{2k}\,\mathbb{Q}$ (M2 component
of v3.65.0 T1) or odd-Riemann $\zeta(\mathrm{odd})\,\mathbb{Q}$ (M3 component);
both sit in $\mathrm{MT}(\mathbb{Q}, 1)$, which embeds in
$\mathrm{MT}(\mathbb{Z}[i, 1/2], 4)$ trivially. The Paper 18 \S III.7
M1/M2/M3 master Mellin engine partition aligns bit-exactly with the
standard MT depth/weight grading at integer-shifted $\zeta$ values:\
$k = 0$ M1 (Hopf-base measure, depth 0, weight 2 in $\pi$), $k = 2$ M2
(Seeley--DeWitt, depth 0, weight $2k'$ in $\pi^{2k'}$), $k = 1$ M3
(vertex-parity-Hurwitz, depth 1, weight $2k'+1$ in $\zeta(2k'+1)$).
The cosmic-Galois $U^*$ acts on $\chi, \eta$ trivially (depth 0
fixed-points), on the M2 component of $F(s)$ as the Tate subgroup
(invariant up to $\pi^{2k}$ scaling), and on the M3 component as the
standard motivic Galois action on $\mathrm{MT}(\mathbb{Q}, 1)$
odd-zeta classes (Brown 2012, Glanois 2015). Total bit-exact zero
residuals:\ 145 (25 + 60 + 60).  Driver
\texttt{debug/compute\_q5p\_fo2\_fo3\_mt\_period\_containment.py};\
data \texttt{debug/data/sprint\_q5p\_fo2\_fo3\_mt\_period.json};\
memo \texttt{debug/sprint\_q5p\_fo2\_fo3\_mt\_period\_memo.md}.
Honest scope:\ Interpretation C closes for the named cocycle classes
($\chi, \eta, F(s)$ panel);\ Interpretations A (Hopf-coaction) and B
(Hopf-automorphism) remain open as named sprint-scale (B via SQ3) and
multi-year (A) follow-ons. The L1 follow-on (a) Tannakian closure of
the pro-unipotent factor in the Connes--Marcolli machinery remains
the canonical multi-year next step.
\end{remark}

\begin{remark}[Rational-residue extension of the case-exhaustion theorem:
the Bernoulli ladder and the functional-equation reading of the two-layer
split (Sprint GB, 2026-05-29)]
\label{rem:bernoulli_ladder}
The case-exhaustion theorem~\ref{thm:pi_source_case_exhaustion} classifies
the $\pi$-content of every GeoVac observable as
$\mathcal{M}[\mathrm{Tr}(D^k\, e^{-tD^2})]$. The companion question is
where the framework's \emph{rational} coefficients come from. A
diagnostic on the gravity sector's small rationals places them on a single
Bernoulli ladder, with each rung $B_{2n}$ glued across two windows by the
Riemann functional equation:
$$
\zeta(2n) = (-1)^{n+1}\frac{B_{2n}(2\pi)^{2n}}{2(2n)!}
   \in \pi^{2n}\,\mathbb{Q}
\qquad\Longleftrightarrow\qquad
\zeta(1-2n) = -\frac{B_{2n}}{2n} \in \mathbb{Q}.
$$
The left (positive-even) window is the M2 transcendental ring of
Layer~2 observables; the right (negative-integer) window is the
$\pi$-free rational residue carried by Layer~1 skeleton coefficients.
Verified placements (all exact in \texttt{sympy}):
\begin{itemize}
\item Conical-defect tip coefficient $\tfrac{1}{12} = -\zeta(-1) =
B_2/2$ (\S4.15); per-$t$ UV target $\tfrac{1}{24} = -\zeta(-1)/2 =
B_2/4$ (the string $c/24$); scalar $S^3$ Casimir $\tfrac{1}{240} =
+\zeta(-3)/2$ ($B_4$ rung); the $\alpha$-conjecture's
$F = \pi^2/6 = \pi^2 B_2 = \zeta(2)$ (Paper~2; the M2 window of the
same $B_2$ whose skeleton window is the conical tip).
\item \emph{Non-trivial internal check}: the replica-entropy
derivative is forced, not free ---
$\frac{d}{d\alpha}\big[\tfrac{1}{12}(\tfrac{1}{\alpha}-\alpha)\big]
\big|_{\alpha=1} = -\tfrac{1}{6} = -B_2 = 2\zeta(-1)$, the
$\alpha$-derivative of the tip's $B_2/2$.
\item \emph{Controls behave}: $\kappa = -1/16$ (the Fock Jacobian
$\Omega^{-4}$) has no clean $\zeta$ form; the Dirac Casimir
$17/480$ misses the scalar $B_4$ rung by exactly $1/32$ and requires
the half-integer Hurwitz shift $\zeta(-3,\tfrac{1}{2})$, i.e.\ the
scalar/spinor distinction is the integer/half-integer (M2/M3) split.
\end{itemize}
The reading is the rational-residue extension of
Theorem~\ref{thm:pi_source_case_exhaustion}: the rational coefficients
are the negative-integer analytic continuation of the \emph{same}
Mellin object whose positive-even values are the Layer-2
transcendentals, and the empirical two-layer split (Paper~34/35) is the
Riemann functional-equation reflection $\zeta(s)\leftrightarrow
\zeta(1-s)$ instantiated per Bernoulli order.

\emph{Honest scope.} This is a structural correspondence at the same
epistemic level as Theorem~\ref{thm:pi_source_case_exhaustion}, not a
forcing proof; it concerns $\zeta$ at \emph{integer} arguments
(special values), not the non-trivial zeros, so it makes no contact
with the critical-line structure (cf.\ WH6 and the three classical-RH
walls RH-N/O/M, which remain in force). The $B_2$, $B_4$, \emph{and}
$B_6$ rungs are now populated (Sprint GB-3): the $B_6$ rung's
falsifiable target was met by the conformal-scalar Casimir on $S^5$
(a GeoVac manifold via the Bargmann--Segal lattice, Paper~24), which
evaluates to $-31/60480 = \tfrac{1}{24}(\zeta(-5) - \zeta(-3))$ --- the
$B_6$ rung $\zeta(-5) = -1/252$ enters with clean coefficient
$+\tfrac{1}{24}$, alongside $B_4$ at $-\tfrac{1}{24}$, with $B_2$
\emph{absent}. The $\zeta(-5) - \zeta(-3)$ spreading (rather than a
pure $\zeta(-5)$) is the \emph{same} quadratic-multiplicity mechanism
that obstructs the functional equation in Paper~28's
functional-equation obstruction remark (the $S^5$ harmonic degeneracy is quartic),
now seen on the Casimir side from an independent observable; the method
reproduces the $S^3$ control $1/240 = \tfrac{1}{2}\zeta(-3)$ (Paper~35).
The integer/half-integer (M2/M3) split recurs: the conformal scalar
gives the clean ladder value, the half-integer HO/Bargmann $S^5$
Casimir gives $-17/3840$ (the ``$17$'' echoing the Dirac-$S^3$ $17/480$).
Next rung: $S^7$ conformal Casimir is predicted to spread across
$\zeta(-7), \zeta(-5), \zeta(-3)$. The ladder is thus supported across
three rungs and two observable types, but remains a structural
correspondence, not a forcing theorem.
Source: \texttt{debug/sprint\_gb\_bernoulli\_ladder\_memo.md},
\texttt{debug/sprint\_gb3\_b6\_rung\_memo.md}.
\end{remark}

\begin{remark}[Bisognano--Wichmann reading of the M1 sub-mechanism on
temporal compactifications (Track D, May 2026)]
\label{rem:bisognano_wichmann_reading}
The M1 sub-mechanism of the master Mellin engine,
$\mathcal{M}[\mathrm{Tr}(D^0\, e^{-tD^2})]$, produces a
$\mathrm{Vol}/\mathrm{Mellin\text{-}measure}$ factor on the spatial
$S^3$ Hopf base ($\mathrm{Vol}(S^2)/\pi^2 = 4/\pi$, the L2 propinquity
rate constant; Sprint MR-A/B in Remark~\ref{rem:master_mellin_domain})
and on a temporal $S^1_\beta$ factor (the bare circumference $\beta$
itself; Sprint TD Track 1, $T_{S^3}\otimes T_{S^1_\beta}$).  Sprint TD
Track 4 (\texttt{debug/sprint\_td\_track4\_memo.md}, May~2026) extends
the temporal-side instantiation to the Euclidean Schwarzschild cigar,
where smoothness at $r=2M$ forces the $\tau$-circle period
$\beta_{\mathrm{cigar}} = 8\pi M$, reproducing the Hawking temperature
$T_H = 1/(8\pi M)$ from M1 alone (sympy residual zero).  Track~D
(\texttt{debug/bisognano\_wichmann\_track\_d\_memo.md}, May~2026) makes
explicit the dual physical reading of M1 on temporal compactifications:
under the standard Wick-rotation chain
(Hartle--Hawking~\cite{hartle_hawking1976};
 Sewell~\cite{sewell1982};
 Bisognano--Wichmann~\cite{bisognano_wichmann1976}),
the framework's M1 $2\pi$ on $S^1_\tau$ IS the period of the
modular-flow / Lorentz-boost orbit of a stationary observer outside
the Schwarzschild horizon.  Equivalently, the M1 sub-mechanism on
temporal compactifications is the Wick-rotated image of the
Bisognano--Wichmann modular-flow mechanism.  This is a
\emph{structural correspondence}, not a literal identification: two
mathematical objects (a Hopf-bundle measure factor on the Riemannian
side; a modular-automorphism period on the Lorentzian side) are
equated by the published Wick-rotation prescription
(Bisognano--Wichmann~1976; Sewell~1982; Hartle--Hawking~1976), not by
an operator-level proof inside the truncated spectral triple
$\mathcal{T}_{n_{\max}}$.  The natural follow-up that would lift the
correspondence to literal identification --- compute the modular
Hamiltonian $K$ on $\mathcal{T}_{n_{\max}}$ for a wedge-restricted
Camporesi--Higuchi state and verify $\sigma_{i\cdot 2\pi} = $ identity
at finite $n_{\max}$, then take the GH limit using the lemmas of
Theorem~\ref{thm:gh_convergence} --- is named as the principal falsifier
in Track~D \S 3.1; estimated cost 4--8 weeks; priority MEDIUM-HIGH.
The case-exhaustion theorem~\ref{thm:pi_source_case_exhaustion} itself
is unchanged: M1 now has \emph{two physical interpretations}
(Hopf-bundle measure / boost-orbit period via Bisognano--Wichmann),
unified by Wick rotation, not two independent constraints on the
master Mellin engine partition.
The corollary on Rindler wedges (Unruh temperature
$T_U = a/(2\pi)$ from the same M1 mechanism with
$\beta_{\mathrm{Rindler}} = 2\pi/a$) is a 1--2 day pendant follow-up
flagged in Track~D \S 5.4; not pursued in the 1-week sprint window.

\textbf{Unruh pendant (May 2026).}\ The Track~D~\S 5.4 follow-up was
executed as a 1-day pendant
sprint~(\texttt{debug/unruh\_pendant\_memo.md}, May 2026).  Wick-rotating
Rindler coordinates $(\eta, \rho)$ for an observer at proper acceleration
$a$ produces Euclidean polar $(\eta_E, \rho)$ with conical regularity at
$\rho = 0$ forcing $\eta_E$-period $2\pi$, hence proper-time period
$\beta_U = 2\pi / a$ and Unruh $T_U = a/(2\pi)$
(Unruh~\cite{unruh1976}).  Sprint~TD Track~1's
\verb|matsubara_spectrum()| applied verbatim with $\beta = 2\pi/a$
reproduces the Unruh-thermal Matsubara spectrum bit-identically (sympy
residual zero on $T_U = 1/\beta_U$, lowest bosonic mode $= a$, lowest
fermionic mode $= a/2 = \pi T_U$).  The single $2\pi$ in $\beta_U$ is
the same M1 Hopf-base measure / $\mathrm{Vol}(S^1)$ signature that drove
the cigar identification, instantiated on the Rindler $\eta_E$-circle
instead of the cigar $\tau$-circle.  The reading is identical:\
under the Wick-rotation chain, the framework's M1 $2\pi$ on
$S^1_{\eta_E}$ IS the period of the modular-flow / boost-orbit of the
accelerated observer.  Hawking, Bisognano--Wichmann, and Unruh are
therefore three faces of one Wick-rotation theorem with one M1
sub-mechanism, all of which the master Mellin engine tracks at $k = 0$;
the published Wick-rotation prescription supplies the bridge to the
Lorentzian-side modular-flow vocabulary in each case.  The
\emph{operator-level falsifier} (modular Hamiltonian on
$\mathcal{T}_{n_{\max}}$ for a wedge-restricted state has analytic
period $2\pi$) named in Track~D~\S 3.1 covers all three faces
simultaneously, since the bridge runs through the same KMS-Tomita-Takesaki
algebra in every case; lifting the period-closure check on any one face
lifts the structural correspondence on Hawking, Bisognano--Wichmann, and
Unruh together.  Same scope (structural correspondence, not literal
identification); same falsifier; same MEDIUM-HIGH 4--8 week priority.

\textbf{Sprint L1 closure (2026-05-16):\ literal identification at the
operator-system level (Riemannian).}\ The principal falsifier named
above is now closed positively at the strongest possible level.\ Sprint
L1 (\texttt{debug/l1\_modular\_hamiltonian\_results\_memo.md}) constructs
the modular Hamiltonian $K$ on $\mathcal{T}_{n_{\max}}$ for the
hemispheric-wedge-restricted (W1 of L1-A) Camporesi--Higuchi state and
verifies the period-closure
$\sigma_{2\pi}(O) = O$ for all four physical witnesses
(Bisognano--Wichmann, Hartle--Hawking, Sewell, Unruh) at every tested
$n_{\max} \in \{2, 3, 4, 5\}$.  Maximum periodicity residual scales as
$O(\sqrt{\dim \mathcal{H}_{n_{\max}}} \cdot \varepsilon_{\mathrm{machine}})$:
$1.8 \times 10^{-16}$ at $n_{\max}=2$ ($\dim_H = 16$),
$4.8 \times 10^{-16}$ at $n_{\max}=3$ ($\dim_H = 40 = g_3^{\mathrm{Dirac}}
= \Delta^{-1}$),
$9.3 \times 10^{-16}$ at $n_{\max}=4$ ($\dim_H = 80$),
$1.6 \times 10^{-15}$ at $n_{\max}=5$ ($\dim_H = 140$).  Cross-witness
collapse is bit-exact (residuals match across witnesses at machine
precision, max consistency residual $= 0.0$ at every tested $n_{\max}$):
the period closure is independent of $\kappa_g$ because the
geometric BW-$\alpha$ realization $K = J_{\mathrm{polar}}$ (the rotation
generator preserving the wedge) has integer spectrum (odd integers
$\text{two\_m}_j$), so $\mathrm{e}^{\mathrm{i}\cdot 2\pi \cdot n} = 1$
holds for all integer $n$.  The closure is finite-cutoff exact, not a
limit-statement, and does not require the Gromov--Hausdorff machinery
of Theorem~\ref{thm:gh_convergence} (although it is consistent with it).
This lifts the BW reading from ``structural correspondence'' to
``literal identification at the operator-system level (Riemannian)''
at every finite $n_{\max}$.  The four-witness Wick-rotation theorem
(HH + Sew + BW + Unruh) becomes a framework-internal theorem at
signature $(3, 0)$.  The Lorentzian extension to signature $(3, 1)$
via the Bizi--Brouder--Besnard~2018~\cite{bizi_brouder_besnard2018}
Krein-space spectral triple is the named Sprint L2 follow-up
(6--12 month scale per Sprint L0 audit
\texttt{debug/lorentzian\_l0\_audit\_memo.md}).
The L1-A architecture obstructions
\S\S10a--10e of \texttt{debug/l1\_modular\_hamiltonian\_architecture\_memo.md}
were addressed: \S10a operator-system leakage is zero by construction
($\sigma_{2\pi}(O) = O$ exactly, so $\sigma_{2\pi}(O)$ trivially stays
in $\mathcal{O}_{n_{\max}}$); \S10b/c chirality / M3 trivialization
consistency: the closure is purely M1-content (the $2\pi$ in the
period is the Hopf-base measure
$\mathrm{Vol}(S^1)$, identified per Sprint TS-E1's master Mellin engine
case-exhaustion theorem~\ref{thm:pi_source_case_exhaustion}); \S10d
Gibbs phenomenon: the hemispheric projector $P_W$ uses a discrete
$m_j \to -m_j$ reflection (no smooth cutoff needed, exact involution);
\S10e $J_{\mathrm{TT}}$-vs-$J_{GV}$ compatibility: the Tomita conjugation
is constructed as a fresh class
\texttt{geovac.modular\_hamiltonian.TomitaConjugation} with verified
$J_{\mathrm{TT}}^2 = +I$, categorically distinct from the Connes
KO-dim 3 charge conjugation $J_{GV}$ of
\texttt{geovac.real\_structure} (which has $J_{GV}^2 = -I$).

\textbf{Sprint L1-tighten (2026-05-16 evening):\ unified BW-$\alpha$/BW-$\gamma$
closure.}\ The L1 closure above used the BW-$\alpha$ geometric K =
$J_{\mathrm{polar}}$ as the load-bearing object;\ the Tomita-Takesaki
BW-$\gamma$ realization was only an expectation-level KMS cross-check.
Sprint L1-tighten
(\texttt{debug/l1\_tighten\_tomita\_results\_memo.md}) rebuilds the
BW-$\gamma$ construction at the full Tomita-Takesaki level on the GNS
Hilbert-Schmidt space $H_{\mathrm{GNS}} = M_{\dim_W}(\mathbb{C})$,
extracts $K_{\mathrm{TT}} = -\log \Delta$ from the polar decomposition
$S = J_{\mathrm{TT}} \Delta^{1/2}$, and verifies
$\sigma_{2\pi}^{\mathrm{TT}}(O) = O$ at the operator-action level for
all four physical witnesses at every tested $n_{\max} \in \{2, 3, 4,
5\}$.  Result:\ \emph{UNIFIED\_STRONG closure} at all tested cutoffs
and witnesses.  Maximum periodicity residual at the Tomita-Takesaki
level: $1.09\times 10^{-16}$ ($n_{\max}=2$), $1.07\times 10^{-15}$
($n_{\max}=3$), $3.03\times 10^{-15}$ ($n_{\max}=4$), $4.79\times
10^{-15}$ ($n_{\max}=5$).  The BW-$\alpha$ and BW-$\gamma$ flows are
operator-action-level conjugates: $\sigma_t^{\mathrm{TT}}(a) =
\sigma_{-t}^{\alpha}(a)$ bit-exactly at general $t$ (verified at $t =
1$ with residual $<10^{-16}$), and identical at the period $t = 2\pi$
because $e^{\pm i 2\pi n} = 1$ for any integer $n$.  The
state-dependent $J_{\mathrm{TT}}(X) = \rho^{1/2} X^* \rho^{-1/2}$ from
the polar decomposition satisfies $J_{\mathrm{TT}}^2 = +I$ to machine
precision at every $n_{\max}$ (the L1 verification only checked the
canonical $U = I$ sanity case; L1-tighten verifies the $\rho^{1/2}$-weighted
genuine Tomita conjugation).  The construction uses the natural local
Hamiltonian $H_{\mathrm{local}} := K_\alpha^W / \beta$, making the
wedge KMS state $\rho_W = e^{-K_\alpha^W}/Z$ the operator-system
analog of the Wightman BW vacuum $\rho_W \propto e^{-2\pi
K_{\mathrm{boost}}}$ on the Rindler wedge.  The four-witness
Wick-rotation theorem is closed at signature $(3, 0)$ via
\emph{both} the geometric BW-$\alpha$ and the Tomita-Takesaki
BW-$\gamma$ readings, not just one;\ the L1 ``BW-$\alpha$-only''
caveat is removed.
\end{remark}

\begin{theorem}[Gromov--Hausdorff convergence on $S^3$]
\label{thm:gh_convergence}
Let $\mathcal{T}_{S^3} = (C^\infty(S^3), L^2(S^3, \Sigma), D_{\mathrm{CH}})$
be the round-$S^3$ metric spectral triple with the Camporesi--Higuchi
Dirac, and let $\mathcal{T}_{n_{\max}} = (\mathcal{O}_{n_{\max}},
\mathcal{H}_{n_{\max}}, D_{n_{\max}})$ be the Connes--van Suijlekom
truncated metric spectral triple at cutoff $n_{\max}$ (operator system
from sprint~R2.1; truthful Camporesi--Higuchi Dirac restricted to the
truncation).  Then the truncated triples converge to $\mathcal{T}_{S^3}$
in the Latr{\'e}moli{\`e}re quantum Gromov--Hausdorff propinquity:
$$
   \Lambda(\mathcal{T}_{n_{\max}}, \mathcal{T}_{S^3})
   \;\le\; C_3 \cdot \gamma_{n_{\max}}
   \;\xrightarrow[n_{\max}\to\infty]{}\; 0,
$$
where $C_3 = 1$ is the Lipschitz comparison constant and
$\gamma_{n_{\max}}$ is the L2 mass-concentration moment of the central
spectral Fej{\'e}r kernel on SU(2).  The bound is qualitative-rate;
the asymptotic rate $\gamma_{n_{\max}} = O(\log n / n)$ is consistent
with but not rigorously proved by the small-$n$ closed forms (deferred
to a quantitative-rate sub-track).
\end{theorem}

\begin{proof}[Proof sketch]
The proof is a Latr{\'e}moli{\`e}re tunneling-pair assembly of five
lemmas closed in WH1 sprint~R2.5 (April--May~2026):

\begin{itemize}
\item[\textbf{L1${}'$}] (Sprint R3.5, 2026-05-04;
\texttt{debug/wh1\_r35\_full\_dirac\_memo.md}): the offdiag
Camporesi--Higuchi operator system substrate has every non-forced
cross-shell Connes-distance pair finite; this provides the
Connes-vS-compatible metric sub-structure on which the propinquity is
defined.

\item[\textbf{L2}] (Sprint R2.5/L2, 2026-05-04;
\texttt{debug/r25\_l2\_proof\_memo.md}): the positive central spectral
Fej{\'e}r kernel $K_{n_{\max}} = Z_{n_{\max}}^{-1} |\sum_{j \le j_{\max}}
\sqrt{2j+1}\,\chi_j|^2$ on SU(2) is normalized, central, with
mass-concentration moment $\gamma_{n_{\max}} \to 0$ (closed-form
algebraic values $\gamma_2, \gamma_3, \gamma_4$, monotone decrease),
Plancherel symbol $\hat{K}_{n_{\max}}(N) = N/Z_{n_{\max}}$ on the
Fock-shell index, and central-multiplier cb-norm
$\|T_K\|_{\mathrm{cb}} = 2/(n_{\max}+1)$ on the central subalgebra
(Bo{\.z}ejko--Fendler / Pisier 2001 Ch.~8 transcription).

\item[\textbf{L3}] (Sprint R2.5/L3, 2026-05-04;
\texttt{debug/r25\_l3\_proof\_memo.md}): the Lipschitz comparison
$\|[D_{\mathrm{CH}}, M_f]\|_{\mathrm{op}} \le C_3 \|\nabla
f\|_{L^\infty}$ holds with $C_3 = 1$ on the natural Avery--Wen--Avery
test panel at $n_{\max} \in \{2, 3, 4\}$, with theoretical bound
$(N-1)/\sqrt{N^2-1} \nearrow 1$ asymptotically.  The
``torsion-and-curvature correction'' anticipated in the Track-A
master memo is identified concretely as the spectral-gap quantization
$|n - n'|$ of the truthful Camporesi--Higuchi Dirac.

\item[\textbf{L4}] (Sprint R2.5/L4, 2026-05-06;
\texttt{debug/r25\_l4\_proof\_memo.md}): the Berezin-type reconstruction
map $B_{n_{\max}}: C(S^3) \to \mathcal{O}_{n_{\max}}$ defined by
$B_{n_{\max}}(f) = \sum_{N \le n_{\max}}\sum_{L,M} \hat{K}_{n_{\max}}(N)\,
c_{NLM}\, M_{NLM}$ (using L2 Plancherel weights and Avery--Wen--Avery
multipliers from sprint~R3.1) is positive (preserves $f \ge 0$),
contractive ($\|B_{n_{\max}}(f)\|_{\mathrm{op}} \le \|f\|_{C(S^3)}$),
an approximate identity ($\|B_{n_{\max}}(f) - P_{n_{\max}} M_f
P_{n_{\max}}\|_{\mathrm{op}} \le \gamma_{n_{\max}} \|\nabla
f\|_{L^\infty}$), and L3-compatible
($\|[D_{\mathrm{CH}}, B_{n_{\max}}(f)]\|_{\mathrm{op}} \le \|\nabla
f\|_{L^\infty}$).  The map is the SU(2) analog of the
Connes--van~Suijlekom 2021~\S3 Fej{\'e}r-kernel reconstruction on
$S^1$, equivalently the non-K{\"a}hler analog of the Hawkins~2000
Berezin map for K{\"a}hler coadjoint orbits.

\item[\textbf{L5}] (Sprint R2.5/L5, 2026-05-06;
\texttt{debug/r25\_l5\_proof\_memo.md}): the tunneling pair
$(B_{n_{\max}}, P_{n_{\max}})$ between $\mathcal{T}_{S^3}$ and
$\mathcal{T}_{n_{\max}}$ has reach $\le \gamma_{n_{\max}}$ in both
directions (L4(c) approximate identity + the dual estimate via L2(g)
Bo{\.z}ejko--Fendler), height contributions bounded constants
(L4(b) contractivity + $\mathrm{height}_P = 0$ as $P$ is a projection),
giving
$\Lambda(\mathcal{T}_{n_{\max}}, \mathcal{T}_{S^3}) \le C_3 \cdot
\gamma_{n_{\max}} \to 0$.
\end{itemize}

The full lemma chain is verified numerically at $n_{\max} \in \{2, 3, 4\}$
with monotone decrease of the bound (2.075, 1.610, 1.322); ratio
$\Lambda(n_{\max}=4) / \Lambda(n_{\max}=2) = 0.637$.  Module
\texttt{geovac/gh\_convergence.py} packages the assembly with 39 unit
tests passing, integrating with the L1${}'$--L4 modules verbatim.
\end{proof}

\begin{remark}[Limit identification]
\label{rem:gh_limit_id}
A companion proposition (Latr{\'e}moli{\`e}re-propinquity-completion +
Kantorovich--Rubinstein duality) identifies the propinquity limit
of $(\mathcal{T}_{n_{\max}}, d_{D_{n_{\max}}})$ with the
Wasserstein--Kantorovich state space $(P(S^3), d_{\mathrm{Wass}})$
against round-$S^3$ geodesic distance.  The argument is bookkeeping
at the L4(a) positivity + L4(b) contractivity level (Rieffel
1999/2004; D'Andrea--Lizzi--Martinetti 2014) and is recorded in
\texttt{debug/r25\_l5\_proof\_memo.md~\S7}.  No new analytical content.
\end{remark}

\begin{remark}[WH1 keystone closure]
\label{rem:wh1_proven}
Theorem~\ref{thm:gh_convergence} closes the WH1 keystone (CLAUDE.md~\S1.7).
The structural claim of WH1 --- that GeoVac is an almost-commutative
spectral triple with Connes--van Suijlekom truncations converging to
the round-$S^3$ spectral triple --- is now \emph{proved} at the
qualitative-rate level.  The ``torsion-and-curvature correction''
anticipated in the Track-A master memo is identified concretely as the
spectral-gap quantization $|n - n'|$ of the truthful Camporesi--Higuchi
Dirac, with per-multiplier ratio bounded by $N - 1$ from the SO(4)
selection rule, and the corresponding asymptotic ratio
$(N-1)/\sqrt{N^2-1} \nearrow 1$ giving the L3 constant $C_3 = 1$. The
remaining open questions (Higgs-gap analysis on the four-way $S^3$
coincidence; full real structure $J$ audit at finite $n_{\max}$, scoping
memo \texttt{debug/real\_structure\_finite\_nmax\_memo.md}) are
downstream of this theorem.
\end{remark}

\begin{remark}[Pythagorean M1 closure of the L2-E residual]
\label{rem:pythagorean_m1_closure}
The case-exhaustion theorem~\ref{thm:pi_source_case_exhaustion}
classifies the transcendental ring in which the $\pi$-content of a
finite chain of Paper~34 projections must sit.  A
post-Sprint-L2-E PSLQ probe (\texttt{debug/h\_local\_residual\_pslq\_memo.md},
2026-05-17) provides a concrete operator-level instantiation of the
theorem at $k = 0$ (the M1 sub-mechanism / Hopf-base measure sector).
On the truncated Lorentzian Krein spectral triple
(Paper~43~\cite{paper43}~\S 10.2), the
Bisognano--Wichmann-aligned local Hamiltonian
$H_{\mathrm{local}} = K^{\alpha, W}_L / \beta$ and the
wedge-restricted Lorentzian Dirac $D_W^L$ are Hilbert--Schmidt
orthogonal,
\[
   \langle H_{\mathrm{local}}, D_W^L \rangle_{\mathrm{HS}}
   \;=\; 0
   \qquad \text{(bit-exact at every tested $(n_{\max}, N_t)$),}
\]
and the squared residual decomposes as a Pythagorean sum in the M1
ring:
\[
   \big\| H_{\mathrm{local}} - D_W^L \big\|_F^2
   \;=\; \frac{\kappa_g^2 \cdot S(n_{\max})}{4\pi^2}
       \;+\; D(n_{\max}),
\]
with $S, D \in \mathbb{Q}[n_{\max}]$ given by cumulative wedge
Casimir traces (Paper~43~\S 10.2,
Cor.~cor:pythagorean\_orthogonality therein).  The sole
transcendental in the residual is $1/\pi^2 = 1/\mathrm{Vol}(S^1)^2$,
the squared Hopf-base-measure signature, exactly the M1 output the
case-exhaustion theorem predicts for a $k = 0$ chain that engages no
M2 heat-kernel content and no M3 vertex-parity Hurwitz content.

This sharpens the engine's predictive content along a complementary
axis to Remark~\ref{rem:master_mellin_domain}.  Sprint MR-A/B/C
established the mechanism-as-domain partition (M1 owns
propinquity-rate observables, M2 owns heat-kernel residuals, M3 owns
vertex-restricted parity observables).  The L2-E result adds a
fourth empirical instantiation:\ M1 owns the modular-Hamiltonian /
spectral-action generator-orthogonality residual on the
hemispheric-wedge KMS state, with the specific ring decomposition
$\mathbb{Q}[n_{\max}] \oplus \pi^{-2}\,\mathbb{Q}[n_{\max}]$ that
the case-exhaustion theorem allows at $k = 0$.  The engine does not
just classify which $\pi$-powers can appear; it now correctly
predicts the exact M1 ring for the modular-Hamiltonian generator
distinction observed in Paper~42~\S 7.2 and Paper~43~\S 10.1.

The orthogonality is universal across all six modular witnesses
(Paper~42~\cite{paper42}~Cor.~cor:single\_construction;\
Paper~43~\S 10.2), via $\kappa_g$-linearity of $H_{\mathrm{local}}$
and $\kappa_g$-independence of $D_W^L$.  The structural mechanism
is the diagonal/off-diagonal subspace decomposition of
$\mathcal{B}(\mathcal{H}_W)$:\ $H_{\mathrm{local}}$ is real and
diagonal in the unfolded wedge spinor basis (eigenvalues
$\mathrm{two\_m}_j$), $D_W^L$ is off-diagonal in the same basis
(intertwines $\pm m_j$ chirality partners).  A formal proof of the
subspace decomposition at general $(n_{\max}, N_t)$ is a named
follow-on;\ the empirical verification across the 18-cell
six-witness panel is strong evidence the decomposition is
structural rather than coincidental.

\emph{Cross-references.}\ Paper~43~\S 10.2
Cor.~cor:pythagorean\_orthogonality (closed form and
witness-universality);\ Paper~42~\S 8
Rem.~rem:pythagorean\_underlies\_collapse (the
six-witness-collapse statement at the inner-product level);\
Paper~24~\S 5.3 four-layer Coulomb/HO asymmetry (the structural
finding that the Pythagorean orthogonality is $S^3$-specific to
the construction --- the $S^5$ Hardy sector has no clean analog,
joining the modular-Hamiltonian level as a fourth Coulomb/HO
asymmetry layer alongside spectrum-computing $L_0$, calibration
$\pi$, and Wilson matter coupling).
\end{remark}

\begin{remark}[BH Phase~0 — clean negative on wedge KMS area-law]
\label{rem:bh_phase0_negative}
The modular-Hamiltonian construction closed by Sprint~L1 (the
operator-system-level four-witness Wick-rotation theorem, this
section) and by Sprint~L2-E (the Lorentzian extension at finite
cutoff) reproduces the Hartle--Hawking and Unruh \emph{temperatures}
via the M1 mechanism (\S~\ref{rem:bisognano_wichmann_reading}).  A
natural follow-up question is whether the same construction also
reproduces the Bekenstein--Hawking \emph{entropy} $S_{BH} = A/(4G)$ at
finite cutoff via the von Neumann entropy of the wedge KMS state
$\rho_W = e^{-K_\alpha^W}/Z$.  Sprint~BH-Phase0
(\texttt{debug/bh\_phase0\_diagnostic\_memo.md}, 2026-05-22) tests
this directly at $n_{\max} \in \{2, 3, 4, 5, 6, 7\}$.

\emph{Result.}\ $S(\rho_W) \approx 1.944\,\log(n_{\max}) + 0.565$ at
$R^2 = 0.99991$ across the panel.  No polynomial scaling appears in
any of the four candidate fits (log-linear, power-law, poly-2,
poly-3).  An $s$-deformation $\rho(s) = e^{-s K_\alpha^W}/Z$ across
$s \in \{0.01, 0.05, 0.1, 0.5, 1, 2, 2\pi\}$ shows power-law exponents
$\alpha \in [0.64, 0.72]$, never close to the $\alpha = 2$ that a
naive Bekenstein--Hawking area-law $S \sim n_{\max}^2$ would predict.

\emph{Mechanism.}\ The structural reading is
\[
   S(\rho_W) \;\approx\; \log\bigl(n_{\mathrm{equator}}\bigr),
   \qquad n_{\mathrm{equator}} = n_{\max}\,(n_{\max} + 1),
\]
where $n_{\mathrm{equator}}$ is the count of wedge states with smallest
$|m_j| = 1/2$ (the equator analog).  The BW canonical state is
effectively the maximally mixed state on the lowest-$K_\alpha$ shell,
with exponentially suppressed contributions $\sim e^{-(2k+1)}$ from
shells $k = 1, 2, \dots$.  The two asymptotic limits saturate
bit-exactly:\ $S(s \to 0) = \log(\dim W)$ (hot, maximally mixed) and
$S(s \to \infty) = \log(n_{\mathrm{equator}})$ (cold, lowest shell
only).  The BW canonical $s = 1$ sits between, closer to the cold
limit.

The slope $\approx 2$ equals $\dim(\partial W) = \dim S^2_{\mathrm{equator}}$
(the dimension of the wedge boundary manifold), giving a
Cardy--Calabrese-style log entropy along the $d = 2$ boundary.  This
is structurally weaker than the QFT Bekenstein--Hawking area-law
$S \sim A \cdot \Lambda^2 / 4$, which would scale as a power of the
UV cutoff times the boundary dimension.

\emph{Structural reading.}\ The operator-system truncation that closes
the four-witness Wick-rotation theorem at machine precision
(\S~VIII.F, Sprint~L1) does \emph{not} preserve the continuum QFT
vacuum entanglement that would give the Bekenstein--Hawking area-law.
The modular-flow period closure (M1 Hopf-base measure mechanism) and
the Bekenstein--Hawking area-law are structurally distinct features
of the gravitational thermodynamics chain;\ the GeoVac framework
captures the former at finite cutoff and not the latter.  Consistent
with the structural-skeleton-scope pattern (CLAUDE.md~\S 2, sprints
H1, LS-8a, HF-3/4/5, multi-focal-composition wall, W3):\ the
framework determines the algebraic skeleton of modular flow but does
not autonomously generate the calibration data (here:\ Newton's
constant) that links wedge-entropy to physical area.

Cross-references:\ Paper~34~\S V.B Bekenstein--Hawking probe row;\
Sprint TD Track~4 (Hawking temperature, machinery-witness, M1
closure);\ the May-16 BH scoping memo
(\texttt{debug/bekenstein\_hawking\_sprint\_scoping\_memo.md},
$8$--$14$ week Connes--Chamseddine cigar route) remains the longer-route
alternative if a literal $S_{BH} = A/(4G)$ closure is pursued, but
with the honest scope statement that Newton's constant is calibration
input not derivation.
\end{remark}

% -------------------------------------------------------------------
\subsection{The unified $U(1) \times SU(2) \times SU(3)$ gauge
content and its limits}
\label{sec:unified_gauge}
% -------------------------------------------------------------------

The case-exhaustion theorem of Sec.~\ref{sec:pi_source_theorem}
identifies M1 (Hopf-base measure) as the mechanism by which the
Wilson plaquette-Haar projection injects $\pi$ into spectral-triple
observables.  This subsection records a complementary structural
fact about the family of GeoVac sub-triples that engage M1: they
realise all three Standard-Model gauge groups as Wilson lattice
constructions, on three structurally distinct sub-manifolds, with a
single unified weak-coupling coefficient $1/(4\,N_c)$.  The
construction is honest about what it is and is not: a unified
\emph{catalogue} of three Wilson-Yang--Mills instances of the
Marcolli--van~Suijlekom gauge-network framework, not a unified
spectral triple in the Connes--Chamseddine
sense~\cite{chamseddine_connes2010}.

\subsubsection*{Three Wilson constructions, one kinetic coefficient}

Three independent sprints have produced gauge-invariant, finite,
Haar-normalisable non-abelian Wilson lattice gauge theories on
GeoVac sub-graphs:

\begin{itemize}\setlength{\itemsep}{2pt}
\item \emph{Paper~25}~\cite{paper25} constructs the abelian
$U(1)$ Wilson--Hodge structure on the Fock-projected $S^3$ Coulomb
graph (Paper~7), with phases $\chi_v$ on nodes and $A_e$ on edges
covariant under the standard incidence operator.
\item \emph{Paper~30}~\cite{paper30} constructs the non-abelian
$\mathrm{SU}(2)$ Wilson lattice gauge theory on the same
$S^3 = \mathrm{SU}(2)$ graph, with link variables $U_e \in
\mathrm{SU}(2)$, plaquettes from primitive non-backtracking walks,
and Haar measure $dU$.  Paper~30~\S 5.2 verifies the weak-coupling
kinetic coefficient symbolically and numerically.
\item \emph{Sprint~ST-SU3} (May~2026, modules
\texttt{geovac/su3\_wilson\_s5.py},
\texttt{tests/test\_su3\_wilson\_s5.py}, $39/39$ pass) constructs
the non-abelian $\mathrm{SU}(3)$ Wilson lattice gauge theory on the
Bargmann--Segal $S^5$ graph of Paper~24~\cite{paper24}, with link
variables $U_e \in \mathrm{SU}(3)$ in the fundamental representation
(independent of the shell index $N$), plaquettes from primitive
non-backtracking walks, and SU(3) Haar measure.  Sprint~ST-SU3~\S 4
verifies Theorem~2 (kinetic coefficient $1/12$) symbolically and
numerically at $N_{\max} = 2, 3$ ($7\,765$ plaquettes), with maximal
Cartan reduction to a coupled $U(1) \times U(1)$ theory verified at
machine precision.
\end{itemize}

The three constructions are summarised in
Table~\ref{tab:unified_gauge}.  The columns are not independent
discoveries:\ each is a specialisation of the gauge-networks
framework of Marcolli--van~Suijlekom~\cite{marcolli_vs2014}, with
the Perez-Sanchez correction~\cite{perez_sanchez2025} that the
continuum limit of such gauge networks is Wilson Yang--Mills
\emph{without} a Higgs sector.

\begin{table}[h]
\centering
\small
\begin{tabular}{|l|c|l|l|c|}
\hline
\textbf{Construction} & $G$ & \textbf{Host manifold} &
\textbf{Vertex Hilbert space} $\mathcal{H}_v$ &
$\frac{1}{4 N_c}$ \\
\hline
Paper~25 & $U(1)$ & $S^3$ Hopf base & $\mathbb{C}$ (scalar) &
$\frac{1}{4}$ \\
Paper~30 & $\mathrm{SU}(2)$ & $S^3 = \mathrm{SU}(2)$ &
$\mathbb{C}^2$ (Dirac) & $\frac{1}{8}$ \\
ST-SU3 & $\mathrm{SU}(3)$ & $S^5$ Bargmann--Segal &
$\bigoplus_N \mathbb{C}^{\dim(N,0)}$ & $\frac{1}{12}$ \\
\hline
\end{tabular}
\caption{The three Wilson lattice gauge constructions as instances of
the Marcolli--van~Suijlekom gauge-network framework on GeoVac
sub-manifolds.  The kinetic coefficient
$\tfrac{1}{4 N_c}$ in the rightmost column is the weak-coupling
expansion coefficient of the Wilson action $S_W = \beta \sum_P \bigl(
1 - \tfrac{1}{N_c} \mathrm{Re}\,\mathrm{tr}\,U_P \bigr)$ around the
trivial vacuum, in the standard normalisation
$\mathrm{tr}(T^a T^b) = \tfrac{1}{2} \delta^{ab}$.  The three rows
verify this coefficient independently:\ Paper~30~\S 5.2 and
Sprint~ST-SU3~\S 4 to machine precision; the abelian Paper~25
case is the trivial $N_c = 1$ limit of either non-abelian
construction.}
\label{tab:unified_gauge}
\end{table}

In each construction the edge Laplacian $L_1 = B^{\top} B$, where
$B$ is the signed node--edge incidence matrix, plays the role of
weak-coupling kinetic operator on $\mathfrak{g}$-valued $1$-cochains.
Paper~25's ``discrete Hodge--1 Laplacian'' is therefore a single
combinatorial object across all three constructions; only its target
Lie algebra changes ($\mathfrak{u}(1)$, $\mathfrak{su}(2)$,
$\mathfrak{su}(3)$).  This is the single mathematical identity
unifying the three sectors at the kinetic level.

\subsubsection*{The two-step SU(3) story}

The $\mathrm{SU}(3)$ row of Table~\ref{tab:unified_gauge} requires a
specific clarification, because an earlier Sprint~5 Track~S5
(April~2026, memo \texttt{debug/s5\_gauge\_structure\_memo.md})
returned a clean negative when asking whether the $S^5$
Bargmann--Segal lattice carries a non-abelian gauge structure
analogous to Paper~25's $U(1)$.  The two findings are not in tension:\
they answer two different questions, and read together they give a
two-step honest picture.

\begin{itemize}\setlength{\itemsep}{2pt}
\item \emph{Step 1 (Sprint~5 Track~S5, negative).}  An
adjacency-preserving $\mathrm{SU}(3)$ action on the $(N, 0)$ tower of
the Bargmann graph fails Wilson's fixed-group-on-every-link
requirement, because the transitions $(N, 0) \to (N+1, 0)$ are
Clebsch--Gordan intertwiners between distinct $\mathrm{SU}(3)$
irreducible representations of growing dimension $\dim(N, 0) =
(N+1)(N+2)/2$, not group elements of a single $\mathrm{SU}(3)$.  The
natural shell-respecting gauge group of the Bargmann graph is
$U(1)$, and this is recorded in Paper~25~\S VII.A as a structural
asymmetry between the $S^3$ and $S^5$ sectors.
\item \emph{Step 2 (Sprint~ST-SU3, positive at the gauge level).}  A
\emph{shell-independent} Wilson lattice gauge theory with link
variables $U_e \in \mathrm{SU}(3)$ in the fundamental representation
(not in any of the $(N, 0)$ irreps) is a structurally different
object that bypasses the Step~1 obstruction.  The links live as
external gauge degrees of freedom on edges; they do not act on
shell labels.  Sprint~ST-SU3 verified gauge-invariance of the
Wilson action under node-local
$U_e \mapsto g_{\mathrm{src}(e)} U_e g_{\mathrm{tgt}(e)}^{\dagger}$
to machine precision and computed the kinetic coefficient $1/12$
both symbolically and numerically.
\item \emph{Matter coupling (Sprint~ST-SU3 Theorem~3, mixed).}  When
one tries to couple this Wilson $\mathrm{SU}(3)$ to the natural
shell matter on the $(N, 0)$ tower, the original Step~1 obstruction
is rediscovered:\ a fundamental-representation link cannot uniformly
couple to representations of growing dimension, and the matter
coupling reverts to the Clebsch--Gordan intertwiner structure.
Wilson $\mathrm{SU}(3)$ on the Bargmann graph is therefore
self-consistent as a \emph{pure-gauge} theory but does not admit a
natural matter coupling in the sense of Paper~30's spin-$1/2$
Dirac fermions on the Coulomb graph.
\end{itemize}

The two-step picture is consistent with Paper~24's Coulomb/HO
asymmetry and sharpens it from three layers to four:\ to the prior
asymmetries (spectrum-computing role of $L_0$; calibration~$\pi$;
Wilson gauge-with-natural-matter) one now adds a fourth layer at
which the $S^3$ Coulomb side is special:\ the gauge group acts in
the same representation as the matter, and matter--gauge coupling
is structurally available.  On the $S^5$ Bargmann side the
combinatorial Wilson--Hodge vocabulary is universal but the
gauge--matter coupling is not.

\subsubsection*{Marcolli--vS lineage of the unified content}

Each construction in Table~\ref{tab:unified_gauge} is a specific
gauge-networks instance in the sense of Marcolli and
van~Suijlekom~\cite{marcolli_vs2014}, with the Perez-Sanchez
correction~\cite{perez_sanchez2025} that the continuum limit of
such a gauge network is Wilson Yang--Mills \emph{without} a Higgs.  The vertex algebra is
$\mathcal{A}_v = \mathbb{C}$ (or $C(V)$ globally) in all three rows;
the vertex Hilbert space and the gauge group differ row-by-row, as
listed in the table.  In each row the gauge-network spectral action
reduces to the standard Wilson action plus a quadratic kinetic term
in the weak-coupling expansion, with coefficient $1/(4 N_c)$.

The lineage placement is identical to that of
Sec.~\ref{sec:lineage}, and each Wilson construction inherits the
same status:\ each is a published-machinery instance with a
specific manifold and gauge group; what is not duplicated in the
literature is the \emph{simultaneous} existence of all three SM
gauge groups as siblings across GeoVac sub-manifolds.

\subsubsection*{What this is not:\ structural distance to a Standard Model}

Table~\ref{tab:unified_gauge} is a catalogue, not a unified theory,
and the rhetoric rule of CLAUDE.md~\S 1.5 forbids treating it as the
latter.  The structural distance from this catalogue to Connes' SM
construction lives in four named gaps.

\begin{enumerate}\setlength{\itemsep}{2pt}
\item[\textbf{G1}] (Three generations of matter.)\ Each
construction in Table~\ref{tab:unified_gauge} carries a single
matter species:\ Paper~30 has one generation of Dirac fermions
(Paper~14 Tier~2~\cite{paper14}); Sprint~ST-SU3 has one tower of
$(N, 0)$ irreps.  Sprint~4H Track~SM-B (April~2026) tested the naive
shell-to-generation map $\sum_f N_c Q_f^2 = 8 = |\lambda_3|$ and
returned a clean negative:\ the per-generation contribution is
charge-universal across the three SM generations whereas every
per-shell $S^3$ invariant varies with $n$, so no shell$\to$generation
map is consistent with charge-universal per-generation structure.
The $8 = 8$ equality was a numerical coincidence with no
representation-theoretic content.  There is no GeoVac mechanism
currently selecting a generation tripling of the Hilbert space.
\item[\textbf{G2}] (Higgs as inner fluctuation of $D$.)\ In
Connes' SM, the Higgs field arises as the off-diagonal block of an
inner fluctuation $D \mapsto D + \omega + \epsilon' J \omega
J^{-1}$ ($\omega = \sum_i a_i [D, b_i]$) for $a_i, b_i \in
\mathcal{A}_F = \mathbb{C} \oplus \mathbb{H} \oplus M_3(\mathbb{C})$
in an almost-commutative algebra.  Per Perez-Sanchez 2024, the
Marcolli--vS gauge-network spectral action gives YM \emph{without}
a Higgs sector by default:\ gauge networks supply only edge-Dirac
hopping data with no internal fiber Dirac, so the off-diagonal
$\mathbb{C} \leftrightarrow \mathbb{H}$ Yukawa block that turns
inner fluctuations into a Higgs is absent.  The natural extension
for GeoVac is $\mathcal{A}_{\mathrm{GV}} \otimes \mathcal{A}_F$
with the electroweak slice $\mathcal{A}_F = \mathbb{C} \oplus
\mathbb{H}$ (deferring color), and with a finite Dirac
$D_F$ on the four-dimensional electroweak fiber whose off-diagonal
block encodes the Yukawa coupling matrix.  Inner fluctuations on
this extension formally split into a gauge-1-form sector (which
must reproduce the unified $U(1) \times \mathrm{SU}(2)$ content of
Papers~25--30) and a candidate Higgs scalar sector living in the
off-diagonal $(\mathbb{C}, \mathbb{H})$ component.  Whether the
required $D_F$ can be selected from GeoVac structure rather than
imposed is the load-bearing open question; we report the result of
this test (Sprint~H1) below in Sec.~\ref{sec:higgs_h1}.  Note that
the operator-system non-multiplicative-closure of
Remark~\ref{rem:operator_system} (witness pair $M_{2,1,0}^2$,
$14.9\%$ residual against $\mathcal{O}_{n_{\max}=2}$) is a separate,
finite-resolution feature of the truncated operator system itself;
inner fluctuations are defined operationally as $\omega + J\omega
J^{-1}$ with $\omega = \sum_i a_i [D, b_i]$, which uses ALGEBRA
elements (the multiplier matrices, which are an algebra over
$\mathbb{C}$ at the level of $M_{\dim H}(\mathbb{C})$ even though
their span $\mathcal{O}_{n_{\max}}$ is not closed under multiplication),
so the $a \cdot b \notin \mathcal{O}$ feature does not block the
fluctuation calculus.
Sprint~H1 (Sec.~\ref{sec:higgs_h1}) closes this gap up to the
Yukawa-undetermined verdict; Sprint~G3 (Sec.~\ref{sec:g3}) supplies
the structural reason that $Y$ is undetermined --- $\gamma_{\mathrm{GV}}$
and $\gamma_F$ are independent commuting $\mathbb{Z}_2$ gradings, the
Yukawa lives on the $\gamma_F$ flip, and no GeoVac-side data couples
to that flip.  G2 and G3 are therefore one structural fact, not two.
\item[\textbf{G3}] (Hypercharge / electroweak unification.)\ The SM
relation $Q = T_3 + Y_W/2$ requires co-located $U(1)_Y$ and
$\mathrm{SU}(2)_L$ on a common Hilbert space with a left--right
chiral asymmetry.  Paper~25's $U(1)$ and Paper~30's
$\mathrm{SU}(2)$ are co-located on the same $S^3$ graph (this is
the most concrete near-term unification target, see below), but
the Camporesi--Higuchi Dirac (Sec.~\ref{sec:construction}
Definition~\ref{def:D_GV_spectral}) is non-chiral by default.  The
chirality grading introduced in Sprint~TS Track~E2 (R3.5, full
Dirac with both chiralities) is a $\mathbb{Z}_2$-grading of the
spinor bundle, not the SM weak-isospin chirality, and it does not
distinguish left- and right-handed sectors of a charged matter
field.
Sprint~G3 (Sec.~\ref{sec:g3}) closed this in the negative on $S^3$
alone:\ $\gamma_{\mathrm{GV}}$ and $\gamma_F$ are independent commuting
$\mathbb{Z}_2$ gradings on $\mathcal{H}_{\mathrm{GV}} \otimes
\mathcal{H}_F$, with operator-norm residual exactly~$2$ at every
$n_{\max} \in \{1, 2, 3\}$.  Identification is structurally
impossible by tensor-product factorisation, and the electroweak
co-location route is therefore pushed to G4.
\item[\textbf{G4}] (Cross-manifold obstruction, the dominant gap.)\
Connes' SM uses a \emph{single} spectral triple
$(\mathcal{A}_{\mathrm{SM}}, \mathcal{H}_{\mathrm{SM}},
D_{\mathrm{SM}})$ on a single base manifold, with all three SM
gauge groups emerging from inner automorphisms of one
almost-commutative algebra.  The GeoVac construction has, at
present, three separate spectral triples on two different graphs
(the $S^3$ Coulomb graph for $U(1)$ and $\mathrm{SU}(2)$; the $S^5$
Bargmann graph for $\mathrm{SU}(3)$).  There is no GeoVac operator
that geometrically unifies the $S^3$ and $S^5$ graphs into a
single spectral triple containing both as commuting factors:\ at
best the Bargmann and Coulomb graphs can be encoded as orthogonal
factors of a tensor-product Hilbert space, but the resulting
object is a direct sum of two Wilson-YM instances rather than a
unification.  This is the most consequential gap.  It is also the
most directly tied to the four-layer Coulomb/HO asymmetry of
Sec.~\ref{sec:coulomb_ho}:\ the absence of cross-manifold
unification is a structural symptom of the same asymmetry that
forbids a calibration-$\pi$ analog and a spectrum-computing $L_0$
on $S^5$.
Sprint~G4 scoping (Sec.~\ref{sec:g4}) sharpens this gap into G4a
(Connes SM on $\mathcal{T}_{S^3}$ alone with
$\mathcal{A}_F = \mathbb{C} \oplus \mathbb{H} \oplus M_3(\mathbb{C})$,
1--2 month sprint, predicted positive-thin) and G4b (genuine
cross-manifold $\mathcal{T}_{S^3} \otimes \mathcal{T}_{S^5}$, blocked
at the NCG-framework level by the Riemannian-vs-Hardy-sector
asymmetry).  After G3 NEGATIVE, G4a is the only sprint-scale
electroweak-unification path remaining.
\end{enumerate}

A footnote on the GUT-embedding direction:\ Sprint~4H Track~SM-C
(April~2026) tested whether $1/40 = \Delta$ in Paper~2's $K = \pi(B
+ F - \Delta)$ formula appears naturally in $\mathrm{SU}(5)$,
$E_8$, or $\mathrm{Spin}(10)$ as a Chern number, $\eta$-invariant,
or Chern--Simons level, and returned a clean negative:\ the dual-%
Coxeter rewrite $1/40 = (3/8)/15 = \sin^2\theta_W^{\mathrm{GUT}} /
\dim(\mathbf{\bar{5}} \oplus \mathbf{10})$ is tautological,
recycling the integers $5$ and $8$ already present in Paper~2.  The
naive shell-to-generation embeddings $\mathrm{shell}_n \to
\mathbf{\bar{5}}$ fail at the $\mathrm{SU}(3) \times \mathrm{SU}(2)$
branching level.  GeoVac therefore has no current contact with the
GUT-embedding lineage of the SM, beyond the tautological numerical
coincidences.

\subsubsection*{The most concrete near-term unification target}

Of the four gaps, G4 (cross-manifold) is the most consequential and
G3 (electroweak unification) is the most reachable.  Paper~25's
$U(1)$ and Paper~30's $\mathrm{SU}(2)$ already live on the same
underlying graph (the Fock-projected $S^3$ Coulomb graph) and
share the same Camporesi--Higuchi Dirac operator, the same
incidence structure $B$, and the same edge Laplacian $L_1 = B^{\top}
B$.  The two constructions differ only at the gauge-group level,
and Paper~30~\S 5.1 (Result~1) establishes that Paper~25 is
exactly the maximal-torus reduction of Paper~30.  The natural
unified construction is therefore an $\mathrm{SU}(2) \times U(1)$
Wilson lattice gauge theory on the same graph, with link variables
$U_e \in \mathrm{SU}(2) \times U(1)$ and a chosen embedding of
hypercharge as the $U(1)$ phase.  Whether such a construction
admits a Higgs-style inner fluctuation through an
almost-commutative extension (closing G2 within the electroweak
sector) is an open structural question; whether the
Camporesi--Higuchi Dirac admits a chirality grading compatible with
weak-isospin (closing G3) is the second open structural question.
This is a clean $S^3$-only target that does not require any
progress on the $S^5$ side and that would, if successful, take
GeoVac from ``three Wilson-YM instances on three sub-manifolds''
to ``unified electroweak Wilson construction on $S^3$, plus a
separate $\mathrm{SU}(3)$ instance on $S^5$.''

We do not claim such a unification has been achieved, or even that
its component pieces have been verified.  The status of this
subsection is:\ a recorded structural fact (the unified $1/(4
N_c)$ coefficient and the lineage placement are real); a recorded
honest gap analysis (G1--G4 are real); and a recorded open
direction ($S^3$ electroweak co-location is the most concrete
target).  Per CLAUDE.md~\S 1.5, no claim is made that GeoVac
contains the Standard Model; the claim is only that it contains
the three SM gauge groups as separate Wilson lattice gauge
theories on its own sub-manifolds, with a unified weak-coupling
kinetic coefficient and a unified gauge-network lineage.

\subsection{Sprint H1:\ Higgs as inner fluctuation on the GeoVac
electroweak almost-commutative extension}
\label{sec:higgs_h1}

The G2 gap of Sec.~\ref{sec:unified_gauge} (Higgs as inner
fluctuation) admits a constructive sprint-scale test:\ build the
minimal electroweak almost-commutative extension
$\mathcal{T}_{\mathrm{AC}} = \mathcal{T}_{\mathrm{GV}} \otimes
\mathcal{T}_F$ with $\mathcal{A}_F = \mathbb{C} \oplus \mathbb{H}$ on
the truthful Camporesi--Higuchi GeoVac triple, compute the inner
fluctuation $D \mapsto D + \omega + \epsilon' J \omega J^{-1}$
explicitly at finite $n_{\max}$, and ask whether the off-diagonal
$\mathbb{C} \leftrightarrow \mathbb{H}$ Yukawa block of the finite
Dirac $D_F$ that turns Marcolli--vS-without-Higgs into a Higgs
construction can be \emph{selected from GeoVac structure} rather
than imposed.  This is the H1 sprint reported in
\texttt{debug/h1\_ac\_extension\_memo.md}; we summarize its content
here.

\paragraph{Construction.}  Following Connes--Marcolli 2008
Ch.~13~\cite{connes_marcolli2008}, the finite electroweak Hilbert
space is doubled into matter and antimatter sectors:\ $\mathcal{H}_F
= \mathcal{H}_F^{\mathrm{matter}} \oplus
\mathcal{H}_F^{\mathrm{antimatter}} = \mathbb{C}^4 \oplus \mathbb{C}^4$,
with $\mathcal{A}_F = \mathbb{C} \oplus \mathbb{H}$ acting on the
matter sector only (the opposite-algebra action on antimatter is
realized via $J$).  The 8 by 8 finite Dirac is $D_F =
\mathrm{block\_diag}(M, \bar{M})$ where the matter-sector $M$ is the
4 by 4 Yukawa-Dirac with off-diagonal $L \leftrightarrow R$ block
$Y = \mathrm{diag}(y_\nu, y_e)$.  The combined Dirac is $D =
D_{\mathrm{GV}} \otimes \mathbb{1}_F + \gamma_{\mathrm{GV}} \otimes
D_F$ with $\gamma_{\mathrm{GV}} = \mathrm{sign}(D_{\mathrm{GV}})$
(Track~2 closure):\ this commutes with $D_{\mathrm{GV}}$ on the
truthful CH Dirac.  The combined real structure is $J = J_{\mathrm{GV}}
\otimes J_F$ with $J_F$ swapping matter and antimatter via complex
conjugation; matter-antimatter doubling makes $J_F^2 = +\mathbb{1}$
(KO-dim 6 of $\mathcal{A}_F$), which combined with KO-dim 3 of the
GV factor gives KO-dim $3 + 6 = 9 \equiv 1 \pmod 8$ for the
combined triple, hence $J^2 = -\mathbb{1}$ and $J D = +D J$.  Both
are verified to machine precision at $n_{\max} \le 3$
(\texttt{tests/test\_almost\_commutative.py}, 38 tests passing).

\paragraph{Inner fluctuations split cleanly into gauge and Higgs
sectors.}  For $\omega = \sum_i a_i [D, b_i]$ with $a_i, b_i \in
\mathcal{A}_{\mathrm{GV}} \otimes \mathcal{A}_F$, we have

\begin{equation}
\label{eq:omega_decomp}
\omega = \underbrace{\sum_i a_i^{\mathrm{GV}} [D_{\mathrm{GV}},
b_i^{\mathrm{GV}}] \otimes a_i^F b_i^F}_{\text{gauge piece}} +
\underbrace{\sum_i a_i^{\mathrm{GV}} b_i^{\mathrm{GV}}
\gamma_{\mathrm{GV}} \otimes a_i^F [D_F, b_i^F]}_{\text{Higgs
piece}}.
\end{equation}

The gauge piece reproduces Papers~25 ($U(1)$ from the
$\mathbb{C}$-summand) and~30 ($\mathrm{SU}(2)$ from the
$\mathbb{H}$-summand) at the operator level.  The Higgs piece is
non-zero exactly when $[D_F, b_i^F]$ has off-diagonal $\mathbb{C}
\leftrightarrow \mathbb{H}$ matrix elements --- which happens
\emph{iff} the imposed Yukawa $Y$ is non-zero.  At
$n_{\max} = 1, 2, 3$, with three test choices of $Y$:

\begin{center}
\begin{tabular}{lccc}
\hline
$D_F$ choice & $\|\Phi\|_{\max}$ & $\|\omega_{\mathrm{gauge}}\|_{\max}$ &
\textbf{Falsifier} \\
\hline
$Y = 0$           &  $0.000$ &  $0.000$--$0.580$ & \textbf{HOLDS}  \\
$y_e = 0.1$, $y_\nu = 0$ & $0.046$--$0.080$ & $0.000$--$0.580$ & FAILS \\
$y_e = 0.3$, $y_\nu = 0.2$ & $0.171$--$0.271$ & $0.000$--$0.580$ & FAILS \\
\hline
\end{tabular}
\end{center}

(50 random generators per cell; data
\texttt{debug/data/h1\_falsifier.json}).

\paragraph{Verdict:\ positive-thin.}  The Higgs construction is
\emph{well-defined} on the GeoVac AC extension at finite $n_{\max}$:
the algebra, Hilbert space, Dirac, and real structure satisfy the
relevant Connes axioms, and inner fluctuations decompose into a
gauge sector recovering Papers~25/30 plus a Higgs sector $\Phi$ that
is non-trivial whenever the imposed Yukawa $Y$ is non-zero.

\emph{However}, GeoVac structure does \textbf{not autonomously
select} a non-trivial $Y$.  Two candidate sources for $D_F$
off-diagonal data were flagged in the H1 scoping memo
(\texttt{debug/almost\_commutative\_scoping\_memo.md}~\S 3.3):\
candidate~A used the offdiag CH chirality bridging variant of
the GeoVac Dirac, but Track~2 (\texttt{debug/real\_structure\_finite\_nmax\_memo.md}~\S 3.D) ruled this out by showing $J D = +D J$ fails at residual~$\sim 2.0$ on offdiag CH; candidate~B was a Sturmian / DUCC inter-shell coupling whose physical interpretation as a Yukawa is itself speculative and was not pursued.

The strong falsifier of the scoping memo~\S 5
(``every Hermitian $D_F$ derivable from GeoVac structure forces
the Higgs sector to vanish'') therefore does not fire as a clean
negative.  The construction admits a Higgs given any $Y$; the
question is which $Y$, and GeoVac data is silent.  This is
exactly the positive-thin reading of the scoping memo~\S 0:\
GeoVac sits on the \emph{Marcolli--vS-without-Higgs} side of the
2014/2024 distinction at the structural level --- not because
inner fluctuations cannot define a Higgs (they can), but because
no GeoVac mechanism selects the Yukawa.

\paragraph{What this means for G2.}  Gap G2 of
Sec.~\ref{sec:unified_gauge} is sharpened to the more precise
statement:\ \emph{the Higgs construction goes through, but the
Yukawa is a free input}, not closed to either ``GeoVac contains a
Higgs'' (which would require a structural selection of $Y$) or
``GeoVac excludes a Higgs'' (which would require a structural
obstruction to non-zero $Y$).  This places G2 in the same status
as G1 (three generations are formally definable but not selected
by GeoVac structure):\ definable but not derivable.

\paragraph{What this means for G3.}  Gap G3 (hypercharge /
electroweak unification, requiring identification of the GV
chirality grading $\gamma_{\mathrm{GV}}$ with the SM weak-isospin
chirality) is \emph{not} closed by H1.  In our construction
$\gamma_{\mathrm{GV}}$ is the sign of $D_{\mathrm{GV}}$ on the
spinor bundle, while the SM weak-isospin chirality is a separate
$\mathbb{Z}_2$-grading on the matter sector of $\mathcal{H}_F$.
The two are added independently; no GeoVac mechanism currently
identifies them.

\paragraph{What this means for the operator-system non-multiplicative-closure
caveat.}  Remark~\ref{rem:operator_system}'s observation that
$\mathcal{O}_{n_{\max}}$ is not closed under multiplication
(witness pair $M_{2,1,0}^2$, $14.9\%$ residual at $n_{\max} = 2$)
is a separate, finite-resolution feature of the operator system
itself.  Inner fluctuations are defined operationally as $\omega +
J \omega J^{-1}$ with $\omega = \sum_i a_i [D, b_i]$, where the
multiplier matrices $a_i, b_i$ are individual elements of
$M_{\dim H}(\mathbb{C})$ (which is an algebra over $\mathbb{C}$);
the lack of closure of \emph{their span} $\mathcal{O}_{n_{\max}}$
under multiplication does not block the fluctuation calculus.
Order-zero $[a, J b J^{-1}] = 0$ holds automatically by the
matter--antimatter doubling (a acts on matter, $J b J^{-1}$ on
antimatter, sectors are disjoint).

\paragraph{Scope and limitations.}  H1 establishes the
constructive existence of the AC extension at finite $n_{\max}$
on the truthful CH GeoVac triple.  H1 does \emph{not} establish a
Higgs vev as a derived quantity (that would require a continuum
spectral-action computation; deferred to Track~1 R2.5~L4
closure for the rigorous interpretation), does \emph{not} include
the SM color factor $M_3(\mathbb{C})$ (deferred per
Sec.~\ref{sec:unified_gauge} G4), and does \emph{not} identify the
GV and weak-isospin chiralities (G3 stays open).  Per CLAUDE.md
\S 1.5, no claim is made that GeoVac contains the Higgs:\ the
claim is that it admits a Higgs given an imposed Yukawa, and that
the Yukawa is not derived from GeoVac structure.

\paragraph{Files.}  Implementation
\texttt{geovac/almost\_commutative.py};\ tests
\texttt{tests/test\_almost\_commutative.py} (38 passing);\ data
\texttt{debug/data/h1\_falsifier.json};\ memo
\texttt{debug/h1\_ac\_extension\_memo.md}.

\paragraph{Thermal extension of the H1 verdict (Sprint TD Track~3,
May 2026).}
The same Yukawa-undetermined structural reading extends to the
thermal Higgs effective potential $V_{\mathrm{eff}}(\phi, T)$ on the
combined triple $T_{\mathrm{GV}} \otimes T_{S^1_\beta} \otimes T_F$.
At one-loop, the thermal mass squared takes the form
\[
  m_T^2(T) \;=\; \frac{T^2}{12} \cdot \mathrm{Tr}_{\mathrm{inner}}\bigl(Y_{\mathrm{diag}}^2\bigr)
                \;+\;\text{(gauge contribution)},
\]
where the outer Stefan--Boltzmann backbone $T^2 / 12$ comes from
Sprint TD Track~1's $M_1 \times M_2$ factorisation
(Paper~35~\S\ref{sec:tensor_product_verification}), and the inner
$\mathrm{Tr}_{\mathrm{inner}}(Y_{\mathrm{diag}}^2)$ is a $k=2$
Dirichlet sum in the Yukawa eigenvalues
(Paper~18~\S~IV.6 and Theorem~\ref{thm:eta_trivialization}-style
filtering applies).  Any candidate $T_F$ whose squared Yukawa
diagonal reproduces the empirical SM spectrum reproduces the
SM-data-fit value $m_T^2(v, T_C) \approx (160~\mathrm{GeV})^2$;
GeoVac admits any compatible spectrum but selects none.

\textbf{Empirical falsification of the strict $Y = 0$ reading.}
The strict Marcolli--vS-without-Higgs case ($Y = 0$) predicts
$\mathrm{Tr}_{\mathrm{inner}}(Y_{\mathrm{diag}}^2) = 0$ and hence
no high-$T$ symmetry-restoring quadratic correction to
$V_{\mathrm{eff}}(\phi, T)$.  This is empirically falsified by the
electroweak phase transition / sphaleron crossover at
$T_C \approx 160$~GeV inferred from baryogenesis phenomenology.  The
empirical universe forces $Y_F > 0$.  This does not derive $Y_F$,
but it closes the H1 verdict's residual ambiguity:\ at finite $T$,
the framework plus observed sphaleron phenomenology rules out the
strict $Y = 0$ slice.  The ``thin'' part of the positive-thin
verdict is now thinner --- five thermal observables (Sprint TD
Track~3 catalogue) are individual probes of distinct inner-factor
structures (Yukawas, color algebra, generation count, charge
assignment, Higgs sector), and the inner factor is required to
deliver a specific value at each.  Three further candidates are
\emph{forbidden} by Theorem~\ref{thm:eta_trivialization} on any
Krajewski-class even spectral triple --- a positive structural
prediction:\ axial chemical potential thermodynamics, parity-odd
transport, and finite-$T$ EDMs vanish identically on the inner
factor of any Connes--Chamseddine-compatible AC extension.

\paragraph{Modular propinquity reformulation of the H1 falsifier
(Sub-sprint M-H1, May 2026).}
The dual modular propinquity construction~\cite{latremoliere2021dual}
provides a natural test of whether the Higgs sub-bimodule structure
imposes a constraint on the Yukawa $Y$ beyond what the Connes axioms
already enforce. The bimodule reformulation of the H1 verdict
(Sub-sprint M-H1, 2026-05-23) closes \emph{negatively} on this
question: the dual modular duality requirement reduces to Hermiticity
of $D_F$ plus matter--antimatter doubling, both already imposed by
Connes' axiom system; the Higgs cross-block $\Phi$ lives entirely in
the fiber sub-bimodule $\mathcal{A}_F = \mathbb{C} \oplus \mathbb{H}$,
structurally decoupled from any GeoVac-side bimodule data; and the
Morita-triviality constraint (``could $\Phi = u^* v$ for
GeoVac-derived $u, v$?'') does not survive matter--antimatter
doubling. The H1 POSITIVE-THIN verdict survives unchanged. The
``dual'' in dual modular propinquity is the
bridge $\leftrightarrow$ tunnel construction
duality~\cite{latremoliere2021dual}, NOT Morita equivalence (a
common misreading); the Morita route was structurally barred
independently. \emph{Net:} the framework gains substantive new
content on the chemistry / bound-state QED side (Sub-sprints M-X,
M-Y, M-Z) but not on the SM-unification side (M-H1). Source:
\texttt{debug/sprint\_modular\_propinquity\_mH1\_higgs\_memo.md},
\texttt{debug/sprint\_modular\_propinquity\_synthesis\_memo.md}.

\paragraph{Sprint M-Z partition tested on a second concrete case
(W1c chemistry wall, Sprint $\beta$-neg 2026-05-23).}
The Sub-sprint M-Z cross-shift / endomorphism partition was tested on
a second concrete case beyond the bound-state QED context where it was
originally derived: the W1c-residual chemistry wall at NaH alkali
hydrides (Paper~19, Sec.~\texttt{sec:w1c\_residual}). The bridge sprint
(\texttt{debug/sprint\_w1c\_mz\_partition\_analysis\_memo.md}, 2026-05-23)
proposed a \emph{three-bucket refinement} --- cross-shift (handled) /
basis-closable cross-shift (handled at sufficient basis size) /
endomorphism (external input) --- and tested it via three prespecified
falsifiable predictions for NaH max\_n=3
(\texttt{debug/data/sprint\_w1c\_mz\_partition\_predictions.json}). The
F1 max\_n=3 sprint
(\texttt{debug/sprint\_f1\_maxn3\_predictions\_test\_memo.md}) returned
\textbf{CLEAN NEGATIVE}: 1 of 3 predictions pass; P1 (internal PES
minimum) fails on the exact falsification criterion, the proposed
basis-closable cross-shift sub-class evaporates. \emph{The two-bucket M-Z
partition (cross-shift / endomorphism) stands as the canonical framing.}

The substantive structural finding the test produced is a sharper
localization of chemistry-side endomorphisms: at max\_n=3 under
combined W1c + multi-zeta architecture, the framework DOES construct
the right bonding orbital (dominant natural orbital $-0.698 \cdot
\phi_\text{Na, 3s} - 0.687 \cdot \phi_\text{H, 1s}$ with $50/50$
amplitude split), but the constructed bonding combination has higher
energy than the separated configuration at every $R$. The framework can
BUILD the right orbital but cannot energetically PREFER it. Sub-layer~3
of the W1c hierarchy reclassifies from ``basis-closable cross-shift''
to ``module endomorphism (basis-irreducible at tested scales)'' and joins
sub-layer~2 in the endomorphism bucket.

\textbf{Same-day same-arc follow-on sprints (F2 + F3) sharpen the
structural reading further into an architectural-extension ladder.}
The F1 max\_n=3 verdict ``module endomorphism, basis-irreducible at
tested scales'' was further refined by two same-day same-arc follow-on
sprints (F2 + F3) executing the diagnostic-before-engineering rule at
the sprint-cycle level. \textbf{F2} (cross-$V_{ne}$ kernel-shape
substitution diagnostic, KERNEL-NOT-IT verdict): the multipole expansion
at default $L_\text{max}=4$ is bit-faithful to converged 3D numerical
quadrature within $\sim 2\times 10^{-5}$~Ha at NaH $R_\text{eq}=3.566$~bohr,
six to ten orders of magnitude below wall depth. Ruled out kernel-shape
substitution as a closure mechanism. Substantive new content: the
framework's h1 in \texttt{composed\_qubit.build\_composed\_hamiltonian}
is strictly block-diagonal --- the cross-block off-diagonal h1 slot is
\emph{architecturally absent}, not just kernel-error-flawed. New sub-wall
name \textbf{W1d} (cross-block h1 architectural extension), structurally
distinct from a missing calibration input. \textbf{F3} (cross-block h1
architectural extension): implemented the missing matrix slot via new
production module \texttt{geovac/cross\_block\_h1.py}; FCI naturals at
NaH max\_n=2 R=3.5 transform $[1.0000, 1.0000] \to [\mathbf{1.9991},
\mathbf{0.0007}]$ --- four-orders-of-magnitude advance in dominant
natural-orbital occupation signature, with the bonding combination as a
doubly-occupied FCI ground state. \textbf{W1d CLOSED at the FCI level.}
But F3 mini-PES remains monotonically descending with $D_e^\text{F3} =
+4.37$~Ha = $58\times$ experimental: cross-block h1 enables bonding
without supplying opposing Pauli repulsion against the [Ne] frozen core.
\textbf{New sub-layer W1e (inner-region overattraction)}, F4 target via
a bonding-orbital Phillips-Kleinman extension.

\emph{Net refined reading of the chemistry sub-walls under M-Z:} chemistry
sub-walls split into THREE structurally distinct sub-categories within
the two-bucket M-Z partition: cross-shifts (handled by existing
architecture), module endomorphisms (split chemistry-side mitigable
vs QED-side strictly-external), and \textbf{architectural-absence}
(missing matrix slots in the existing architecture, conditionally
closable via architectural extension). The architectural-absence
category is a structurally distinct refinement WITHIN the two-bucket
partition, not a contradiction of it. The architectural-extension
ladder (W1c-cross-screening $\to$ W1c-multi-zeta-basis $\to$
W1d-cross-block-h1 $\to$ W1e?) is the chemistry-side analog of refinements
that the QED side does not natively support: QED endomorphisms (LS-8a
$Z_2 - 1$, $\delta m$; Sprint H1 Yukawa) are strictly externally
specified via UV-completion physics, while chemistry endomorphisms are
intrinsically external but framework-internally computable (FrozenCore
+ atomic-physics fits) AND admit architectural extensions (cross-block
h1, bonding-orbital PK). This explains why the W1c arc produced a
five-sub-layer hierarchy (W1c, W1c-multi-zeta, W1d, W1e, plus the
FALSIFIED three-bucket M-Z refinement candidate) while the QED-side
LS-8a wall has remained at a single-counterterm-input characterization.

\emph{Sources:}
\texttt{debug/sprint\_beta\_neg\_synthesis\_memo.md} (Sprint $\beta$-neg
synthesis, F1-maturity; superseded for framing by the full-arc memo);
\texttt{debug/sprint\_f2\_cross\_vne\_kernel\_memo.md} (F2 KERNEL-NOT-IT
+ architectural-absence finding, W1d named);
\texttt{debug/sprint\_f3\_cross\_block\_h1\_memo.md} (F3 W1d CLOSED at
FCI level + W1e newly named);
\texttt{debug/sprint\_w1c\_full\_arc\_synthesis\_memo.md} (comprehensive
full-arc synthesis at F3 maturity, supersedes Sprint $\beta$-neg
synthesis at the framing level). Cross-references Paper~19
\S\texttt{sec:w1c\_residual} (W1c-residual wall + F1-F2-F3 arc extension
with revised five-sub-layer hierarchy), Paper~17 \S6.10 (composed-architecture
cross-reference, updated to F3 maturity), Paper~34
\S\texttt{sec:conv\_three\_bucket\_falsified} (running-catalogue entry,
parent), and Paper~34 \S\texttt{sec:conv\_w1c\_f2\_f3} (running-catalogue
entry, F2+F3 follow-on).

\textbf{F4--F6 refinement of W1e characterization
($\beta$-update, 2026-05-23).} Three same-day diagnostic sprints
(F4, F5, F6) further refined the W1e (inner-region overattraction)
sub-layer characterization beyond F3's working hypothesis of
``missing Pauli repulsion against the [Ne] frozen core,'' tested
all three candidate closure mechanisms F3 §6 had named, and
sharpened the chemistry-side sub-mechanism dissection of W1e in
the M-Z partition. \textbf{F4} (bonding-orbital
Phillips--Kleinman extension): RULED OUT at Step~1 diagnostic ---
rank-1 PK on the F3-constructed bonding orbital saturates at
${\sim}43\%$ closure ceiling in the FCI sensitivity test, far
below the $96\%$ required for the 2$\times$ experimental closure
window. \textbf{W1e is therefore NOT a single-particle Pauli-
orthogonality wall} (which would have been the cross-shift-class
mechanism within M-Z); it is a multi-determinant FCI correlation
effect. \textbf{F5} (explicit-core Hartree treatment of [Ne]
electrons via cross-block $\sum_c (J - K)$): RULED OUT at Step~1
diagnostic --- right SIGN (repulsive, $+1.12$~Ha) but only
$25.7\%$ of wall depth; predicted residual $D_e^\text{F5} \approx
+3.25$~Ha = 43$\times$ experimental. \textbf{W1e is NOT
addressable by mean-field Hartree-class core-bonding interaction}
(which would have been closable by extended-bimodule machinery
of standard-Hartree class operators within M-Z). \textbf{F6}
(valence basis enlargement to NaH max\_n=4 with full F3 stack):
PARTIAL\_IMPROVEMENT --- $10.2\%$ PES-level closure (the 2-point
gate suggests $26.1\%$ but is artifactual; the inner-region collapse
continues from $R = 3.5$ down to $R = 2.0$ with another $0.696$~Ha
descent the 2-pt gate misses). \textbf{W1e is NOT addressable by
basis enlargement alone at max\_n=4} (which would have been a
basis-truncation cross-shift candidate within M-Z).

\emph{Net structural reading of the chemistry-side sub-mechanism
dissection of W1e under M-Z:} the chemistry-side dissection of W1e
refines the framework's understanding of where ``correlation'' sits
in the M-Z partition. W1e correlation is neither single-particle
Pauli (which is the cross-shift-class mechanism within M-Z, ruled
out by F4) nor mean-field Hartree-level (which would be closable
by extended-bimodule machinery for standard-Hartree-class operators,
ruled out by F5), but \textbf{genuine multi-determinant content
that requires multi-week+ architectural lift} (Schmidt
orthogonalization at the basis-construction level, DMRG-class
beyond-FCI methods, or explicit correlation factors --- all are
multi-week+ architectural investments). W1e characterizes a
\textbf{three-sub-sub-mechanism decomposition with
structurally-independent closure ceilings}: basis-truncation
(${\sim}10\%$ PES, F6 measured) + Hartree-pressure (${\sim}25\%$
2-pt, F5 predicted) + deep-correlation-residual (${\sim}65$--$90\%$
of wall, untested by sprint-scale compute). The chemistry-side
architectural-extension ladder (W1c-cross-screening $\to$
W1c-multi-zeta-basis $\to$ W1d-cross-block-h1) has plausibly hit a
natural pause point at W1e: the next architectural extension
candidates are multi-week+ commitments rather than sprint-scale
diagnostic-then-implementation cycles, and the chemistry arc's
forward direction shifts from ``find the right closure mechanism''
to ``characterize the framework's structural-skeleton scope limits
empirically for second-row hydride binding'' --- consistent with
the broader CLAUDE.md~§1.7 multi-focal-composition wall taxonomy
finding (the framework's structural-skeleton scope does not
autonomously deliver binding energetics at this precision).

\emph{Methodological side finding for future M-Z sub-mechanism
sprints.} The F6 2-point-vs-PES discrepancy reveals that 2-point
gates at fixed reference geometry ($R = 3.5$~bohr vs $R = 10$~bohr)
overestimate PES-derived closure by ${\sim}2.5\times$ on
monotonically-descending PES in the W1c-residual regime. Future
sprint-cycle diagnostics in this regime should use PES well depth
at $R_\text{min}$ as the load-bearing closure metric rather than
2-point energy differentials. The F3-reported $D_e^\text{F3} =
+4.37$~Ha is itself a 2-pt value; the corresponding F3 PES well
depth at $R_\text{min} = 2.0$ is likely larger (untested but
expected by analogy to F6). This is a methodological refinement
of the diagnostic-before-engineering discipline (CLAUDE.md~§1.7)
at the sprint-cycle level: not just ``run a diagnostic before
production wiring,'' but also ``choose the diagnostic geometry
that captures the load-bearing physical content rather than a
convenient reference.''

\emph{Sources (F4--F6 refinement):}
\texttt{debug/sprint\_f4\_bonding\_pk\_memo.md} (F4 STOP-at-Step-1,
rank-1 PK insufficient);
\texttt{debug/sprint\_f5\_explicit\_core\_memo.md} (F5
STOP-at-Step-1, mean-field Hartree insufficient);
\texttt{debug/sprint\_f6\_maxn4\_nah\_memo.md} (F6
PARTIAL\_IMPROVEMENT, basis enlargement to max\_n=4);
\texttt{debug/sprint\_beta\_update\_f4f6\_memo.md} ($\beta$-update
F4-F6-APPENDED-TO-F3-BASELINE synthesis). Cross-references
Paper~19 \S\ref{sec:w1e_refinement_f4_f6} (refined
six-sub-layer hierarchy with three W1e sub-sub-mechanisms),
Paper~17 \S6.10 (composed-architecture cross-reference, updated
to $\beta$-update F4--F6 maturity), Paper~34
\S\texttt{sec:conv\_w1e\_f4\_f6\_refinement} (running-catalogue
entry, F4--F6 follow-on to F3 entry).

\textbf{W1e sprint-scale candidate space closed at Schmidt + core
correlation Day-1 diagnostics (2026-05-23).} The two multi-week+
candidates F4--F6 named (Schmidt orthogonalization at
basis-construction level; fully correlated [Ne] cores) were tested
by Day-1 diagnostics. \textbf{Both ruled out.} Schmidt-orthogonalizing
H 1s against the Na [Ne] core at $R = 3.5$~bohr produces
$\sum_c |S_c|^2 = 1.00 \times 10^{-3}$ (only $0.1\%$ projection mass;
the H atom at $R = 3.5$ is structurally outside the [Ne] core's
significant amplitude) and a total diagonal-element differential
$\Delta_\text{total} = +8.6$~mHa --- $\sim 12\times$ below the
prespecified GO threshold and $\sim 500\times$ below the wall depth.
[Ne] core-correlation differential at NaH $R_\text{eq}$ converges
across three independent literature anchors (Iron-Oren-Martin~2003,
Sylvetsky-Martin~2018, atomic Na correlation $\times$ active-fraction)
to $\sim 0.001$--$0.003$~Ha, $10^3\times$ below the wall depth.
\emph{Five independent sprint-scale W1e closure mechanisms are now
ruled out} (F4, F5, F6, Schmidt, core correlation). The chemistry-side
architectural-extension ladder
(W1c-cross-screening $\to$ W1c-multi-zeta-basis $\to$
W1d-cross-block-h1) has reached the sprint-scale boundary; further
closure candidates are multi-month commitments. The W1e wall is the
\textbf{sixth structurally-independent instance of the
multi-focal-composition wall pattern} across the framework, joining
five physics observables (Sprint H1 Yukawa, LS-8a $Z_2/\delta m$,
HF-3 recoil, HF-4 Zemach, HF-5 multi-loop $a_e$). \emph{The pattern's
cross-domain transferability --- five physics + one chemistry, two
structurally distinct calibration-data classes hitting the same
structural-skeleton-scope boundary --- is the substantive new content
of this closeout, extending the structural-skeleton scope statement
from physics-only to physics+chemistry.} \emph{Sources:}
\texttt{debug/sprint\_w1e\_schmidt\_diagnostic\_memo.md} (Schmidt
Day-1 STOP),
\texttt{debug/sprint\_w1e\_core\_correlation\_estimate\_memo.md}
(core correlation DEFER). Cross-references Paper~19
\S\ref{sec:w1e_schmidt_core_correlation} and Paper~34
\S\ref{sec:conv_w1e_cross_domain_wall}.

% -------------------------------------------------------------------
\subsection{Sprint G3:\ Electroweak chirality co-location on $S^3$}
\label{sec:g3}
% -------------------------------------------------------------------

Sprint~H1 (Sec.~\ref{sec:higgs_h1}) closed G2 up to the
Yukawa-undetermined verdict and left G3 (electroweak unification,
requiring identification of $\gamma_{\mathrm{GV}}$ with the SM
weak-isospin chirality on $\mathcal{H}_F$) as the next physics
question.  The G3 sprint (May~2026) is a constructive test:\ build
$\gamma_{\mathrm{GV}}$ on the GeoVac side and $\gamma_F$ on the AC
side as separate $\mathbb{Z}_2$ gradings each satisfying the
standard NCG axioms, lift both to the tensor-product Hilbert space
$\mathcal{H}_{\mathrm{GV}} \otimes \mathcal{H}_F$, and ask whether
the residual

\begin{equation}
\Delta := \gamma_{\mathrm{GV}} \otimes \mathbb{1}_F -
          \mathbb{1}_{\mathrm{GV}} \otimes \gamma_F
\end{equation}

vanishes (or vanishes on the physical sector under standard
left/right convention swap).  $\Delta = 0$ would close G3 in the
positive; $\Delta \neq 0$ robustly would close G3 in the negative
on $S^3$ alone and push electroweak unification onto G4
cross-manifold territory.

\paragraph{Constructions.}  The sprint executed three parallel
tracks (memos
\texttt{debug/g3a\_chirality\_memo.md},
\texttt{debug/g3b\_chirality\_F\_audit\_memo.md},
\texttt{debug/g3c\_tensor\_chirality\_memo.md}).  Track~G3-A
constructed
$\gamma_{\mathrm{GV}} = \sigma_x \otimes \mathbb{1}_w$ on the
R3.5 full Dirac operator system at $n_{\max} \in \{1, 2, 3\}$, with
$w = \mathrm{spinor\_dim}(n_{\max})$ and the basis ordering
$[\chi = +1, \chi = -1]$ from \texttt{full\_dirac\_basis}.  All three
$\mathbb{Z}_2$-grading axioms are bit-exact at every tested
$n_{\max}$:\ $\gamma_{\mathrm{GV}}^2 = \mathbb{1}$,
$\{\gamma_{\mathrm{GV}}, D_{\mathrm{truthful}}\} = 0$, and
$\gamma_{\mathrm{GV}} M \gamma_{\mathrm{GV}} = M$ as a matrix
equality for every scalar multiplier $M \in \mathcal{O}_{n_{\max}}$.
The natural $\sigma_z$ candidate (the $D$-eigenvalue sign) commutes
with $D_{\mathrm{truthful}}$ rather than anticommuting and is exposed
only as a diagnostic.  Module
\texttt{geovac/chirality\_grading.py} (50/50 tests passing).

Track~G3-B added \texttt{ElectroweakFiniteTriple.chirality\_F()} to
\texttt{geovac/almost\_commutative.py} with the standard
Connes--Marcolli SM convention at KO-dim~6 of $\mathcal{T}_F$:

\begin{equation}
\gamma_F^{\mathrm{matter}} = \mathrm{diag}(+1, +1, -1, -1), \qquad
\gamma_F^{\mathrm{anti}} = -\gamma_F^{\mathrm{matter}}.
\end{equation}

The antimatter sign-flip is forced by $\{J_F, \gamma_F\} = 0$:\ the
naive ``identical on both copies'' convention would put
$\mathcal{T}_F$ at KO-dim~0 and silently break the H1 module's
$\{J_F, D_F\}$ axiom suite.  All grading axioms
($\gamma_F^2 = \mathbb{1}$, $\{\gamma_F, D_F\} = 0$ at any imposed
Yukawa, $\gamma_F \pi_F(a) \gamma_F = \pi_F(a)$,
$\gamma_F = \gamma_F^*$, $\{J_F, \gamma_F\} = 0$) hold to bit-exact
zero residual; 15 new tests in \texttt{TestChiralityF}, all passing.
The combined KO-dimension of
$\mathcal{T}_{\mathrm{AC}} = \mathcal{T}_{\mathrm{GV}} \otimes
\mathcal{T}_F$ is therefore $3 + 6 = 9 \equiv 1 \pmod 8$, consistent
with H1's combined-Dirac construction.

\paragraph{Tensor-product diagnostic.}  Track~G3-C computed $\Delta$
at $n_{\max} \in \{1, 2, 3\}$ using the production
$\gamma_{\mathrm{GV}}$ and $\gamma_F$ from G3-A and G3-B.  Sanity
invariants are bit-exact:\ $\gamma_{\mathrm{GV}}^2 = \mathbb{1}$,
$\gamma_F^2 = \mathbb{1}$, and
$[\gamma_{\mathrm{GV}} \otimes \mathbb{1}_F,
\mathbb{1}_{\mathrm{GV}} \otimes \gamma_F] = 0$ (the two operators
commute because they act on disjoint tensor factors).  The residual
itself is the structurally cleanest possible non-zero object:

\begin{table}[h]
\centering
\small
\begin{tabular}{|c|c|c|c|c|c|}
\hline
$n_{\max}$ & $\dim \mathcal{H}_{\mathrm{GV}}$ &
$\dim \mathcal{H}_{\mathrm{total}}$ &
$\|\Delta\|_{\mathrm{op}}$ & $\|\Delta\|_F$ &
$\mathrm{spec}(\Delta)$ \\
\hline
$1$ & $4$  & $32$  & $2.000$ & $8.000$  &
$\{-2{:}8,\ 0{:}16,\ +2{:}8\}$    \\
$2$ & $16$ & $128$ & $2.000$ & $16.000$ &
$\{-2{:}32,\ 0{:}64,\ +2{:}32\}$  \\
$3$ & $40$ & $320$ & $2.000$ & $25.298$ &
$\{-2{:}80,\ 0{:}160,\ +2{:}80\}$ \\
\hline
\end{tabular}
\caption{Tensor-product chirality residual $\Delta =
\gamma_{\mathrm{GV}} \otimes \mathbb{1}_F -
\mathbb{1}_{\mathrm{GV}} \otimes \gamma_F$ at
$n_{\max} \in \{1, 2, 3\}$, computed in
\texttt{debug/g3c\_tensor\_chirality.py} from the production
$\gamma_{\mathrm{GV}}$ (G3-A, $\sigma_x \otimes \mathbb{1}_w$) and
$\gamma_F$ (G3-B, KO-dim~6).  Operator norm is exactly~$2$ at every
$n_{\max}$; eigenvalue spectrum is $\{-2, 0, +2\}$ with
multiplicities in ratio $1{:}2{:}1$, $n_{\max}$-independent.
Frobenius norm scales as
$\sqrt{2 \cdot \dim \mathcal{H}_{\mathrm{total}}}$ exactly.  Data
\texttt{debug/data/g3c\_tensor\_chirality.json}.}
\label{tab:g3c_residual}
\end{table}

\paragraph{Convention swaps do not absorb the residual.}  Three
standard convention swaps were tested at every $n_{\max}$:\
(a)~global sign-flip $\gamma_F \to -\gamma_F$ (left/right
convention), (b)~restriction to the matter sector of
$\mathcal{H}_F$, and (c)~restriction to the Weyl chirality of
$\mathcal{H}_{\mathrm{GV}}$.  Swap~(a) leaves the baseline invariant.
Swaps~(b) and~(c) reduce $\|\Delta\|_F$ by dropping basis vectors,
but $\|\Delta\|_{\mathrm{op}}$ remains~$2$ (or~$1$ under swap~(c)
with $\sigma_x$-$\gamma_{\mathrm{GV}}$, where the Weyl restriction
collapses $\gamma_{\mathrm{GV}}|_w$ to $0_w$ and $\Delta_c =
-\mathbb{1}_w \otimes \gamma_F$, still non-zero).  Combined swaps
yield the same conclusion.  No standard NCG sign convention absorbs
$\Delta$.

\paragraph{Robustness.}  A $\sigma_x$-vs-$\sigma_z$ cross-check at
$n_{\max} = 2$ confirms that the spectrum of $\Delta$ is identical
under both $\gamma_{\mathrm{GV}}$ conventions (both produce
$\{-2{:}32,\ 0{:}64,\ +2{:}32\}$, since $\sigma_x$ and $\sigma_z$
share eigenvalues $\pm 1$ with identical multiplicities and both
commute with $\gamma_F$).  The structural verdict is therefore
convention-independent.

\paragraph{Sanity check on $\gamma_5 := \gamma_{\mathrm{GV}} \otimes
\gamma_F$.}  The product operator
$\gamma_5 := \gamma_{\mathrm{GV}} \otimes \gamma_F$ is the canonical
Connes--Chamseddine combined chirality on $\mathcal{T}_{\mathrm{AC}}$
and satisfies $\gamma_5^2 = \mathbb{1}$ to bit-exact precision at
every $n_{\max}$, with spectrum $\{-1, +1\}$ in equal multiplicities.
The product works; the \emph{difference} does not.  $\gamma_5$ exists
and is well-defined as the combined grading; $\gamma_{\mathrm{GV}}$
and $\gamma_F$ do not identify across the tensor product.

\paragraph{Verdict:\ NEGATIVE on $S^3$ alone.}  The two chirality
gradings $\gamma_{\mathrm{GV}}$ (Camporesi--Higuchi Dirac sign on
$S^3$) and $\gamma_F$ (weak-isospin grading on the AC factor) are
\emph{independent commuting} $\mathbb{Z}_2$ gradings on
$\mathcal{H}_{\mathrm{GV}} \otimes \mathcal{H}_F$.  They cannot be
identified by any standard NCG sign convention because, by
construction, $\gamma_{\mathrm{GV}}$ acts as identity on
$\mathcal{H}_F$ and $\gamma_F$ acts as identity on
$\mathcal{H}_{\mathrm{GV}}$;\ no basis change on the tensor product
turns one into the other.  Electroweak chirality co-location does
not close on $S^3$ alone.

\paragraph{The G2--G3 corollary.}  Sprint~H1 left G2 (Higgs as inner
fluctuation) at the verdict ``construction admits a Higgs given an
imposed Yukawa $Y$, but no GeoVac mechanism selects $Y$.''  G3's
NEGATIVE supplies the structural reason.  The Yukawa sits in the
off-diagonal $\mathbb{C} \leftrightarrow \mathbb{H}$ block of $D_F$
that anticommutes with $\gamma_F$;\ the candidate GeoVac sources for
that block (offdiag CH chirality bridging, Sturmian inter-shell
coupling) have no kinematic relation to $\gamma_F$, because
$\gamma_{\mathrm{GV}}$ is structurally independent of $\gamma_F$.
Equivalently:\ no GeoVac-side data fixes $Y$ because the GeoVac
chirality lives on a different $\mathbb{Z}_2$ grading from the one
the Yukawa flips.  G2 and G3 are therefore one structural fact, not
two.  This is a \emph{sharpening} of H1's G2 verdict (the
Yukawa-undetermined gap is named structurally), not a new gap.

\paragraph{Files.}  Modules
\texttt{geovac/chirality\_grading.py} (G3-A, 50/50 tests) and the
\texttt{chirality\_F} additions to
\texttt{geovac/almost\_commutative.py} (G3-B);\ tests
\texttt{tests/test\_chirality\_grading.py} and
\texttt{tests/test\_almost\_commutative.py} (53 passing,
\texttt{TestChiralityF} block);\ diagnostic
\texttt{debug/g3c\_tensor\_chirality.py};\ data
\texttt{debug/data/g3a\_chirality.json},
\texttt{debug/data/g3c\_tensor\_chirality.json};\ memos
\texttt{debug/g3a\_chirality\_memo.md},
\texttt{debug/g3b\_chirality\_F\_audit\_memo.md},
\texttt{debug/g3c\_tensor\_chirality\_memo.md}.

% -------------------------------------------------------------------
\subsection{Sprint G4:\ Cross-manifold scoping (G4a reachable, G4b
NCG-framework-blocked)}
\label{sec:g4}
% -------------------------------------------------------------------

With G2 and G3 collapsed into one structural fact and electroweak
unification on $S^3$ alone closed in the negative, G4
(cross-manifold) becomes \emph{the} remaining route to a
GeoVac-internal SM construction rather than the dominant gap among
several.  A scoping pass (May~2026, memo
\texttt{debug/g4\_cross\_manifold\_scoping\_memo.md}) sharpens G4
into two sub-gaps with sharply different reachability.

\paragraph{G4a (reachable):\ Connes SM on $\mathcal{T}_{S^3}$ alone
with $\mathcal{A}_F = \mathbb{C} \oplus \mathbb{H} \oplus
M_3(\mathbb{C})$.}  The minimal extension of Sprint~H1's electroweak
slice $\mathcal{A}_F = \mathbb{C} \oplus \mathbb{H}$ to the full
Connes SM finite algebra is reachable as a 1--2~month sprint
extending \texttt{geovac/almost\_commutative.py}.  The
$M_3(\mathbb{C})$ summand supplies $\mathrm{SU}(3)$ at the
inner-fluctuation level on the existing GeoVac $S^3$ triple, without
recourse to the $S^5$ Bargmann construction.  Predicted
outcome:\ \emph{positive-thin}, in the exact sense of Sprints~H1
and~G3.  Inner fluctuations decompose into a gauge-1-form sector
(recovering $U(1) \times \mathrm{SU}(2) \times \mathrm{SU}(3)$ at the
operator level, with the abelian and $\mathrm{SU}(2)$ parts already
verified by H1) plus a Higgs sector that is non-trivial whenever the
Yukawa-Dirac off-diagonal blocks are imposed.  The G2--G3 corollary
applies verbatim:\ G4a defines the construction, but no GeoVac
mechanism selects the Yukawa, hypercharge assignments, or
generations.  G4a therefore does not close the SM-derivation
question;\ it relocates the open data-selection items
$(Y, \text{hypercharge}, \text{generation count})$ from
``electroweak slice'' (H1, G3) to ``full Connes finite algebra,''
without changing their status.  G4a remains worth pursuing because
it makes Sprint~ST-SU3's Bargmann SU(3) construction structurally
redundant for SM-derivation purposes (one $S^3$ construction with
$M_3(\mathbb{C})$ delivers the same gauge content) and because it
is the most published-precedent-compatible target available.

\paragraph{G4b (NCG-framework-blocked):\ genuine cross-manifold
$\mathcal{T}_{S^3} \otimes \mathcal{T}_{S^5}$.}  The structurally
deeper version of G4 asks whether the GeoVac $S^3$ Coulomb triple
and the Paper~24 $S^5$ Bargmann lattice can be combined into a
single spectral triple with both as commuting factors, so that the
SU(3) of Sprint~ST-SU3 is honestly cross-manifold-located rather
than recapitulated on $S^3$.  This sub-gap is blocked at the
\emph{NCG-framework level}, not at the GeoVac level.  Connes' tensor
product of spectral triples~\cite{connes_marcolli2008} requires
both factors to be Riemannian spectral triples in the same sense.
$\mathcal{T}_{S^3}$ carries the Camporesi--Higuchi Riemannian spinor
Dirac (KO-dim~3).  Paper~24's $S^5$ Bargmann construction is built
from the holomorphic Hardy sector $H^2(S^5)$ restricted to symmetric
SU(3) irreps $(N, 0)$;\ its natural diagonal operator is the
first-order complex Euler operator
$\hat{N} = \sum_i z_i \partial_i$, not a Riemannian Dirac.  Two
candidate $\mathcal{T}_{S^5}$ exist, and neither preserves both
standard NCG axioms and the GeoVac-specific content of Sprints
25--30 / ST-SU3:

\begin{itemize}\setlength{\itemsep}{2pt}
\item \emph{Option~A (Riemannian round-$S^5$).}  Take $D_{S^5}$ to
be the standard Riemannian spinor Dirac on round $S^5$ (KO-dim~5).
The tensor product is a textbook Connes object with combined
KO-dim $3 + 5 = 8 \equiv 0 \pmod 8$, but it discards the
$(N, 0)$-tower structure that made Paper~24's $S^5$ ``natural,''
and the bit-exact $\pi$-free certificate of Paper~24~\S III is
irrelevant to round-$S^5$ (calibration $\pi$ reappears via standard
Weyl asymptotics).
\item \emph{Option~B (Bargmann--Euler-based).}  Promote the Euler
operator $\hat{N}$ on the Bargmann graph to a Dirac-like object via
chiral doubling.  This preserves Sprint~ST-SU3 SU(3) Wilson content
and the $\pi$-free certificate, but the resulting object is
\emph{not} a Riemannian spectral triple in the published sense:\
KO-dimension is not classically defined for a Hardy-sector
first-order complex operator, and there is no published prescription
in the NCG literature for tensoring a Riemannian spectral triple
with a Hardy-sector / Berezin--Toeplitz first-order complex object.
The closest analogs (Hawkins~2000 on the unit disk;
Hekkelman~2022, 2024; Hekkelman--McDonald~2024 on flat structures)
stay within Berezin--Toeplitz on a single manifold or within
Riemannian on a single manifold, never crossing.
\end{itemize}

The structural reading is that G4b is the Coulomb/HO
asymmetry of Paper~24~\S V resurfacing at the spectral-triple level.
Paper~24~\S 5.3 now records four layers of this asymmetry, all
Coulomb-specific:\ (i)~spectrum-computing role of $L_0$;\
(ii)~calibration~$\pi$;\ (iii)~Wilson gauge with natural matter
coupling, Sec.~\ref{sec:unified_gauge};\ and (iv)~modular-Hamiltonian
Pythagorean orthogonality on the hemispheric wedge, the structural
mechanism behind the L2-E signature-independence finding
(Paper~43~\S 10.2~\cite{paper43};
Remark~\ref{rem:pythagorean_m1_closure} above).  G4b is a
\emph{sibling} of these four layers, not the lone fourth layer:\ the
Coulomb $S^3$ side carries a Riemannian spectral triple in the
standard NCG sense and supports the BW-aligned wedge-modular
construction at signatures $(3, 0)$ and $(3, 1)$, while the Bargmann
$S^5$ side has none of the four ingredients (no spinor bundle on
the $(N, 0)$ Hardy tower;\ no second-order-vs-first-order distinction
with a separate Dirac;\ no half-integer wedge;\ no calibration $\pi$
or natural Wilson matter coupling). The Connes tensor-product
framework cannot bridge them as currently published.  Closing G4b is
therefore not a sprint-scale GeoVac task;\ it requires either an
extension of NCG to handle Riemannian--Hardy mixed tensor products,
or the identification of a unifying ambient manifold whose GeoVac
sub-structure contains both $S^3$ Coulomb and $S^5$ Bargmann as
natural sub-objects.  Both directions are recorded as honest open
structural questions, not as sprint targets.

\paragraph{Net.}  After Sprints H1 + G3 + G4 scoping, the
four-gap analysis of Sec.~\ref{sec:unified_gauge} compresses to:

\begin{itemize}\setlength{\itemsep}{2pt}
\item G2 and G3 are one structural fact:\ the Yukawa is
undetermined because $\gamma_{\mathrm{GV}}$ and $\gamma_F$ are
independent commuting $\mathbb{Z}_2$ gradings on the tensor product
(G3 NEGATIVE).
\item G4 splits cleanly:\ G4a (Connes SM on $\mathcal{T}_{S^3}$
alone with $M_3(\mathbb{C})$) is \textbf{closed at POSITIVE-THIN}
(2026-05-31):\ $\mathrm{U}(1) \times \mathrm{SU}(2) \times
\mathrm{SU}(3)$ recovered from inner fluctuations, all Connes axioms
bit-exact, Yukawa/hypercharge/generations remain imposed;\ G4b
(genuine cross-manifold) is blocked at the NCG-framework level by the
Coulomb/HO structural asymmetry.
\item G1 (three generations) is unchanged from H1:\ formally
definable, not selected by GeoVac structure.
\end{itemize}

GeoVac does not autonomously generate the Standard Model;\ the
construction admits SM gauge content given imposed Yukawa,
hypercharge, and generation choices, and the GeoVac data is silent
on each of those choices.  This is a clean structural verdict, in
the sense of CLAUDE.md~\S 1.5:\ the framework is consistent with a
GeoVac SM construction, but the framework does not select one.

\paragraph{G4a closure (Sprint G4a, 2026-05-31).}
The full Connes SM finite algebra
$\mathcal{A}_F = \mathbb{C} \oplus \mathbb{H} \oplus M_3(\mathbb{C})$
has been implemented and verified on
$\mathcal{T}_{S^3}$ at $n_{\max} \in \{1, 2, 3\}$ with
$\dim \mathcal{H}_F = 32$ (16 matter + 16 antimatter;
$\dim \mathcal{H}_{\text{total}} = 1280$ at $n_{\max} = 3$).
All six Connes axioms ($J^2 = -I$, $JD = +DJ$, $D = D^\dagger$,
$\gamma^2 = I$, order-zero, order-one) pass at literal
machine-zero residual at every tested $n_{\max}$, with no
finite-resolution degradation.
Inner fluctuations decompose into
$\mathrm{U}(1) \times \mathrm{SU}(2) \times \mathrm{SU}(3)$
gauge 1-forms plus a Higgs scalar that is non-trivial whenever the
Yukawa-Dirac off-diagonal blocks are imposed and identically zero
otherwise.  Matter-antimatter and lepton-quark cross-sectors decouple
exactly.  The predicted POSITIVE-THIN verdict is confirmed:
the construction admits the full SM gauge content with imposed
Yukawa, hypercharge, and generation choices, and the GeoVac data is
silent on each of those choices.
Module \texttt{geovac/standard\_model\_triple.py} ($\sim 530$ lines),
45 tests passing.

\paragraph{H1 Higgs inner fluctuation (Sprint H1, May 2026).}
The full fluctuated Dirac $D_A = D + A + JAJ^{-1}$ has been
constructed at $n_{\max} \in \{2, 3\}$
($\dim \mathcal{H}_{\text{total}} = 512, 1280$). Inner
fluctuations decompose cleanly into gauge
($\mathrm{U}(1) \times \mathrm{SU}(2) \times \mathrm{SU}(3)$,
matching G4a) plus Higgs (complex $\mathrm{SU}(2)$ doublet $\Phi$
in the L--R off-diagonal block, non-trivial iff Yukawa $Y \ne 0$).
The spectral-action Mexican-hat potential is confirmed:
$\Tr(D_F^2) > 0$ and $\Tr(D_F^4) > 0$ give the standard CCM
sign structure $V(\Phi) = -\mu^2 |\Phi|^2 + \lambda |\Phi|^4$.

\textbf{Yukawa non-selection theorem.} The space of admissible
$D_F$ satisfying all Connes axioms has $128$ real parameters
after Hermiticity ($1024 \to 512$), chirality grading
($512 \to 128$), and $J$-reality ($128 \to 128$; the order-one
condition does not further constrain). For one generation,
$Y = \mathrm{diag}(y_\nu, y_e, y_u, y_d)$ carries $8$ free
real parameters. The gauge sector is \textbf{forced};
the Yukawa is \textbf{free}. This is grounded in the G3
structural finding: $\gamma_{\text{GV}}$ and $\gamma_F$ are
independent commuting $\mathbb{Z}_2$'s, so no GeoVac-side data
couples to the $\gamma_F$ flip that carries the Yukawa.
Source: \texttt{debug/h1\_higgs\_inner\_fluctuation\_memo.md}.

\paragraph{The forced/free seam (forcing-catalogue Door~4, 2026-06-01).}
The Yukawa non-selection theorem admits a sharper structural reading.
The almost-commutative factorization $D^2 = D_{\mathrm{GV}}^2 \otimes 1
+ 1 \otimes D_F^2$ (the cross term vanishes by chirality
anticommutation) factorizes the master Mellin engine into an
\emph{outer} ring --- the M1/M2/M3 transcendentals of the GeoVac
sector --- tensored with an \emph{inner} Yukawa Dirichlet ring
$\mathbb{Q}[y_i^{-2s}]$ generated by the Yukawa eigenvalues. The two
rings share no generator: the $\eta$-trivialization
($\{\gamma_F, D_F\} = 0$ annihilates the odd-$k$ heat traces that
would couple them), together with the G3 commuting-$\mathbb{Z}_2$
result, makes the inner Yukawa ring \emph{disjoint} from every forced
outer ring. This is the precise content of the structural-skeleton
scope (\S 1.5 of the project guidelines): the bone (the forced,
integer/rational skeleton) and the calibration data (free) live in
provably non-overlapping rings, and the Yukawa is free \emph{because}
it generates its own ring --- a quantity cannot be derived from a ring
it generates. The reading also explains, after the fact, why the H1,
W3, and Koide attempts to select the Yukawa all returned clean
negatives: each probed \emph{inside} the almost-commutative category,
where the seam already forbids selection.

\paragraph{The inner algebra structure is mostly forced (Door~4b,
  2026-06-01).} The seam places the Yukawa \emph{values} on the free
side; it does \emph{not} place the inner \emph{algebra structure}
there. A finite-algebra enumeration shows that among all finite
semisimple real $*$-algebras with at most three summands and factor
size at most three, exactly \emph{two} reproduce the gauge group
$\mathrm{U}(1) \times \mathrm{SU}(2) \times \mathrm{SU}(3)$ as their
unitary/inner-automorphism content: the Standard Model finite algebra
$\mathbb{C} \oplus \mathbb{H} \oplus M_3(\mathbb{C})$ and
$\mathbb{C} \oplus M_2(\mathbb{C}) \oplus M_3(\mathbb{C})$. The same
Bertrand $\times$ complex-Hopf-tower truncation at $n \le 3$ that
forces the gauge content (Sec.~\ref{sec:unified_gauge}) therefore also
forces the factor count ($3$, one per Hopf rung $n = 1, 2, 3$), the
$n = 1$ factor ($\mathbb{C}$), and the $n = 3$ factor
($M_3(\mathbb{C})$); the residual non-uniqueness collapses to a single
binary fork --- $\mathbb{H}$ versus $M_2(\mathbb{C})$ at the $n = 2$
rung. This fork was conjectured (Door~4b) to have a GeoVac-native handle:
the outer triple's $J_{\mathrm{GV}}^2 = -1$ at KO-dimension~$3$ (the
real-structure result of Sec.~\ref{sec:construction}, the pseudoreal
$\mathrm{Spin}(3) = \mathrm{SU}(2)$ two-spinor sign --- the
\emph{same} $S^3 = \mathrm{SU}(2)$ fact that produces the gauge
$\mathrm{SU}(2)$) might select $\mathbb{H}$ over $M_2(\mathbb{C})$. A
direct $J$ sign-table audit (Door~4c, 2026-06-01) closed this
conjecture \emph{negative}: $J_{\mathrm{GV}}^2 = -1$ \emph{admits} but
does not \emph{force} $\mathbb{H}$. The decisive fact is that over
$\mathbb{C}$ the two candidates are the same complex algebra
($\mathbb{H} \otimes_{\mathbb{R}} \mathbb{C} \cong M_2(\mathbb{C})$);
the $\mathbb{H}$-vs-$M_2(\mathbb{C})$ distinction is a choice of
\emph{real form}, fixed by the sign of an \emph{internal}
$\mathbb{C}^2$ antilinear involution
($J_{\mathrm{int}}^2 = -1 \Leftrightarrow \mathbb{H}$;
$J_{\mathrm{int}}^2 = +1 \Leftrightarrow M_2(\mathbb{C})$), which lives
on the algebra side and is invisible to the combined
$J = J_{\mathrm{GV}} \otimes J_F$. The combined $J^2$, the combined
$(\varepsilon, \varepsilon')$ signs, and the combined KO-dimension are
\emph{bit-identical} for both candidates at $n_{\max} \in \{1, 2, 3\}$,
and the Chamseddine--Connes--Marcolli order-one-with-Yukawa selector
is satisfied by both (matter/antimatter decoupling). The $\mathbb{H}$
choice is therefore a \emph{literature import} (the CCM
complex-fermion-representation / second-order condition on the finite
factor), not a GeoVac forcing; the fork stays open. The honest
ledger: \emph{forced} are the factor count and the $\mathbb{C}$ and
$M_3(\mathbb{C})$ summands; \emph{admitted but not forced} is the
$n = 2$ real form ($\mathbb{H}$ vs $M_2(\mathbb{C})$); \emph{free} are
the generation count (invisible to inner automorphisms), the inner
KO-dimension, and the Yukawa values. The net sharpening is real but
bounded: the geometry forces the Standard Model gauge group, the
factor count, and two of the three matrix factors --- more internal
algebraic \emph{shape} than the gauge group alone --- but the third
factor's real form, the generation multiplicity, and the Yukawa values
remain outside the forcing. Sources:
\texttt{debug/door4\_gauge\_yukawa\_boundary\_memo.md},
\texttt{debug/door4b\_inner\_algebra\_forcing\_memo.md},
\texttt{debug/door4c\_j\_signtable\_audit\_memo.md}.

\paragraph{A geometric handle on the $n=2$ fork (Door~4d, 2026-06-02).}
A second handle, structurally orthogonal to the combined-$J$ sign table
of Door~4c, is the \emph{Division-Algebra-of-the-Sphere} (DAS) criterion:
\emph{at each Hopf rung $n$, the inner algebra is the natural associative
real normed division algebra whose unit-norm sphere IS the rung's Hopf
bundle total space $S^{2n-1}$.} The identifications are bit-exact at
all three rungs:\ $S^1$ is the unit-norm sphere of $\mathbb{C}$ under
complex multiplication (closure $1.1 \times 10^{-16}$); $S^3$ is the
unit-norm sphere of $\mathbb{H}$ under quaternion multiplication and is
bit-exactly isomorphic to $\mathrm{Sp}(1) = \mathrm{SU}(2)$ via the
standard $\mathbb{H} \to M_2(\mathbb{C})$ embedding (closure
$2.2 \times 10^{-16}$, $\mathrm{SU}(2)$-isomorphism error
$3.3 \times 10^{-16}$); and by Hurwitz's classification of finite-dimensional
real normed division algebras (only $\mathbb{R}, \mathbb{C}, \mathbb{H}$
are associative, with unit spheres $S^0, S^1, S^3$ respectively), $S^5$
has no associative-division-algebra realization, so the inner algebra at
$n=3$ falls back to a matrix algebra. DAS \emph{agrees} with Door 4b's
elimination forcings at $n=1$ ($\mathbb{C}$) and $n=3$ ($M_3(\mathbb{C})$),
and \emph{closes} the $n=2$ fork by selecting $\mathbb{H}$ --- the
direct realization with $\mathrm{U}(\mathbb{H}) = \mathrm{Sp}(1) =
\mathrm{SU}(2)$ on the 3-dimensional $S^3$, in contrast to the
$M_2(\mathbb{C})$ realization whose unit-norm sphere is the 4-dimensional
$\mathrm{U}(2)$ with $\mathrm{SU}(2)$ reached only via the
$\mathrm{U}(2)/\mathrm{U}(1)$ unimodularity quotient. The DAS argument
never invokes any $J$ or sign-table data; it is a purely geometric
criterion identifying the rung-$n$ sphere with the unit-norm subset of
the rung's associative division algebra. The verdict is
\emph{partial-door}: the existing argument of
Sec.~\ref{sec:unified_gauge} extracts only the gauge group from the
Hopf-rung sphere; adopting DAS as a forcing requires the additional
structural theorem that the AC tensor product
$\mathcal{H} = \mathcal{H}_{\mathrm{GV}} \otimes \mathcal{H}_F$ transfers
the rung-$n$ sphere's division-algebra structure to a constraint on
$\mathcal{H}_F$'s inner algebra. This theorem is sprint-reachable and
paper-level; if it closes, the inner-algebra status upgrades from
``mostly forced with the $\mathbb{H}$-vs-$M_2(\mathbb{C})$ fork closed
by literature import (CCM second-order condition)'' to ``mostly forced
with the fork closed by a GeoVac-coherent geometric criterion,'' leaving
only the generation count, the inner KO-dimension, and the Yukawa values
outside the forcing. Source:
\texttt{debug/door4d\_division\_algebra\_sphere\_memo.md}.

\paragraph{Falsifier A$'$ on DAS (Door 4e, 2026-06-02).}
The DAS handle introduced above is partial-door because the existing
construction of Sec.~\ref{sec:unified_gauge} extracts only the gauge
group from the Hopf-rung sphere. The natural test was whether any
standard AC tensor-product compatibility condition fills the gap, i.e.,
distinguishes $\mathbb{H}$ from $M_2(\mathbb{C})$ as the $n=2$ inner
algebra. Three candidates were tested literally:\ (C1)
SU(2)-equivariance under the Ad action $g \cdot m \cdot g^{-1}$;\
(C2) Hopf U(1)-equivariance under scalar phase $e^{i\theta}I$ acting
from the right;\ (C3) principal-bundle compatibility with the Hopf U(1)
fiber. All three returned \emph{negative}: (C1) both $\mathbb{H}$ and
$M_2(\mathbb{C})$ are SU(2)-Ad-invariant real $*$-subalgebras of
$M_2(\mathbb{C})$ (the former is the unique minimal non-abelian such,
the latter is the maximal; the distinction requires a minimality
criterion that is itself extra input); (C2) every $2 \times 2$ complex
matrix commutes with the scalar phase $e^{i\theta}I$ trivially, so the
condition is vacuous; (C3) both candidate actions are
principal-bundle compatible at machine precision by associativity of
matrix multiplication and centrality of scalar phases. \emph{No standard
AC tensor-product compatibility condition forces} $\mathbb{H}$
\emph{over} $M_2(\mathbb{C})$. The substantive content of the test is
not the negative outcome (which was the expected one) but the
\emph{import comparison} it enables:\ the standard
Chamseddine--Connes--Marcolli derivation of $\mathbb{H}$ at the $n=2$
factor invokes \emph{three} axioms beyond the bare Connes spectral-triple
data (second-order condition with nonzero finite Dirac, complex chiral
fermion representation, and the $2N^2 = 32$ dimension count), whereas
DAS invokes \emph{one} (the rung-$n$ sphere's division-algebra
structure is inherited by the inner algebra). Both are imports relative
to the bare axioms; DAS is leaner. The path to upgrading partial-door
to door is therefore not a theorem about existing CCM machinery but a
paper-level structural axiom:\ adopting ``the inner algebra at rung $n$
is the $*$-algebra whose unit group IS the rung-$n$ sphere $S^{2n-1}$ as
a Lie group, when one exists; else fallback $M_n(\mathbb{C})$.'' This
reads the Hopf-rung sphere fully (Lie-group structure plus topology)
where the existing construction reads it partially (gauge group only),
and gives $\mathbb{C}$ at $n=1$ (via $S^1 = \mathrm{U}(1) =
\mathrm{U}(\mathbb{C})$), $\mathbb{H}$ at $n=2$ (via $S^3 =
\mathrm{Sp}(1) = \mathrm{U}(\mathbb{H})$), and $M_3(\mathbb{C})$ at
$n=3$ (since $S^5$ is not a Lie group at all, by Adams 1958). Source:
\texttt{debug/door4e\_falsifier\_A\_prime\_memo.md}.

\paragraph{Closure of the Door 4 series (Door 4f, 2026-06-02; net axiom
savings audited).} The natural next test was whether adopting Upgrade B
introduces any new downstream constraint on existing GeoVac observables.
Twelve sub-tests were run:\ inner KO-dimension (T1), chirality grading
$\gamma_F$ (T2), Yukawa eigenvalues (T3), generation count (T4),
Connes order-zero (T5), order-one (T6), and second-order (T7)
conditions, bosonic spectral-action gauge-sector coefficient (T8),
Higgs vacuum manifold (T9), cross-rung Yukawa coupling (T10),
extensibility to rung $n \ge 4$ via Adams 1958 (T11), and consistency
with the Bertrand $\times$ Hopf-tower reading of
Sec.~\ref{sec:unified_gauge} (T12). \emph{All twelve} return either
SILENT (Upgrade B unconstrained on the target), COMPATIBLE (the
pre-existing constraint passes for $\mathbb{H}$, the mechanism
unaffected), IDENTICAL to CCM (same observable from either
$\mathbb{H}$ or $M_2(\mathbb{C})$ post-unimodularity), or a CLEAN
extensibility/strengthening rule (T11, T12). Zero sub-tests introduce
a new constraint; zero contradict any existing GeoVac result. Upgrade B
is therefore a \emph{clean reformulation} of the CCM $\mathbb{H}$-selection
with no new predictive content.

\textbf{Honest axiom-savings audit.} The naive comparison in the
preceding paragraph (CCM 3 axioms vs DAS 1 axiom) over-counts. Sub-test
T7 (Connes second-order condition) shows that the second-order
condition is \emph{still} required in the Upgrade-B formulation,
because it constrains the off-diagonal Yukawa structure of $D_F$
(which $\mathbb{H}$-selection alone does not determine).
Upgrade B therefore replaces \emph{two} of CCM's three
$\mathbb{H}$-selection axioms --- the complex chiral fermion
representation requirement and the $2N^2 = 32$ dimension count ---
with the single sphere-Lie-group axiom; the second-order condition
persists in both formulations. Net axiom savings: \emph{one}. The
Door~4 series is structurally complete at \textbf{partial-door (final)}.
The remaining open questions are paper-level (whether to adopt
Upgrade B's axiom-minimality advantage as the working construction
principle) and the deep wall (force $N_{\mathrm{gen}}$ and the inner
KO-dim from a packing principle; no known handle; deferred). Source:
\texttt{debug/door4f\_falsifier\_A\_double\_prime\_memo.md}.

\paragraph{The deep wall is a relative necessity, not a gap (Direction~2
  scoping, 2026-06-03).} The ``deep wall'' named above --- forcing
$N_{\mathrm{gen}}$ and the inner KO-dimension from a packing principle ---
is sharper than an open gap:\ relative to the packing construction of
Paper~0 as presently formulated, it is a \emph{necessity}. Two
established facts compose. First, the Door~4 seam:\ the inner
Yukawa/real-structure data lives in a ring (and a set of
$(\mathcal{H}_F, D_F, J_F, \gamma_F)$ data) provably disjoint from every
forced outer ring. Second, Paper~0~\S VII:\ the packing construction is
\emph{kinematic} --- it emits the label set $(n, \ell, m, s)$ and the
graph topology, and emits no operator, no real structure, no coupling,
and no Hilbert-space multiplicity. But $N_{\mathrm{gen}}$ is exactly a
Hilbert-space multiplicity (the factor in $\mathcal{H}_F =
\mathbb{C}^{N_{\mathrm{gen}}} \otimes \mathcal{H}_F^{(1)}$, invisible to
every algebra- and automorphism-level mechanism), and the inner
KO-dimension is exactly a $(J_F, \gamma_F)$ datum --- both of the type
the packing construction does not produce. The packing axiom therefore
cannot constrain either:\ the residue is packing-inaccessible by
composition of these two facts. The single logically-open escape is a
\emph{packing-native multiplicity} --- a datum the construction emits on
the outer labels that would map to a factor-$N_{\mathrm{gen}}$ tripling
of $\mathcal{H}_F$. It is empty:\ Paper~0's entire multiplicity inventory
is the $2\ell+1$ angular slots (odd, and \emph{growing} with $\ell$,
hence not a generation-universal multiplicity) together with a single
$\mathbb{Z}_2$ already consumed by spin; no constant factor of three
appears anywhere. This promotes the structural-skeleton scope from an
observation to a \emph{relative theorem of necessity} for the
$N_{\mathrm{gen}}$/inner-KO-dimension residue, and gives a single
structural reason the H1, W3, and Koide selection attempts failed:\ all
probed inside the almost-commutative category, on the side of the seam
where the disjointness forbids selection and the packing axiom has no
reach. The qualifier ``relative'' is load-bearing:\ the statement
forecloses the residue for the \emph{current} packing construction, not
for every conceivable discrete one. Lifting it would require a
\emph{sibling axiom} that emits fiber multiplicities --- a new primitive
Paper~0 does not contain --- recorded as the open (multi-year) direction
with no present handle. The inner \emph{algebra shape} $\mathbb{C}
\oplus \mathbb{H} \oplus M_3(\mathbb{C})$ is the one inner datum the
construction does reach (Door~4b/4d, via the Hopf tower and Hurwitz's
classification); the residue it cannot reach is precisely the
multiplicity- and real-structure-valued data. Source:\
\texttt{debug/seam\_packing\_scoping\_memo.md}.

\paragraph{The Forced-Count Theorem (synthesis, 2026-06-03).}
The eight preceding paragraphs (Sprint H1 through Door 4f and the
Direction 2 scoping) compose into a single Bertrand-shaped statement
that crystallizes the structural content already proved across the
sub-sprints.  Stated as a theorem it makes explicit what the
distributed prose shows in pieces:\ the framework forces the
\emph{dimension} of the $D_F$ moduli space, but does not (and at the
current packing primitive cannot) force the points within it.

\begin{theorem}[Forced $D_F$ moduli dimension]
\label{thm:forced_count}
Let $\mathcal{T} = (\mathcal{A}_{\mathrm{GV}} \otimes \mathcal{A}_F,\;
\mathcal{H}_{\mathrm{GV}} \otimes \mathcal{H}_F,\;
D_{\mathrm{GV}} \otimes 1 + 1 \otimes D_F,\; J,\; \gamma)$ be the
GeoVac almost-commutative spectral triple of
Sec.~\ref{sec:construction}, assuming:
\begin{enumerate}
\item The Bertrand-with-complex-Hopf-tower truncation $n \le 3$ of
  Sec.~\ref{sec:unified_gauge}, with the Upgrade B sphere-Lie-group
  axiom of Door~4e, forcing $\mathcal{A}_F = \mathbb{C} \oplus
  \mathbb{H} \oplus M_3(\mathbb{C})$.
\item The canonical Chamseddine--Connes--Marcolli SM representation of
  $\mathcal{A}_F$ on $\mathcal{H}_F$, with generation multiplicity
  $N_{\mathrm{gen}}$ imported as an external input (Direction 2:\ not
  packing-derived).
\item The standard real-spectral-triple axioms:\ Hermiticity,
  KO-dim-6 chirality grading $\{\gamma_F, D_F\} = 0$,
  $J$-reality $[J_F, D_F] = 0$, and Connes order-one.
\end{enumerate}
Then the space of admissible finite Dirac operators $D_F$ has real
dimension forced by these inputs alone, with the reduction chain at
$N_{\mathrm{gen}} = 1$:
\begin{equation}
\label{eq:forced_count_chain}
\dim_{\mathbb{R}}\, \mathcal{M}(D_F)
:\; 1024 \xrightarrow{\textnormal{Herm.}} 512
     \xrightarrow{\{\gamma_F, D_F\}=0} 128
     \xrightarrow{[J_F, D_F]=0\,+\,\text{order-one}} 128 .
\end{equation}
The physical-slice dimension under restriction to diagonal Yukawa
$Y = \mathrm{diag}(y_\nu, y_e, y_u, y_d)$ is $8$ per generation.  For
$N_{\mathrm{gen}}$ generations with full generation-mixing Yukawa
matrices the count is $128 \cdot N_{\mathrm{gen}}^2$.
\end{theorem}

\begin{proof}[Proof sketch]
The chain in eq.~\eqref{eq:forced_count_chain} is the H1 sprint
accounting at $n_{\max} = 2$ on the canonical
$\mathcal{A}_F$-representation ($32 \times 32$ at
$N_{\mathrm{gen}} = 1$):\ Hermiticity halves the free parameters;
$\{\gamma_F, D_F\} = 0$ kills the chirality-diagonal blocks;
$J$-reality identifies the matter and antimatter sectors; the
order-one condition is satisfied for the SM representation without
further reduction (the lepton and quark off-chirality blocks each
satisfy order-one because $\mathcal{A}_F$ acts block-diagonally in
lepton/quark and $M_3(\mathbb{C})$ commutes with the Yukawa flavour
structure).  Each step is verified bit-exactly at
$n_{\max} \in \{2, 3\}$ in module
\texttt{geovac/standard\_model\_triple.py} (45 tests passing).  The
$N_{\mathrm{gen}}^2$ scaling follows from the off-chirality block
dimension scaling linearly with $N_{\mathrm{gen}}$ in each chirality.
The inputs (1)--(3) are themselves theorems / proven structural
forcings:\ (1) by Door~4b + Door~4d/4e/4f (the Hopf-tower truncation
forces the factor count; the sphere-Lie-group axiom forces ${\mathbb
  H}$ at $n = 2$); (3) is the standard real-spectral-triple axiom
set; (2) is the one external input, recorded by the Direction 2
paragraph as packing-unreachable in the current Paper 0 construction.
\end{proof}

\begin{corollary}[Bertrand analogy]
\label{cor:bertrand_analogy}
Theorem~\ref{thm:forced_count} is the inner-factor analogue of
Bertrand's theorem on closed-orbit central forces.  Bertrand forces
the \emph{form} of the central potential ($-k/r$ or $k r^2$) without
forcing the coupling $k$; Theorem~\ref{thm:forced_count} forces the
\emph{dimension} of the $D_F$ moduli space without forcing the
$D_F$-entries themselves.  The two seams are structurally identical:\
the framework determines what the moduli space \emph{is}, not where
in it the realised system sits.
\end{corollary}

\begin{remark}[Paired with the Door 4 negative half]
\label{rem:forced_count_with_seam}
The positive content of Theorem~\ref{thm:forced_count} sits beside two
negative companions already recorded above:
\begin{itemize}\setlength{\itemsep}{2pt}
\item \textbf{Door 4 seam theorem (Yukawa-value wall):}\ the
  $\eta$-trivialization plus tensor factorization plus the G3
  commuting-$\mathbb{Z}_2$ result prove that the $D_F$-entries inside
  $\mathcal{M}(D_F)$ are in a Dirichlet ring disjoint from every
  forced outer ring, so the framework cannot derive them.  Sprint
  Yukawa-PSLQ (2026-06-03) added empirical confirmation:\ a
  $162$-cell sweep of the nine measured Yukawa values against the
  forced inner-factor pure-Tate basis $\mathbb{Q}[\pi, \pi^{-1}]$
  (basis sharpened by the $\eta$-trivialization audit ruling out
  $M_3$/Dirichlet $L$ on the inner factor) returned zero hits at
  $M \le 10^3$, at both MS-bar $M_Z$ and MS-bar
  $\mu = 2 \times 10^{16}$ GeV, across three transforms
  ($y_f,\, y_f^2,\, \log y_f$).  Source:\
  \texttt{debug/sprint\_yukawa\_pslq\_memo.md}.
\item \textbf{Direction 2 packing-reach theorem:}\
  $N_{\mathrm{gen}}$ and the inner KO-dimension are
  packing-unreachable in the current Paper 0 construction
  (preceding paragraph, source
  \texttt{debug/seam\_packing\_scoping\_memo.md}).
\end{itemize}
Theorem~\ref{thm:forced_count} thus crystallises the \emph{full}
forced/free seam at the level of $D_F$:\ the forced data is the
moduli space and its dimension; the free data is (i)~the point in
the moduli space (Yukawa values, Door 4 wall), (ii)~the
multiplicity factor ($N_{\mathrm{gen}}$, Direction 2), and (iii)~the
real-structure signature (inner KO-dimension, Direction 2).  All three
free data are calibration inputs in the structural-skeleton-scope
sense (\S 1.5 of the project guidelines).
\end{remark}

\paragraph{Files.}  Scoping memo
\texttt{debug/g4\_cross\_manifold\_scoping\_memo.md}
($\sim 5\,100$ words).  G4a sprint memo
\texttt{debug/sprint\_g4a\_connes\_sm\_memo.md}.
H1 sprint memo
\texttt{debug/h1\_higgs\_inner\_fluctuation\_memo.md}.
G4b is recorded as an open structural
question.

% -------------------------------------------------------------------
\subsection{Frontier-of-field framing: GeoVac-internal walls vs.
  spectral-action open problems}
\label{sec:frontier_framing}
% -------------------------------------------------------------------

The five walls catalogued in the multi-focal-composition audit
(\texttt{debug/multifocal\_phase\_a\_synthesis\_memo.md}, May~2026)
are not all GeoVac-specific.  Reading them through the spectral-triple
construction of Sec.~\ref{sec:construction} and the Marcolli--van~Suijlekom
lineage of Sec.~\ref{sec:lineage}, three are GeoVac-internal architectural
gaps that are tooling-addressable, and two sit at the broader frontier
of the spectral-action / NCG programs.  We name them here for honesty
and to fix the scope of the open questions of Sec.~\ref{sec:open}.

\paragraph{GeoVac-internal walls (W1a, W1b, W1c).}

\begin{itemize}\setlength{\itemsep}{2pt}
\item \textbf{W1a (cross-register coordinate operator):}\ no two-body
spatial operator $V(\hat{\mathbf{r}}_e, \hat{\mathbf{R}}_n)$ has been
implemented across the nuclear and electronic registers of the
composed nuclear--electronic Hamiltonian (Paper~23~\cite{paper23}~\S
VI; \texttt{geovac/nuclear/nuclear\_electronic.py}).  The framework
couples the registers through spin-spin hyperfine and a classical
finite-size shift but evaluates spatial operators at the classical
proton position.  This is an architectural gap, not a no-go theorem;
the calibrated atomic-physics target is the
Pachucki--Patko\v{s}--Yerokhin recoil Hamiltonian (PRL 130, 023004,
2023), which is exact in the mass ratio at $(Z\alpha)^6$.
\item \textbf{W1b (magnetization-distribution operator):}\ no
operator on the proton register represents the Zemach radius
$r_Z$ as a Layer-2 magnetization focal length distinct from the
charge radius $r_p$ (Paper~34~\cite{paper34}~\S\,III; Sprint HF-4,
May~2026).  Calibrated against Eides~et~al.\ \emph{Phys.\ Lett.\ B}
(2024).  Architectural gap; QCD-internal magnetization density is
input data, not GeoVac-internal.
\item \textbf{W1c (cross-center frozen-core screening):}\ the
\texttt{FrozenCore} $Z_{\mathrm{eff}}(r)$ provides same-center
screening only; cross-center screening of bare $V_{ne}$ from the
opposite nucleus is not implemented (Sprint~7 NaH/MgH$_2$ balanced
FCI; CLAUDE.md~\S2).  Sprint~7b made partial progress for the
spin-orbit splitting; the PES extension is mechanical once cross-center
screening is wired through the multipole expansion.  Architectural
gap.
\end{itemize}

These three walls are tooling-addressable inside the present
framework: the missing operators exist in principle, are calibrated
against published atomic-physics targets, and require no new
mathematics in the broader NCG literature.  Each can be closed by a
sprint at the 2--10 week scale.  Their existence does not undermine
the spectral-triple framing of Sec.~\ref{sec:construction}; they
catalogue what has not yet been built rather than what cannot be.

\paragraph{Frontier-of-field walls (W2a, W2b).}

\begin{itemize}\setlength{\itemsep}{2pt}
\item \textbf{W2a (multi-loop UV/IR composition):}\ the bare
iterated Connes--Chamseddine spectral action on the
Camporesi--Higuchi spectrum reproduces UV-divergent integrands of
multi-loop QED (Sprint LS-8a, May~2026; Paper~36~\cite{paper36}~\S
VII) but cannot autonomously generate $Z_2$, $\delta m$, $Z_3$
counterterms for finite extraction at two loops.  This is not
GeoVac-specific.  The Marcolli--van~Suijlekom 2014 rationality theorem
for Robertson--Walker spectral actions~\cite{marcolli_vs2014}
partially addresses spectral-action coefficient structure on
homogeneous compact backgrounds, and the subsequent Hekkelman--McDonald
2024 noncommutative integral approximation~\cite{hekkelman_mcdonald2024}
provides a single-cutoff Szeg\H{o}-style asymptotic theory in the
Connes--van~Suijlekom paradigm; neither closes the multi-cutoff
matching question that renormalization counterterm generation
requires.  The published QFT archetype is the matched effective-field-theory
chain (e.g.\ HQET~$\to$~NRQED~$\to$~pNRQED~$\to$~SCET, JHEP 05 (2025)
171), where each matching step is at a single physical scale and
imports $\overline{MS}$ counterterms by hand.  No published spectral-action
framework reproduces this autonomously.  W2a is therefore an open
question of the spectral-action program, of which GeoVac is one
specific instance; it is not a defect of the present construction
relative to the lineage.

\item \textbf{W2b (cross-manifold spectral-triple composition):}\ the
tensor product of two infinite metric spectral triples
$(\mathcal{T}_{S^3}, d_{D_{S^3}}) \otimes (\mathcal{T}_{S^5}, d_{D_{S^5}})$,
needed to compose the GeoVac Coulomb sector (Sec.~\ref{sec:construction})
with the Bargmann--Segal HO sector of Paper~24~\cite{paper24}, has no
published Latr\'emoli\`ere-propinquity convergence theorem.  The
almost-commutative case (one infinite Riemannian factor + one finite
factor) is closed by Latr\'emoli\`ere 2026~\cite{latremoliere2026}
(\emph{Spectral continuity of almost commutative manifolds for the
$C^1$ topology on Riemannian metrics}); the two-infinite case is
explicitly not handled in that paper, nor in the closely related
Farsi--Latr\'emoli\`ere 2024 inductive-sequence work (Adv.\ Math.\ 437,
109442).  This is the structural content of the G4b cross-manifold
blocker (Sec.~\ref{sec:g4}):\ not a GeoVac-specific limitation but a
genuinely open NCG question.  Closing it would itself be a substantial
mathematical contribution, parallel to but separate from Paper~38's
SU(2) propinquity theorem.
\end{itemize}

\paragraph{Forecast (without claiming closure).}

We do not claim W1a/W1c can be closed without surprises;\ the
multi-$\lambda$ Shibuya--Wulfman algebraic check that gates the W1a
attack (multipole termination across mismatched exponents)
is not in the published Avery-school literature and could fail to
preserve Gaunt-style closure (Track~3 literature memo, May~2026,
\S 1).  W1b is bounded by external QCD input.  W2a and W2b are open
in the broader literature, not just in the GeoVac record.  The honest
forecast is that closing W1a/b/c is internal-effort plausible at the
multi-month scale per wall;\ closing W2a or W2b would be a
contribution to the spectral-action / NCG programs rather than to
GeoVac alone, and we do not commit to it as a deliverable here.

\paragraph{Implication for the Marcolli--vS lineage placement.}

Sec.~\ref{sec:lineage} read the present construction as a graph
spectral triple in the lineage of \cite{marcolli_vs2014, perez_sanchez2024,
perez_sanchez2025}.  The frontier-of-field framing above sharpens this:\
the lineage placement is correct, and the W2a/W2b walls are the lineage's
own open problems, not GeoVac's separate concerns.  This is consistent
with Observation~\ref{obs:lineage} (each structural ingredient has
published precedent;\ what is not duplicated is the $\alpha$
prediction at $8.8 \times 10^{-8}$).  The W1 walls remain GeoVac-side
because they are about the architectural extension of the present
triple to multi-particle real-space coupling, which the SM-style almost-commutative
construction does not address.

% -------------------------------------------------------------------
\subsection{Lorentzian transfer audit and the master Mellin engine
under signature extension (Sprint L0, May 2026)}
\label{sec:lorentzian_audit_paper32}
% -------------------------------------------------------------------

The case-exhaustion theorem
(Theorem~\ref{thm:pi_source_case_exhaustion}) classifies every $\pi$
in a finite chain of Paper~34 projections as engaging one of three
sub-mechanisms of a master Mellin engine
$\mathcal{M}[\mathrm{Tr}(D^k\,e^{-tD^2})]$ with
$k \in \{0, 1, 2\}$: M1 (Hopf-base measure), M2 (Seeley--DeWitt
heat-kernel coefficients), M3 (vertex-parity Hurwitz / Dirichlet-$L$
content).  The GH-convergence theorem
(Theorem~\ref{thm:gh_convergence}) and the Bisognano--Wichmann
reading (Remark~\ref{rem:bisognano_wichmann_reading}, with Unruh
pendant) close WH1 on the Riemannian side.  We close out
Sec.~\ref{sec:pi_source_theorem} with a brief structural appendix
that records what Sprint L0 (May 2026) found when the audit was
extended to the Riemannian $(3, 0) \to$ Lorentzian $(3, 1)$
extension of the framework.

\paragraph{Sprint L0 verdict on the master Mellin engine.}
Sprint L0 (\texttt{debug/lorentzian\_l0\_audit\_memo.md}, May 2026;
machine-readable table at
\texttt{debug/data/lorentzian\_l0\_audit.json}) classifies all
twenty-eight entries of Paper~34's projection
dictionary~\cite{paper34} by their transfer behavior under signature
extension; the full per-projection table sits at
Paper~34~\S~V.E (\texttt{sec:lorentzian\_transfer\_audit}), the
structural framing of the Sig/Op partition is in Paper~31~\S~8
(\texttt{sec:signature\_partition})~\cite{paper31}.  Two
sharpenings of the master Mellin engine fall out of the L0 audit.

\medskip
\noindent\emph{(i) M1 sub-mechanism is the load-bearing free-transfer
engine.}  The M1 sub-mechanism ($k = 0$) on a temporal $S^1_\beta$ or
proper-time Rindler $\eta_E$-circle inherits a Lorentzian reading via
the four-witness Wick-rotation theorem (Hawking + Sewell + BW + Unruh,
codified Sprint Unruh-pendant 2026-05-10;
Remark~\ref{rem:bisognano_wichmann_reading} and its Unruh-pendant
extension).  Four entries of Paper~34's dictionary sit in the
\textsc{wick-map free} bucket via this mechanism:\ observation /
temporal-window (\S~III.15), vector-photon promotion (\S~III.11) at
the $1/(4\pi)$ per-loop level, gauge choice (\S~III.26) at the
Coulomb-gauge per-loop level, and the Wick-rotation projection itself
(\S~III.27).  None requires Krein-space machinery for its
metric-functional-level reading.  Equivalently:\ the
$2\pi = \mathrm{Vol}(S^1)$ signature of M1 in the master Mellin engine
is also the modular-flow / Lorentz-boost orbit period via published
QFT, with the same generator on both sides of the Wick rotation.

\medskip
\noindent\emph{(ii) M3 sub-mechanism predicted to trivialize at
$(3, 1)$.}  A non-trivial structural prediction emerges from the
Bizi--Brouder--Besnard~\cite{bizi_brouder_besnard2018} signature
classification:\ at $(m, n) = (3, 1)$, the BBB sign table flips the
$\{J, \gamma_5\}$ anticommutation relation that drives M3's
vertex-parity content on the Riemannian side (Paper~28's
$\zeta(s, 3/4)$ and $\zeta(s, 5/4)$ Hurwitz reductions producing
Catalan's $G$ and $\beta(4)$).  On chirality-symmetric spectrum at
$(3, 1)$, the vertex-parity asymmetry vanishes, M3's
half-integer-Hurwitz content becomes structurally empty, and the
case-exhaustion theorem
(Theorem~\ref{thm:pi_source_case_exhaustion}) trivializes to a
\emph{two-mechanism partition} M1 $\cup$ M2 on the Lorentzian side.
This is not a contradiction with the Riemannian theorem:\ the theorem
is stated at signature $(3, 0)$ and the case-exhaustion is over
mechanisms producing $\pi$ in finite chains \emph{at that signature}.
The $(3, 1)$ trivialization is a structural sharpening that becomes
a Sprint L2 consistency check (3--6 months, gated on blocker B2; see
Paper~31~\S~8.5).

\paragraph{What this does and does not change.}
The Riemannian case-exhaustion theorem stands;\ WH1 PROVEN
(Theorem~\ref{thm:gh_convergence}) stands.  The L0 audit is a
documentation-level extension that classifies the framework's transfer
behavior under signature change without re-deriving any of the
Riemannian results.  Sprint L1 (4--8 weeks, MEDIUM-HIGH) would
deepen the Wick-rotation reading from structural correspondence to
literal operator-system identification by computing the modular
Hamiltonian $K$ on $\mathcal{T}_{n_{\max}}$ for a Rindler-wedge-restricted
Camporesi--Higuchi state, verifying $\sigma_{i \cdot 2\pi} = $ identity
at finite $n_{\max}$, and taking the GH limit using
Theorem~\ref{thm:gh_convergence}'s lemmas.  That closure would lift the
correspondence on all four Wick-rotation faces simultaneously via the
shared KMS-Tomita-Takesaki algebra.  The L0 audit names this as the
principal Sprint L1 falsifier without itself closing it.

The Nieuviarts~2025~\cite{nieuviarts2025a} twisted-spectral-triple
morphism is the named trigger that motivated reopening the May 2026
NO-GO scoping;\ a parallel L2-Nieuviarts-scoping pass tests whether
the construction applies cleanly to $S^3 = \mathrm{SU}(2)$
(odd-dimension caveat) and would shorten blocker B1 (Lorentzian
propinquity, 6--12 months) substantially if applicable.  The current
audit does \emph{not} assume the morphism applies.

% -------------------------------------------------------------------
\subsection{Krein space construction at signature $(3, 1)$
(Sprint L2-B, 2026-05-16)}
\label{sec:l2b_krein}
% -------------------------------------------------------------------

The first concrete deliverable of the Sprint L2 arc lifts the Hilbert
space and the fundamental symmetry of the truncated triple from
signature $(3, 0)$ to signature $(3, 1)$, following the van~den~Dungen
2016 prescription~\cite{vandungen2016}.  The construction is at the
\emph{Hilbert-space level only}; the Lorentzian Dirac and the Connes
axiom audit are deferred to L2-C and L2-D below.

\paragraph{Krein space.}
For $n_{\max} \ge 1$ and $N_t \ge 1$ define
\begin{equation}
\label{eq:l2b_krein_space}
   \mathcal{K}_{n_{\max}, N_t}
   \;:=\; \mathcal{H}_{\mathrm{GV}}^{n_{\max}}
     \otimes L^2(\mathbb{R}_t)_{\mathrm{cutoff}, N_t},
\end{equation}
where $\mathcal{H}_{\mathrm{GV}}^{n_{\max}}$ is the Riemannian
Hilbert space of Definition~\ref{def:H_GV} and $L^2(\mathbb{R}_t)_{\mathrm{cutoff}, N_t}$
is the truncated temporal slot: a uniform real grid
$t_k = -T_{\max} + k \cdot 2 T_{\max}/(N_t - 1)$ for $N_t \ge 2$, or the
singleton $\{0\}$ for $N_t = 1$.  The temporal slot is NOT compactified
to $S^1_\beta$ (which would create closed timelike curves per Geroch's
theorem and is structurally incompatible with the bounded-wedge
reading of Paper~42).

\paragraph{Conventions.}
We use the West-coast metric $\eta = \mathrm{diag}(+1, -1, -1, -1)$
and the Peskin--Schroeder chiral (Weyl) basis for the four Cl$(3, 1)$
gamma matrices.  In this basis, $\gamma^5 = \mathrm{diag}(-I_2, +I_2)$
is diagonal and \verb|FullDiracLabel.chirality| is identified with
the $\gamma^5$ eigenvalue, so the Riemannian limit
identification with $\mathcal{H}_{\mathrm{GV}}$ is a transparent
basis-label match (no similarity transform required).  The Dirac
basis alternative (where $\gamma^0$ is diagonal instead) is
unitarily equivalent at the operator level but would have required
a unitary similarity transform for the Riemannian-limit check;
the chiral basis was chosen on Riemannian-limit-clarity grounds.

\paragraph{Fundamental symmetry $J$.}
The fundamental symmetry is
\begin{equation}
\label{eq:l2b_J_def}
   J \;=\; J_{\mathrm{spatial}} \otimes I_{N_t},
\end{equation}
where $J_{\mathrm{spatial}}$ is the lift of $\gamma^0$ to the
chirality-doubled $\mathcal{H}_{\mathrm{GV}}$: on
$|n_D, \kappa, m_j, \chi\rangle$ the action is the chirality swap
$|n_D, \kappa, m_j, -\chi\rangle$.  In the chiral basis,
$\gamma^0 = \bigl(\begin{smallmatrix} 0 & I_2 \\ I_2 & 0 \end{smallmatrix}\bigr)$
is the off-diagonal block-swap, so $J_{\mathrm{spatial}}$ is a real
$\{0, 1\}$-valued permutation matrix.  The Krein inner product is
$\langle \psi, \phi \rangle_{\mathcal{K}} :=
\langle \psi, J \phi \rangle$.

\paragraph{Krein axiom verification.}
At every panel cell $(n_{\max}, N_t) \in \{1, 2, 3\} \times \{1, 11,
21\}$, the construction satisfies
\begin{enumerate}\setlength{\itemsep}{2pt}
\item $J^2 = +I$ (bit-exact Frobenius residual $= 0.0$),
\item $J^* J = I$ ($J$ Hilbert-unitary, bit-exact),
\item $J^* = J$ ($J$ Hilbert-Hermitian, bit-exact),
\item $\mathcal{K} = \mathcal{K}^+ \oplus \mathcal{K}^-$ with
$\dim \mathcal{K}^\pm = \dim \mathcal{K}/2$ and $J|_{\mathcal{K}^\pm}
= \pm I$ on the eigenspaces (rank $\dim \mathcal{K}/2$ each, bit-exact).
\end{enumerate}
The bit-exact-zero residuals follow from the structural fact that
$J = J_{\mathrm{spatial}} \otimes I_{N_t}$ is the Kronecker product of
a $\{0, 1\}$-valued integer permutation matrix and an identity matrix,
both involving only operations exact in IEEE~754 arithmetic on real
integers.

\paragraph{Load-bearing Riemannian-limit check.}
At $N_t = 1$ the temporal slot collapses to the singleton $\{0\}$ and
$\mathcal{K}_{n_{\max}, N_t = 1}$ recovers $\mathcal{H}_{\mathrm{GV}}^{n_{\max}}$
\emph{bit-identically} at every $n_{\max} \in \{1, 2, 3\}$:\ the
dimension formula $\dim \mathcal{K}_{n_{\max}, 1} = N_{\mathrm{Dirac}}(n_{\max})
= \tfrac{2}{3} n_{\max} (n_{\max} + 1)(n_{\max} + 2) = 4, 16, 40$ matches
the Paper~32~\S~III formula exactly, the basis labels are bit-identical
under \verb|FullDiracLabel.__eq__|, and the $J$ matrix at $N_t = 1$
equals $J_{\mathrm{spatial}}$ as numpy arrays
(Frobenius difference $= 0.0$).  This confirms that the
Camporesi--Higuchi spatial spinor bundle is structurally compatible
with the Cl$(3, 1)$ gamma-matrix embedding in the West-coast chiral
basis;\ WH1 PROVEN
(Theorem~\ref{thm:gh_convergence}) is not re-opened by this result.
Source:\ \verb|debug/l2_b_krein_construction_memo.md|, raw data
\verb|debug/data/l2_b_krein_construction.json|, implementation
\verb|geovac/krein_space_construction.py|, regression tests
\verb|tests/test_krein_space_construction.py| (109 cases, all pass;
zero regression in 132 upstream tests).

\paragraph{Honest scope.}
The Krein-space construction lifts the Hilbert space and the
fundamental symmetry only.  The Lorentzian Dirac, the algebra
$\mathcal{A}_L$ on $\mathcal{K}$, the operator-algebra Krein-adjoint
structure, and the modular Hamiltonian on the Krein space are deferred
to L2-C, L2-D, and L2-E respectively.

% -------------------------------------------------------------------
\subsection{Lorentzian Dirac operator $D_L$ on $S^3 \times \mathbb{R}$
(Sprint L2-C, 2026-05-16)}
\label{sec:l2c_lorentzian_dirac}
% -------------------------------------------------------------------

Sprint L2-C constructs the Lorentzian Dirac operator on the Krein
space~\eqref{eq:l2b_krein_space} via the van~den~Dungen Proposition~4.1
prescription~\cite{vandungen2016}:\ Wick-rotate the Lorentzian metric
$g$ on $M = S^3 \times \mathbb{R}$ to its Riemannian counterpart
$g_r$, construct the standard Riemannian Dirac
$\not{\!\!D}_{g_r}$ on $(M, g_r)$, and define the Krein-self-adjoint
Lorentzian Dirac as $D_L = i^t \not{\!\!D}_{g_r}$ with $t = 1$ for
signature $(s, t) = (3, 1)$.

\paragraph{Construction.}
The Wick-rotated Riemannian Dirac decomposes via spatial-temporal
splitting:
\begin{equation}
\label{eq:l2c_dirac_decomposition}
   \not{\!\!D}_{g_r}
   \;=\; \gamma^0 \otimes \partial_t
     + D_{\mathrm{GV}} \otimes I_{N_t},
\end{equation}
with $\gamma^0$ from the Sprint L2-B chiral basis,
$D_{\mathrm{GV}}$ the Paper~32~\S~III Camporesi--Higuchi truthful
full-Dirac matrix on $\mathcal{H}_{\mathrm{GV}}^{n_{\max}}$, and
$\partial_t$ the temporal-derivative matrix on the $L^2(\mathbb{R}_t)_{\mathrm{cutoff}, N_t}$
grid.  The Lorentzian Dirac is then
\begin{equation}
\label{eq:l2c_lorentzian_dirac}
   D_L \;=\; i \cdot \bigl[ \gamma^0 \otimes \partial_t
     + D_{\mathrm{GV}} \otimes I_{N_t} \bigr]
   \;\in\; \mathbb{C}^{\dim \mathcal{K} \times \dim \mathcal{K}}.
\end{equation}

\paragraph{Temporal derivative.}
The temporal-derivative $\partial_t$ uses centered finite differences
with Dirichlet (zero) boundary conditions:
\begin{equation}
\label{eq:l2c_partial_t}
   (\partial_t)_{k, l}
   \;=\; \frac{1}{2 \Delta_t} \bigl(\delta_{l, k+1} - \delta_{l, k-1}\bigr),
   \qquad \Delta_t = \frac{2 T_{\max}}{N_t - 1}.
\end{equation}
This choice is \emph{structurally forced}, not a free parameter:\ the
centered-FD plus Dirichlet-zero combination is the unique discretization
on a finite grid that gives anti-Hermitian $\partial_t$
($\partial_t^\dagger = -\partial_t$), which is the property required
for Krein-self-adjointness $D_L^\times = D_L$ to close via
the vdD recipe.  Forward, backward, and upwind schemes give
$\partial_t^\dagger \neq -\partial_t$ (breaking Krein-self-adjointness);
periodic boundary conditions would re-introduce closed timelike curves
(Sprint L2-A audit~\S~3.8).

\paragraph{Sign of $i^t$.}
With $t = 1$, the prescription gives $i^t = i$ (not $-i$).  The sign
is derived from the structural requirement $D_L^\times = D_L$ with
$D_L^\times := \gamma^0 D_L^\dagger \gamma^0$, using
$\{\gamma^0, D_{\mathrm{GV}}\} = 0$ (the chirality swap $\gamma^0$
anticommutes with the chirality-diagonal $D_{\mathrm{GV}}$) and
$\partial_t^\dagger = -\partial_t$.  Explicit symbolic calculation
(sympy verification) gives $D_L^\times = +D_L$ for $i^t = +i$ and
$D_L^\times = -D_L$ for $i^t = -i$.

\paragraph{Verification at finite cutoff.}
The construction verifies at every panel cell $(n_{\max}, N_t) \in
\{1, 2, 3\} \times \{1, 11, 21\}$:
\begin{enumerate}\setlength{\itemsep}{2pt}
\item \emph{Krein-self-adjointness.}  $\|D_L^\times - D_L\|_F = 0$
bit-exact at every cell.
\item \emph{Real spectrum on Krein-positive cone.}  Restricting $D_L$
to $\mathcal{K}^+$ (the $+1$ eigenspace of $J$) gives real eigenvalues
to machine precision ($\max |\mathrm{Im}(\lambda)| \le 2.32 \times
10^{-17}$).
\item \emph{Riemannian-limit recovery (LOAD-BEARING).}  At $N_t = 1$,
$D_L|_{N_t = 1} = i \cdot D_{\mathrm{GV}}$ bit-identically.  The
absolute-value spectrum $\{|\lambda_n| = n_{\mathrm{Fock}} + 1/2\}$
with multiplicities $g_n^{\mathrm{Dirac}} = 2 n (n+1)$ matches
$D_{\mathrm{GV}}$ bit-exactly at $n_{\max} = 1, 2, 3$.
\end{enumerate}

\paragraph{Chirality finding (non-zero $\{\gamma^5, D_L\}$).}
A non-trivial structural finding emerges:\ the BBB universal
relation $\chi D = -D\chi$ (Sec~5(v) of~\cite{bizi_brouder_besnard2018})
does \emph{not} hold on truthful $D_{\mathrm{GV}}$.  The
anticommutator $\{\gamma^5, D_L\}$ has Frobenius norm
$2 \|\gamma^5 D_{\mathrm{GV}}\|_F = 2 \|D_{\mathrm{GV}}\|_F$ at
$N_t = 1$ (values $6.00, 18.33, 38.88$ at $n_{\max} = 1, 2, 3$),
because GeoVac's truthful $D_{\mathrm{GV}}$ is chirality-diagonal in
the Peskin--Schroeder chiral basis and therefore commutes with
$\gamma^5$.  This finding is structural, not a bug:\ it is a
consequence of two locked Riemannian-side conventions (chirality-as-$\gamma^5$
from L2-B and chirality-diagonal $D_{\mathrm{GV}}$ from Paper~32~\S~III).
The full diagnostic and the three resolutions are documented in
Sec.~\ref{sec:l2d_axiom_audit_31} below.

\paragraph{Sources.}
Implementation:\ \verb|geovac/lorentzian_dirac.py|.  Regression tests:\
\verb|tests/test_lorentzian_dirac.py| (108 cases, all pass;
zero regression in 205 upstream tests).  Sprint memo:\
\verb|debug/l2_c_lorentzian_dirac_memo.md|; raw data
\verb|debug/data/l2_c_lorentzian_dirac.json|.

% -------------------------------------------------------------------
\subsection{Connes axiom audit at $(m, n) = (4, 6)$
(Sprint L2-D, 2026-05-16)}
\label{sec:l2d_axiom_audit_31}
% -------------------------------------------------------------------

Sprint L2-D performs the Connes axiom audit on the pair $(D_L,
\gamma^5, J_L)$ built in Sprints L2-B/C, against the
Bizi--Brouder--Besnard~2018~\cite{bizi_brouder_besnard2018}
classification at $(m, n) = (4, 6)$.

\paragraph{BBB sign determination.}
Reading the BBB Table~1 entries at $(m, n) = (4, 6)$ directly from
arXiv:1611.07062~v2:
$\varepsilon = +1$, $\varepsilon'' = -1$, $\kappa = -1$, $\kappa'' = +1$.
These four primitive signs, combined with the four BBB defining
relations $J^2 = \varepsilon$, $J\chi = \varepsilon''\chi J$,
$J\eta = \varepsilon\kappa\,\eta J$, $\eta\chi = \varepsilon''\kappa''\,\chi\eta$,
predict at $(4, 6)$:\ $J^2 = +I$, $\{J, \chi\} = 0$,
$\{J, \eta\} = 0$, $\{\eta, \chi\} = 0$.  At the bare Cl$(3, 1)$
4-spinor level with charge conjugation $U_4 = i\gamma^2$, all four
relations hold bit-exact in the chiral basis (Frobenius residual
$= 0.0$).  The translation $(s, t) = (3, 1) \leftrightarrow (m, n) =
(4, 6)$ uses BBB Table~3:\ $m \equiv t + s$, $n \equiv t - s \pmod 8$,
giving $m = 4$, $n = -2 \equiv 6 \pmod 8$.

\paragraph{Lift to the Krein space.}
The 4-spinor charge conjugation $i\gamma^2 = (i\sigma_y)_{\mathrm{chir}}
\otimes (i\sigma^2)_{\mathrm{spin}}$ decomposes into a chirality-swap-with-sign
factor times the standard spin-1/2 charge conjugation.  The lift to
$\mathcal{K}_{n_{\max}, N_t}$ is
\begin{equation}
\label{eq:l2d_J_lorentzian}
   J_L \;=\; J_{L,\mathrm{spatial}} \otimes K_t,
\end{equation}
where $J_{L,\mathrm{spatial}}$ acts on $|n, l, m_j, \chi\rangle$ as
$\sigma_{\mathrm{chir}}(\chi) \cdot \sigma_{\mathrm{spin}}(l, -m_j)
\cdot |n, l, -m_j, -\chi\rangle$ with
$\sigma_{\mathrm{chir}}(\pm 1) = \mp 1$ and
$\sigma_{\mathrm{spin}}(l, m_j) = (-1)^{(2l + 1 - 2m_j)/2}$, and
$K_t$ is complex conjugation on $\mathbb{C}^{N_t}$.  The matrix
$U_L := J_L \cdot K^{-1}$ is a real $\{0, \pm 1\}$-valued unitary;
the antilinear $J_L = U_L K$ satisfies $J_L^2 = +I$ bit-exact via
$U_L \overline{U_L} = U_L^2$ in IEEE~754.

\paragraph{Six-axiom panel.}
At every panel cell $(n_{\max}, N_t) \in \{1, 2, 3\} \times \{1, 11,
21\}$ on truthful $D_{\mathrm{GV}}$:

\begin{center}
\begin{tabular}{lcc}
\toprule
Axiom & Predicted at $(4, 6)$ & Residual \\
\midrule
(i) $J_L^2 = +I$ (BBB $\varepsilon = +1$) & $+I$ & $0.0$ bit-exact \\
(ii) $\{J_L, \chi\} = 0$ (BBB $\varepsilon'' = -1$) & $0$ & $0.0$ bit-exact \\
(iii) $\{J_L, \eta\} = 0$ (BBB $\varepsilon\kappa = -1$) & $0$ & $0.0$ bit-exact \\
(iv) $J_L D_L = +D_L J_L$ (BBB universal Sec~5(v)) & $JD = +DJ$ & $0.0$ bit-exact \\
(v) $\{\chi, D_L\} = 0$ (BBB universal Sec~5(v)) & $\chi D = -D\chi$
  & STRUCTURAL FINDING (non-zero) \\
(vi) order-zero $[a, J_L b J_L^{-1}] = 0$ & $0$
  & $\le 0.0675$ finite-resolution (sample-of-3) \\
(vii) order-one $[[D_L, a], J_L b J_L^{-1}] = 0$ & $0$
  & $\le 0.101$ finite-resolution (sample-of-3) \\
\bottomrule
\end{tabular}
\end{center}

Axioms (i)--(iv) hold bit-exactly at every cell.  Axioms (vi)--(vii)
exhibit the same 5--10\% finite-resolution residuals as the
Riemannian-side audit (Table~\ref{tab:axiom_audit}, ``Order zero
$[a, J b J^{-1}] = 0$''); the residuals are sample-of-3 estimates on
the canonical $\mathcal{A}_{\mathrm{GV}}$ multiplier basis and
match the structural-artifact reading of Connes--van~Suijlekom
operator-system multiplicative-closure failure.  Axiom (v) is the
structural finding documented in
Table~\ref{tab:axiom_audit_lorentzian} and discussed below.

\paragraph{Structural finding on $\{\chi, D_L\}$.}
The BBB universal axiom $\chi D = -D\chi$ fails on truthful
$D_{\mathrm{GV}}$:\ $\|\{\gamma^5_K, D_L\}\|_F = 2\|D_{\mathrm{GV}}\|_F$
at $N_t = 1$, scaling as
$\|D_{\mathrm{GV}}\|_F \sqrt{N_t} \cdot (2 + \mathrm{small})$ at
$N_t > 1$.  The mechanism is the chirality-diagonal structure of
$D_{\mathrm{GV}}$ in the chiral basis (\S~\ref{sec:l2c_lorentzian_dirac}).
This is the load-bearing scope finding of Sprint L2-D and is the same
structural fact captured in Table~\ref{tab:axiom_audit_lorentzian}.
The three resolutions (R1 accept, R2 use offdiag CH, R3 redefine
$\gamma^5$) are documented in the discussion immediately following
Table~\ref{tab:axiom_audit_lorentzian}, with R1 + R2 in parallel as
the recommended path for Sprint L2-E.

\paragraph{M3 trivialization is convention-dependent.}
Sprint L0 (\S~\ref{sec:lorentzian_audit_paper32} above) predicted that
the M3 sub-mechanism of the master Mellin engine would
\emph{trivialize at $(3, 1)$ on chirality-symmetric spectrum}, via the
BBB sign-flip $\{J, \gamma^5\} = 0$ removing the asymmetry that drives
M3's vertex-parity content.  Sprint L2-D tests this prediction and
finds:\ \textbf{the prediction is convention-dependent}.

Under the \emph{$n_{\mathrm{Fock}}$-parity convention} (the Paper~28
\S QED-vertex convention that accesses Catalan $G$ and Dirichlet
$\beta(4)$ via quarter-integer Hurwitz shifts), the vertex-parity sum
$D_{\mathrm{even}}(4) - D_{\mathrm{odd}}(4)$ on the chirality-symmetric
$(3, 1)$ truncation is bit-identical to the Riemannian value.  Each
full-Dirac level $n_{\mathrm{Fock}}$ contributes its full degeneracy
$g_n = 2 n (n+1)$ summed over both chirality blocks;\ the parity is
read from the chirality-independent $n_{\mathrm{Fock}}$ label.
The L0 prediction is therefore \textbf{FALSIFIED} under this reading
at every $n_{\max} \in \{1, \ldots, 5\}$ (bit-identical difference).

Under the \emph{chirality-pairing convention} (pair ``even'' with
$\chi = +1$ Weyl and ``odd'' with $\chi = -1$ anti-Weyl), the
chirality-symmetric truncation gives $D_+ = D_-$ identically by
chirality symmetry, so the difference vanishes bit-exact.  The L0
prediction is \textbf{trivially confirmed} under this reading, but
tautologically (any chirality-symmetric truncation gives the result
by symmetry alone, independent of signature).

The structural refinement is therefore: \textbf{the BBB sign-flip
$\{J, \gamma^5\} = 0$ at $(m, n) = (4, 6)$ is a spectral-triple-axiom
statement, not a vertex-parity-sum statement.}  It does not directly
imply M3 trivialization at any natural parity convention.
M3's mechanism is sectional in $n_{\mathrm{Fock}}$-parity, NOT in
spacetime signature.  Paper~18~\S~III.7 records the corresponding
master-Mellin-engine sharpening.

\paragraph{Sources.}
Implementation:\ \verb|geovac/connes_axiom_audit_31.py| (700 lines
incl.\ docstrings).  Regression tests:\
\verb|tests/test_connes_axiom_audit_31.py| (75 cases, all pass).
Sprint memo:\ \verb|debug/l2_d_connes_axiom_audit_31_memo.md|; raw data
\verb|debug/data/l2_d_connes_axiom_audit_31.json| (BBB-Table-1 signs,
4-spinor verification, full axiom panel, M3 trivialization panel).

\paragraph{Sprint L2-E handoff.}
The four BBB-predicted-sign axioms hold bit-exactly, so the
operator-algebra structure on the Krein space is fully BBB-compatible
at the $J$-relation level.  Sprint L2-E (the Lorentzian-level Paper~42
redo) should construct the Krein-side modular Hamiltonian
$K_L = -\log \Delta_L$ on the Krein-positive cone $\mathcal{K}^+$ and
verify the analog of Paper~42's $\sigma_{2\pi}(O) = O$ at the
operator-system level on the Krein space.  The resolution choice
R1 (truthful $D_{\mathrm{GV}}$, preserves Riemannian-limit recovery)
versus R2 (offdiag CH, satisfies BBB axiom~(v)) is recommended to be
tested in parallel.  WH1 PROVEN
(Theorem~\ref{thm:gh_convergence}) is not re-opened by these results.

% -------------------------------------------------------------------
\subsection{Krein-level unified-strong four-witness theorem at finite
cutoff (Sprint L2-E, 2026-05-17)}
\label{sec:l2e_krein_modular_hamiltonian}
% -------------------------------------------------------------------

Sprint L2-E (\verb|geovac/modular_hamiltonian_lorentzian.py|, $\sim860$
lines including docstrings; tests in
\verb|tests/test_modular_hamiltonian_lorentzian.py|, 74 fast plus 3
slow, all passing) builds the operator-system-level Krein-side modular
Hamiltonian on the hemispheric wedge of $S^3 \times \mathbb{R}$ at
signature $(3, 1)$ and verifies the four LOAD-BEARING falsifiers of
the Sprint L2-F catalogue
(\texttt{debug/sprint\_l2\_falsifiers.md}~\S~3) plus the
Riemannian-limit recovery LOAD-BEARING falsifier plus the headline
$H_{\mathrm{local}}$ verdict and the six-witness collapse, all bit-exact
at every $(n_{\max}, N_t) \in \{1, 2, 3\} \times \{1, 11, 21\}$.
The closure mirrors Paper~42 \S\S~4--8~\cite{paper42} verbatim with
the temporal slot added via the Sprint L2-B Krein space.

\paragraph{Lorentzian wedge, KMS state, and BW choice.}
The hemispheric wedge on the Krein space is
$W_L = P_W^{\mathrm{spatial}} \otimes P_{t \geq 0}$, where
$P_W^{\mathrm{spatial}} = (1/2)(I + R_{\mathrm{polar}})$ is the
Paper~42 hemispheric wedge on $\mathcal{H}_{\mathrm{GV}}$ (the
$m_j$-reflection involution) and $P_{t \geq 0}$ is the
$t \geq 0$ half-line projector on $\mathbb{C}^{N_t}$.  The wedge KMS
state under the BW choice $H_{\mathrm{local}} := K_L^{\alpha, W}/\beta$
at $\beta = 2\pi$ is
\[
\rho_W^L = \frac{e^{-K_L^{\alpha, W}}}{Z},
\]
$\beta$-independent at the algebra-action level, with $K_L^{\alpha, W}$
the wedge-restricted BW-$\alpha$ generator inheriting integer
eigenvalues $\mathrm{two}\_m_j$ from Paper~42 Definition~5.1.

\paragraph{The Krein-level four-witness theorem at finite cutoff.}
The Lorentzian-side BW-$\alpha$ (geometric) and BW-$\gamma$
(Tomita--Takesaki on the Krein-GNS Hilbert--Schmidt space
$M_{\dim W_L}(\mathbb{C})$) constructions both verify
$\sigma_{2\pi}^L(O) = O$ bit-exactly at every panel cell:

\begin{center}
\begin{tabular}{rrrcc}
\toprule
$(n_{\max}, N_t)$ & $\dim K$ & $\dim W_L$ & $r_W^{L, \alpha}$ & $r_W^{L, \mathrm{TT}}$ \\
\midrule
$(1, 1)$ & 4   & 2   & $0.0$                  & $\le 10^{-16}$ \\
$(1, 11)$ & 44  & 12  & $\le 10^{-16}$         & $\le 10^{-16}$ \\
$(1, 21)$ & 84  & 22  & $\le 2 \times 10^{-16}$ & $\le 2 \times 10^{-16}$ \\
$(2, 1)$  & 16  & 8   & $\le 10^{-16}$         & $\le 10^{-16}$ \\
$(2, 11)$ & 176 & 48  & $\le 3 \times 10^{-16}$ & $\le 3 \times 10^{-16}$ \\
$(2, 21)$ & 336 & 88  & $\le 4 \times 10^{-16}$ & $\le 4 \times 10^{-16}$ \\
$(3, 1)$  & 40  & 20  & $\le 5 \times 10^{-16}$ & $\le 5 \times 10^{-16}$ \\
$(3, 11)$ & 440 & 120 & $\le 4 \times 10^{-15}$ & $\le 4 \times 10^{-15}$ \\
$(3, 21)$ & 840 & 220 & $\le 7 \times 10^{-15}$ & $\le 7 \times 10^{-15}$ \\
\bottomrule
\end{tabular}
\end{center}

Verdict at every cell:\ \textsc{strong\_identification\_lorentzian}.
Residuals scale as $O(\sqrt{\dim} \cdot \varepsilon_{\mathrm{machine}})$
— pure round-off, no structural drift.  Flow conjugacy
$\sigma_t^{L, \mathrm{TT}}(O) = \sigma_{-t}^{L, \alpha}(O)$ at general
$t \in \{1, \pi, 2\pi\}$ holds bit-exactly with the same residual bound,
and the six-witness cross-consistency residual is exactly $0.0$ at
every cell.

\paragraph{Theorem (Krein-level four-witness Wick-rotation theorem at
finite cutoff).}\label{thm:krein_four_witness_paper32}
\emph{On the Lorentzian Krein space
$\mathcal{K}_{n_{\max}, N_t} = \mathcal{H}_{\mathrm{GV}}^{n_{\max}}
\otimes \mathbb{C}^{N_t}$, with the wedge
$W_L = P_W^{\mathrm{spatial}} \otimes P_{t \geq 0}$, the wedge KMS
state $\rho_W^L = e^{-K_L^{\alpha, W}}/Z$, and the BW unit normalisation
$\kappa_g = 1$, the four-witness Wick-rotation theorem
(Hartle--Hawking $+$ Sewell $+$ Bisognano--Wichmann $+$ Unruh) closes at
the operator-system level via:}
\begin{itemize}\setlength{\itemsep}{2pt}
\item \emph{BW-$\alpha$ period closure:} $\sigma_{2\pi}^{L, \alpha}(O) = O$
bit-exactly for all $O \in B(\mathcal{K}_{W_L})$.
\item \emph{BW-$\gamma$ Tomita period closure:} $\sigma_{2\pi}^{L, \mathrm{TT}}(a) = a$
bit-exactly for all $a \in B(\mathcal{K}_{W_L})$ on the Krein-GNS space.
\item \emph{Flow conjugacy:} $\sigma_t^{L, \mathrm{TT}}(a) = \sigma_{-t}^{L, \alpha}(a)$
bit-exactly at general $t$.
\item \emph{Six-witness collapse:} all six instantiations
$\{\mathrm{BW},\mathrm{HH}_{M=1},\mathrm{HH}_{M=2},\mathrm{Sew}_{M=1},
\mathrm{Unruh}_{a=1},\mathrm{Unruh}_{a=2}\}$ give bit-identical
$\Delta_L$ and $K_L^{\mathrm{TT}}$.
\end{itemize}
\emph{All four statements verified bit-exact at every tested
$(n_{\max}, N_t) \in \{1, 2, 3\} \times \{1, 11, 21\}$.  Proofs mirror
Paper~42 Theorems 5.4, 6.3, 7.1, and Cor.~8.1~\cite{paper42}.}

The structural reading is clean:\ the same M1 mechanism (Hopf-base
measure $\mathrm{Vol}(S^1)/2\pi$, master Mellin engine $k = 0$
sub-mechanism per the case-exhaustion theorem of \S~VIII) that
produces the Riemannian-side period closure on the Wightman-BW
vacuum analog $\rho_W$ produces the Lorentzian-side period closure on
$\rho_W^L$ at $(3, 1)$.  The Wick-rotation theorem (Hawking $+$ Sewell
$+$ BW $+$ Unruh) \textbf{lifts from structural correspondence at the
metric-functional level (Sprint TD~Track~4, Unruh-pendant) to literal
identification at the operator-system level (Lorentzian, finite
cutoff)} at every Krein cutoff.

\paragraph{Riemannian-limit recovery (load-bearing falsifier).}
At $N_t = 1$, $P_{W_L} = P_W^{\mathrm{spatial}}$,
$K_L^{\alpha, W} = K_\alpha^W$, $\rho_W^L = \rho_W$,
$\Delta_L = \Delta$, and $K_L^{\mathrm{TT}} = K_{\mathrm{TT}}$
bit-identically (Frobenius residual exactly $0.0$, not merely
machine-precision) at every $n_{\max} \in \{1, 2, 3\}$.  This is the
load-bearing falsifier that ensures Sprint L2-E is the Lorentzian
\emph{extension} of Paper~42~\S~6 (the Tomita BW-$\gamma$
construction)~\cite{paper42}, not a parallel-but-different construction.

\paragraph{$H_{\mathrm{local}}$ signature-independence finding
(headline structural finding of Sprint L2-E).}
Paper~42~\S~7.2 / open question O3 names a load-bearing scope finding:\
the framework's intrinsic Dirac $D_W$ is NOT the right local
Hamiltonian for the BW vacuum at $\beta = 2\pi$;\ $H_{\mathrm{local}} :=
K_\alpha^W / \beta$ is.  Sprint L2-E tests whether this finding holds
at signature $(3, 1)$.  \textbf{Result:} at the Riemannian limit
($N_t = 1$) on the Lorentzian Krein space,
$\|H_{\mathrm{local}}|_{(3, 1), N_t = 1} - D_L^W|_{N_t = 1}\|_F$
EQUALS the Riemannian-side residual bit-exact at every tested
$n_{\max} \in \{1, 2, 3\}$ (values $2.1332$, $6.5275$, $13.854$).
At $N_t > 1$ the Lorentzian-side residual is \emph{refined upward}
by the temporal-derivative content of $D_L = i (\gamma^0 \otimes
\partial_t + D_{\mathrm{GV}} \otimes I)$.  The
\textbf{Paper~42~\S~7.2 finding is signature-INDEPENDENT at the
Riemannian limit}.  The spectral-action--vs--modular-Hamiltonian
generator distinction is therefore NOT a peculiarity of the round
$S^3$ Riemannian truncation;\ it is a deeper structural feature of
the framework's modular content.  Paper~42~\S~10 O3 is refined from
``open question about signature-dependence'' to ``signature-independent
structural distinction at the Riemannian limit, refined upward by
temporal-derivative content at $N_t > 1$''~\cite{paper42}.

\paragraph{Honest scope.}
Closure is at finite $n_{\max}$ and finite $N_t$, NOT in the
continuum (GH-limit / Lorentzian propinquity) limit.  No published
Lorentzian propinquity construction exists as of May~2026
(Latr\'emoli\`ere 2017/2026 and Hekkelman--McDonald 2024 are strictly
Riemannian).  Constructing a Lorentzian propinquity is the named
Sprint L3 target.  The Sprint L2-D structural finding (BBB universal
$\chi D = -D\chi$ fails on truthful $D_{\mathrm{GV}}$) is preserved
in L2-E:\ the wedge construction is $K_L^{\alpha, W}$-driven and does
not invoke this axiom.  Cross-manifold W2b
($\Tcal_{S^3} \otimes \Tcal_{\mathrm{Hardy}(S^5)}$) remains
\textbf{open} as a separate frontier (see \S~VIII.D
``frontier-of-field framing''), structurally distinct from the
$(3, 1)$ extension within $\Tcal_{S^3}$ closed here.

\paragraph{Sources.}
Memo:\ \verb|debug/l2_e_modular_hamiltonian_lorentzian_memo.md|.
Computational driver:\
\verb|debug/l2_e_modular_hamiltonian_lorentzian_compute.py|.
Raw data:\
\verb|debug/data/l2_e_modular_hamiltonian_lorentzian.json|.
Implementation:\
\verb|geovac/modular_hamiltonian_lorentzian.py|.
Tests:\
\verb|tests/test_modular_hamiltonian_lorentzian.py| (74 fast + 3 slow,
all pass; zero regression in 355 upstream tests).
The math.OA-style standalone writeup of the Lorentzian extension is
outlined in Paper~43 (\verb|papers/group1_operator_algebras/paper_43_lorentzian_extension_outline.md|,
companion document);\ when drafted, Paper~43 sits alongside Papers~38,
39, 40, 42 as the fourth math.OA paper in the series and does not
supersede Paper~42 — Paper~42 is the Riemannian closure, Paper~43 is
the Lorentzian extension.

% -------------------------------------------------------------------
\subsection{Modular Hamiltonian at finite $n_{\max}$: Sprint L1 closure
on the Riemannian side (2026-05-16)}
\label{sec:modular_hamiltonian_l1}
% -------------------------------------------------------------------

The Sprint L1 falsifier named in
Remark~\ref{rem:bisognano_wichmann_reading} and
\S~\ref{sec:lorentzian_audit_paper32} is closed in this subsection at
the strongest possible level:\ the modular Hamiltonian $K$ on the
truncated Camporesi--Higuchi spectral triple $\mathcal{T}_{n_{\max}}$
satisfies the period closure $\sigma_{2\pi}(O) = O$ \emph{bit-exactly}
at finite $n_{\max}$ for all four physical witnesses
(Bisognano--Wichmann, Hartle--Hawking, Sewell, Unruh).  The result
lifts the four-witness Wick-rotation theorem from ``structural
correspondence'' (Remark~\ref{rem:bisognano_wichmann_reading} prior
verdict) to ``literal identification at the operator-system level
(Riemannian)'' at every finite cutoff.

\paragraph{Architecture.}
Sprint L1 implements the locked architectural decisions of L1-A
(\texttt{debug/l1\_modular\_hamiltonian\_architecture\_memo.md}), L1-B
(\texttt{debug/l1\_infrastructure\_audit\_memo.md}), and L1-C
(\texttt{debug/l1\_witness\_spec\_memo.md}):

\begin{itemize}
\item \textbf{Wedge (W1):}\ hemispheric $P_W$ on $S^3$ aligned with
the Hopf-base axis, defined as $P_W = (1/2)(I + R_{\mathrm{polar}})$
where $R_{\mathrm{polar}}: m_j \mapsto -m_j$ is the equatorial
reflection involution on the spinor basis;\ $P_W^2 = P_W$ and
$P_W^* = P_W$ are bit-exact (no smoothing needed; the reflection is
an exact involution on the discrete basis).
\item \textbf{Unit normalisation (U-1):}\ natural rapidity units
$\kappa_g = 1$ for the canonical BW witness;\ the four witnesses
parameterize as Hartle--Hawking $\kappa_g = 1/(4M)$, Unruh
$\kappa_g = a$, Sewell $\kappa_g = 1/(4M)$.
\item \textbf{Construction:}\ the BW-$\alpha$ geometric realization
$K = J_{\mathrm{polar}}$ (rotation generator around the wedge axis),
with integer eigenvalues $\mathrm{two\_m}_j$ on the full-Dirac basis
(odd integers in $\{-(2l+1), \ldots, +(2l+1)\}$).  The Tomita BW-$\gamma$
realization is also instantiated as a sanity check on the trace
cyclicity of the KMS state at expectation level; both realizations
close consistently.
\item \textbf{Dirac:}\ truthful Camporesi--Higuchi
(\verb|camporesi_higuchi_full_dirac_matrix|), the only Dirac choice
satisfying $J_{GV} D = +D J_{GV}$ on the truncated operator system
(see \S~\ref{sec:axiom_audit}).  R3.2's $n$-degeneracy obstruction for
the Connes-distance SDP does \emph{not} bite for modular flow, which
requires only cyclic-separatingness of the cyclic vector.
\end{itemize}

\paragraph{Result.}
The maximum periodicity residual
$\max_O \|\sigma_{2\pi}(O) - O\|_F$, taken over the wedge-restricted
versions of the first five non-identity multipliers in the
\verb|FullDiracTruncatedOperatorSystem| at cutoff $n_{\max}$, is at
machine precision $O(\sqrt{\dim \mathcal{H}_{n_{\max}}} \cdot
\varepsilon_{\mathrm{machine}})$ for every tested $n_{\max}$ and every
witness:

\begin{center}
\begin{tabular}{rrrcc}
\toprule
$n_{\max}$ & $\dim_H$ & $\mathrm{wedge\,dim}$ & $\max\|\sigma_{2\pi}(O) - O\|_F$ & verdict\\
\midrule
2 & 16 & 8 & $1.80 \times 10^{-16}$ & STRONG\_IDENTIFICATION\\
3 & 40 & 20 & $4.76 \times 10^{-16}$ & STRONG\_IDENTIFICATION\\
4 & 80 & 40 & $9.33 \times 10^{-16}$ & STRONG\_IDENTIFICATION\\
5 & 140 & 70 & $1.59 \times 10^{-15}$ & STRONG\_IDENTIFICATION\\
\bottomrule
\end{tabular}
\end{center}

The $n_{\max} = 3$ cutoff is structurally distinguished:\ $\dim_H = 40
= g_3^{\mathrm{Dirac}} = \Delta^{-1}$ (Paper~2's structurally
distinguished selection-principle cutoff).  The L1 closure at this
cutoff is the operator-system-level confirmation that the M1 mechanism's
$2\pi$ closure holds on the Paper-2-relevant truncation.

\paragraph{Cross-witness collapse.}
The cross-witness consistency residual
$\bigl|\|\sigma_{2\pi}^{(\text{witness})}(O) - O\| -
\|\sigma_{2\pi}^{\mathrm{BW}}(O) - O\|\bigr|$ is exactly $0$ at every
tested $n_{\max}$ between BW ($\kappa_g = 1$), HH at $M = 1$ ($\kappa_g
= 1/4$), HH at $M = 2$ ($\kappa_g = 1/8$), Sewell at $M = 1$
($\kappa_g = 1/4$), Unruh at $a = 1$ ($\kappa_g = 1$), and Unruh at
$a = 2$ ($\kappa_g = 2$).  The period closure is independent of
$\kappa_g$ because the geometric K\_boost has integer spectrum:
$\mathrm{e}^{\mathrm{i}\cdot 2\pi \cdot n} = 1$ for any integer $n$.
The four-witness theorem collapses to a single operator-system test
at the unit normalisation U-1.

\paragraph{Structural reading.}
The period closure $\sigma_{2\pi}(O) = O$ is finite-cutoff exact, not
a limit-statement.  It does not require the Gromov--Hausdorff machinery
of Theorem~\ref{thm:gh_convergence} (although it is consistent with it
via the propinquity rate cross-check: residual / $\gamma_{n_{\max}}$
ratio is $O(\varepsilon_{\mathrm{machine}})$, several orders below the
qualitative-rate prediction).  The structural reason is the
spectral-spectrum property:\ the BW-$\alpha$ geometric K has integer
spectrum on the full-Dirac basis (odd integers $\mathrm{two\_m}_j$),
inherited from the SO(4) Casimir grading of the Camporesi--Higuchi
spinor harmonics.  This is the framework-internal version of
Bisognano--Wichmann's continuum statement.

The $\pi$-source case-exhaustion theorem
(Theorem~\ref{thm:pi_source_case_exhaustion}) is enriched, not
weakened, by this result.  The $2\pi$ in the modular period is M1
content (Hopf-base measure $\mathrm{Vol}(S^1)$) at the operator-system
level, not just at the metric-functional level.  The dual physical
reading (Riemannian Hopf-base measure / Lorentzian boost-orbit period)
is now unified at the operator level inside the spectral triple, not
only by the published Wick-rotation chain.

\paragraph{Honest scope.}
The L1 closure is at signature $(3, 0)$ (Riemannian).  Three
distinctions remain open:

\begin{itemize}
\item The Lorentzian extension to signature $(3, 1)$ via the
Bizi--Brouder--Besnard~2018~\cite{bizi_brouder_besnard2018}
Krein-space spectral triple is the named Sprint L2 follow-up
(6--12 month scale per Sprint L0 audit
\texttt{debug/lorentzian\_l0\_audit\_memo.md}).  Sprint L0 predicts
M3 trivializes at $(3, 1)$ but M1 stays (so the four-witness theorem
extends to Lorentzian);\ Sprint L2 would verify this directly.
\item Non-tracial states would require explicit Tomita polar
decomposition of the modular $S$.  The tracial Gibbs state on
$\mathcal{T}_{n_{\max}}$ gives the trivial reduction
$K_{\mathrm{mod}} = \beta H_{\mathrm{local}} + \log Z$ (BW-$\gamma$);
non-tracial extension is an open structural question.
\item Cross-manifold modular structure (the Paper~24~\S~V cross-manifold
$\mathcal{T}_{S^3} \otimes \mathcal{T}_{\mathrm{Hardy}(S^5)}$ blocker)
remains unaddressed and is not the subject of Sprint L1.
\end{itemize}

The implementation is at
\verb|geovac/modular_hamiltonian.py| (production module, fresh) with
\verb|tests/test_modular_hamiltonian.py| (41 tests, 39 fast + 2 slow,
all pass).  Computational driver:
\verb|debug/l1_modular_hamiltonian_compute.py|.  Machine-readable
results JSON:\
\verb|debug/data/l1_modular_hamiltonian_results.json|.
Closure memo:\ \verb|debug/l1_modular_hamiltonian_results_memo.md|.

\paragraph{Sprint L1-tighten (same day):\ load-bearing Tomita--Takesaki
closure.}
The L1 closure above uses the BW-$\alpha$ geometric realization
$K = J_{\mathrm{polar}}$ as the load-bearing object; the Tomita BW-$\gamma$
realization was instantiated only as an expectation-level KMS sanity
check.  Sprint L1-tighten (same day, 2026-05-16 evening) rebuilds the
BW-$\gamma$ construction at the full Tomita--Takesaki level on the
GNS Hilbert--Schmidt space $H_{\mathrm{GNS}} = M_{\dim_W}(\mathbb{C})
\cong \mathbb{C}^{\dim_W^2}$, computes the modular operator $\Delta$
as a $(\dim_W^2 \times \dim_W^2)$ matrix via the polar decomposition
$S = J_{\mathrm{TT}} \Delta^{1/2}$, extracts $K_{\mathrm{TT}} = -\log
\Delta$, and verifies $\sigma_{2\pi}^{\mathrm{TT}}(O) = O$ at the
operator-action level for the same five generators tested by BW-$\alpha$.

The construction uses the natural local Hamiltonian $H_{\mathrm{local}}
:= K_{\alpha}^W / \beta$ on the wedge sub-algebra (so $\beta
H_{\mathrm{local}} = K_{\alpha}^W$ is the wedge-restricted boost-class
generator with integer spectrum $\mathrm{two\_m}_j$); this matches the
operator-system analog of the Wightman BW vacuum $\rho_W \propto
e^{-2\pi K_{\mathrm{boost}}}$ on the Rindler wedge.

\emph{Result (UNIFIED\_STRONG):}\ at every tested $n_{\max} \in \{2, 3,
4, 5\}$ for all six witness instantiations (BW, HH$_{M=1}$, HH$_{M=2}$,
Sewell$_{M=1}$, Unruh$_{a=1}$, Unruh$_{a=2}$):

\begin{center}
\begin{tabular}{rrrcccc}
\toprule
$n_{\max}$ & $\dim_W$ & $\dim_{\mathrm{GNS}}$
& $\sigma^{\alpha}_{2\pi}$ residual & $\sigma^{\mathrm{TT}}_{2\pi}$ residual
& $\sigma_t^{\mathrm{TT}} - \sigma_{-t}^{\alpha}$ at $t=1$ & $J_{\mathrm{TT}}^2 - I$\\
\midrule
2 & 8  & 64   & $1.10 \times 10^{-16}$ & $1.09 \times 10^{-16}$ & $4.18 \times 10^{-17}$ & $0$\\
3 & 20 & 400  & $2.22 \times 10^{-16}$ & $1.07 \times 10^{-15}$ & $9.79 \times 10^{-17}$ & $5.00 \times 10^{-17}$\\
4 & 40 & 1600 & $3.51 \times 10^{-16}$ & $3.03 \times 10^{-15}$ & $3.18 \times 10^{-16}$ & $7.22 \times 10^{-17}$\\
5 & 70 & 4900 & $4.97 \times 10^{-16}$ & $4.79 \times 10^{-15}$ & $4.03 \times 10^{-16}$ & $2.18 \times 10^{-17}$\\
\bottomrule
\end{tabular}
\end{center}

All residuals at machine precision.  The four-witness cross-collapse is
preserved at the Tomita--Takesaki level (in fact, the density matrix
$\rho_W = e^{-K_\alpha^W}/Z$ is itself $\beta$-independent, so $\Delta$
and $K_{\mathrm{TT}}$ are identical across the six witnesses);\ all
six instantiations give bit-identical $\sigma^{\mathrm{TT}}_{2\pi}$
residuals at every $n_{\max}$.

\emph{Structural relation between BW-$\alpha$ and BW-$\gamma$.}\ The two
constructions generate the same modular automorphism on the algebra,
up to the sign of $t$:
\begin{equation}
\sigma_t^{\mathrm{TT}}(a) \;=\; \rho^{it}\, a\, \rho^{-it}
\;=\; e^{-it K_\alpha^W}\, a\, e^{+it K_\alpha^W}
\;=\; \sigma_{-t}^{\alpha}(a).
\label{eq:tt_alpha_conjugate}
\end{equation}
The cross-flow difference $\|\sigma_1^{\mathrm{TT}}(O) -
\sigma_1^{\alpha}(O)\|_F \sim 0.16$ at non-period $t = 1$ is non-zero
in general (the conjugate flows are not equal at arbitrary $t$); but
$\|\sigma_1^{\mathrm{TT}}(O) - \sigma_{-1}^{\alpha}(O)\|_F < 10^{-16}$
verifies equation~\eqref{eq:tt_alpha_conjugate} bit-exactly.  At $t =
2\pi$ the sign is invisible because $e^{\pm i 2\pi n} = 1$ for any
integer $n$, so both flows close to the identity bit-exactly.

\emph{$J_{\mathrm{TT}}^2 = +I$ at the FULL Tomita level.}\ The
polar-decomposition antilinear conjugation $J_{\mathrm{TT}}(X) =
\rho^{1/2} X^* \rho^{-1/2}$ satisfies $J_{\mathrm{TT}}^2 = +I$ to
machine precision at every $n_{\max}$ (residuals in the table above).
This is structurally distinct from the L1 verification, which only
checked the canonical sanity case $U = I$; L1-tighten verifies the
$\rho^{1/2}$-weighted Tomita conjugation itself, which is the genuine
state-dependent antilinear operator defined by the polar decomposition.
This separates $J_{\mathrm{TT}}$ from the Connes KO-dim 3 charge
conjugation $J_{\mathrm{GV}}$ (which has $J_{\mathrm{GV}}^2 = -I$);
the two antilinear operators coexist on $H_{n_{\max}}$ without
conflation, as required by L1-A~\S~10e.

\emph{Verdict graduation.}\ The L1 closure is at the operator-system
level via the geometric BW-$\alpha$ realization \emph{and} via the
Tomita--Takesaki BW-$\gamma$ polar decomposition.  The two
constructions agree at the operator-action level (up to sign of $t$;
identical at the period).  The four-witness Wick-rotation theorem is
operator-system-level closed at signature $(3, 0)$ via both readings,
not just via the geometric one.  The L1~\S~7.3 ``BW-$\alpha$-only''
caveat is removed.

\emph{Honest scope.}\ The unified closure rests on the choice
$H_{\mathrm{local}} = K_{\alpha}^W / \beta$, which makes the wedge KMS
state the operator-system analog of the Wightman BW vacuum $\rho_W
\propto e^{-2\pi K_{\mathrm{boost}}}$.  An alternative choice
$H_{\mathrm{local}} = D_W$ (wedge-restricted Camporesi--Higuchi Dirac)
would not have integer spectrum at $\beta H_{\mathrm{local}}$ and the
bit-exact closure would not hold;\ in that case the framework's intrinsic
Dirac is simply not the right local Hamiltonian for the wedge KMS state
at $\beta = 2\pi$.  The L1-tighten unified closure is for the
BW-aligned wedge state, the load-bearing choice for the Bisognano--Wichmann
identification.

The L1-tighten implementation extends \verb|geovac/modular_hamiltonian.py|
with the \verb|TomitaModularStructure| class (load-bearing BW-$\gamma$)
and accessor methods (\verb|tomita_structure|, \verb|k_alpha_geometric|,
\verb|k_tomita|, \verb|compare_alpha_vs_tomita|) on the
\verb|ModularHamiltonian| class.  Test coverage:\
\verb|tests/test_modular_hamiltonian.py| extended with 26 fast + 2 slow
new tests (total 67 tests, 0 regressions of the 41 L1 tests).
Computational driver:\ \verb|debug/l1_tighten_tomita_compute.py|.
Machine-readable results:\
\verb|debug/data/l1_tighten_tomita_results.json|.  Closure memo:\
\verb|debug/l1_tighten_tomita_results_memo.md|.

% ===================================================================
\section{Marcolli--van~Suijlekom Lineage}
\label{sec:lineage}
% ===================================================================

The construction of Sec.~\ref{sec:construction} is structurally a
specialization of the gauge networks framework of Marcolli and
van~Suijlekom~\cite{marcolli_vs2014}, with the
Perez-Sanchez correction~\cite{perez_sanchez2025}
that the continuum limit is Yang--Mills without Higgs.  The key
parallels:

\begin{itemize}\setlength{\itemsep}{2pt}
\item \emph{Vertex algebra.}  Marcolli--vS take $\mathcal{A}_V = C(V)$
on a vertex set $V$; we take $\mathcal{A}_{\mathrm{GV}} = \mathbb{C}^{V_{\mathrm{Fock}}}$
on the Fock-projected $S^3$ vertex set (Sec.~\ref{sec:construction}).
\item \emph{Edge connection.}  Marcolli--vS take a connection on
edges valued in a finite group~$G$; Paper~25 takes $G = U(1)$ and
Paper~30 takes $G = \mathrm{SU}(2)$.
\item \emph{Spectral action.}  Marcolli--vS show that the spectral
action of the resulting almost-commutative spectral triple equals a
Wilson lattice gauge action on the connection; Paper~30 verifies this
explicitly for $\mathrm{SU}(2)$ at $n_{\max} = 3, 4$.
\item \emph{Continuum limit.}  Per Perez-Sanchez, the continuum limit
of the gauge-network spectral action is Yang--Mills without a Higgs
sector; Paper~28's heat-kernel Seeley--DeWitt expansion is the
spectral-action expansion in this limit.
\end{itemize}

\begin{observation}[Lineage placement]
\label{obs:lineage}
Each structural ingredient of the GeoVac spectral triple has published
precedent: finite almost-commutative spectral triples
(\cite{krajewski1998,paschke_sitarz2000,cacic2009}), graph spectral
triples (\cite{marcolli_vs2014,perez_sanchez2024}), spectral
truncations (\cite{connes_vs2021,hekkelman2022,hekkelman_mcdonald2024}), and
the Standard Model spectral action (\cite{chamseddine_connes2010}).
What is not duplicated in the literature is the $\alpha$ prediction at
$8.8 \times 10^{-8}$ from a discrete spectral construction; that
remains the original empirical claim of Paper~2.
\end{observation}

This is what the working hypothesis WH1 in the project register
upgrades to ``MODERATE--STRONG'' on (CLAUDE.md~\S 1.7): the structural
framing is in the published lineage; the predictive claim is not
duplicated.

% ===================================================================
\section{Open Questions}
\label{sec:open}
% ===================================================================

The construction of Sec.~\ref{sec:construction} and the audit of
Sec.~\ref{sec:axiom_audit} leave the following structural questions
genuinely open.

\paragraph{Q1.  Order one beyond rank 1.}  The order-one condition
holds trivially for the abelian $\mathcal{A}_{\mathrm{GV}}$ and
non-trivially constrains the rank-1 non-abelian extension
$\mathcal{A}_{\mathrm{GV}} \otimes M_2(\mathbb{C})$ of Paper~30.  At
rank~$\ge 2$ (e.g., $M_3(\mathbb{C})$ or sums of matrix algebras as in
the Standard Model triple), order-one becomes a strong constraint and
typically requires modification of the Dirac operator
($\textit{first-order condition} \Rightarrow$ Higgs sector
in~\cite{chamseddine_connes2010}).  Whether GeoVac admits a natural
rank-$\ge 2$ extension---and what its spectral action would predict---is
an open question.  The Bargmann--Segal $S^5$ HO triple of
Paper~24 was probed for an $\mathrm{SU}(3)$ extension in
Sprint~5 (CLAUDE.md~\S 2) with negative outcome; that is one negative
data point, not the full picture.  The electroweak slice
$\mathcal{A}_F = \mathbb{C} \oplus \mathbb{H}$ (rank up to 2 via
$\mathbb{H} \cong M_2(\mathbb{C})$ over $\mathbb{C}$, a real
$*$-subalgebra) is the most reachable test, and Sprint~H1
addresses it; see Sec.~\ref{sec:higgs_h1} below for the verdict.

\paragraph{Q2.  Real structure at finite $n_{\max}$ vs in the limit.}
Proposition~\ref{prop:reality} establishes a real structure on the
continuum spinor bundle that restricts to the truncation.  Whether the
restricted $J_{\mathrm{GV}}^{(n_{\max})}$ exactly satisfies $J^2 =
+\mathbb{1}$ at finite $n_{\max}$ in finite-precision arithmetic, or
only in the continuum limit, is a structural question that the audit
of Sec.~\ref{sec:axiom_audit} did not pin down to machine precision.
A sprint-scale verification at $n_{\max} \le 5$ would close this.

\paragraph{Q3.  KO-dimension and the four-way coincidence.}  Working
hypothesis WH4 asks: is the four-way $S^3$ coincidence (Sec.~\ref{sec:intro})
forced by the $\mathrm{KO}$-dimension~$3$ structure of the present
triple, or is it independent geometric data?  The relevant test is
whether any other rank-1 spectral triple of $\mathrm{KO}$-dimension~3
fails to produce all four roles.  $S^3$ is unique among compact rank-1
non-abelian simply-connected groups, so the question reduces to:
within the Marcolli--vS gauge-network setting, does
$\mathrm{KO}$-dimension~3 + non-abelian fiber + non-trivial Hopf
bundle uniquely select $S^3$?  The lineage literature does not
appear to address this directly.

\paragraph{Q4.  Spectral-action selection of $\Lambda$.}  Paper~28's
spectral-action computation produced an exact two-term form
$\mathrm{Tr}\,e^{-D^2/\Lambda^2} = (\sqrt{\pi}/2)\Lambda^3 -
(\sqrt{\pi}/4)\Lambda + O(e^{-\pi^2/\Lambda^2})$ but did not
autonomously select $\Lambda$.  Setting the trace equal to $K/\pi$
gave a closed-form Cardano solution $\Lambda_\infty = 3.7102\ldots$
(50 dps), which was tautological in the sense that it was determined
by demanding the equality.  Whether there is a structural principle
that selects $\Lambda$ is open; if so, $K/\pi$ would become a derived
quantity rather than a coincidence.  This is the most direct route by
which Paper~2 could be promoted from observation to derivation, and it
is also the one most resistant to incremental progress.

\paragraph{Q5.  Higher-form structure (Breit / rank-2 connections).}
Paper~25~\S VII.B asked whether the Breit interaction (rank-2 in
electron-electron coupling) lifts the $U(1)$ 1-form picture to a
higher-form gauge theory.  In spectral-triple language this is the
question whether the GeoVac triple admits a natural extension to a
rank-2 fiber that produces the Breit retarded integrals.  Paper~30's
$\mathrm{SU}(2)$ result is rank~1 (vector); rank~2 (tensor) is a
distinct extension; the answer is unknown.

% ===================================================================
\section{Conclusion}
\label{sec:conclusion}
% ===================================================================

This paper has constructed an explicit finite spectral triple
$(\mathcal{A}_{\mathrm{GV}}, \mathcal{H}_{\mathrm{GV}},
D_{\mathrm{GV}})$ on the Fock-projected $S^3$ graph, audited it
against the standard Connes axioms (Table~\ref{tab:axiom_audit},
Theorem~\ref{thm:axiom_complete}), and identified Papers~25, 28, 30,
31 as projections of the resulting object
(Table~\ref{tab:subsectors}).  The construction sits inside the
published Marcolli--van~Suijlekom gauge-network framework with the
Perez-Sanchez correction.  The combination rule
$K = \pi(B + F - \Delta)$ is now phrased as a sum across three
structurally distinct sectors of the same triple
(Observation~\ref{obs:three_sector}), without changing its
conjectural status.  The Coulomb/HO asymmetry of Paper~31 is the
statement that the GeoVac triple is one specific real odd
$\mathrm{KO}$-dimension-3 triple, not a generic class
(Observation~\ref{obs:pi_coulomb}).

Three things this paper did not do.  It did not derive $\alpha$;
Paper~2 stays in Observations.  It did not claim the Standard Model;
GeoVac is rank~1, the SM is higher rank.  It did not prove that the
four-way $S^3$ coincidence is forced by $\mathrm{KO}$-dimension~3;
that is open question Q3.

Three things this paper did do.  It made the spectral-triple framing
explicit and auditable.  It placed the framework in a published
mathematical lineage.  It made the open questions about $\alpha$, the
fiber rank, and the four-way coincidence into precise spectral-triple
questions instead of free-floating conjectures.

% ===================================================================
\begin{thebibliography}{99}

\bibitem{connes1995}
A.~Connes, ``Noncommutative geometry and reality,''
\textit{J.~Math.~Phys.}\ \textbf{36}, 6194--6231 (1995).

\bibitem{connes_marcolli2008}
A.~Connes and M.~Marcolli, \textit{Noncommutative Geometry, Quantum
Fields and Motives}, AMS Colloquium Publications~\textbf{55} (2008).

\bibitem{vansuijlekom_book2015}
W.~D.~van~Suijlekom, \textit{Noncommutative Geometry and Particle
Physics}, Mathematical Physics Studies, Springer (2015).

\bibitem{chamseddine_connes2010}
A.~H.~Chamseddine and A.~Connes,
``Spectral action principle,''
\textit{Commun.~Math.~Phys.}\ \textbf{186}, 731--750 (1997);
``Why the Standard Model,''
\textit{J.~Geom.~Phys.}\ \textbf{58}, 38--47 (2008);
``The spectral action principle in noncommutative geometry and the
superconnection,''
\textit{Phys.~Rev.~D}\ \textbf{83}, 045001 (2011).

\bibitem{marcolli_vs2014}
M.~Marcolli and W.~D.~van~Suijlekom,
``Gauge networks in noncommutative geometry,''
\textit{J.~Geom.~Phys.}\ \textbf{75}, 71--91 (2014),
arXiv:1301.3480.

\bibitem{fathizadeh_marcolli2016}
F.~Fathizadeh and M.~Marcolli,
``Periods and motives in the spectral action of Robertson--Walker
spacetimes,''
\textit{Comm.~Math.~Phys.}\ \textbf{356}, 641--671 (2017),
arXiv:1611.01815.

\bibitem{deligne2010}
P.~Deligne,
``Le groupe fondamental unipotent motivique de
$\mathbb{G}_m - \mu_N$, pour $N = 2, 3, 4, 6$ ou $8$,''
\textit{Publ.~Math.~IHES} \textbf{112}, 101--141 (2010);
preprint arXiv:math/0302267 (2005).

\bibitem{glanois2015}
C.~Glanois,
``Motivic unipotent fundamental groupoid of
$\mathbb{G}_m \setminus \mu_N$ for $N = 2, 3, 4, 6, 8$ and
Galois descents,''
arXiv:1411.4947 (v2, 2 Sep 2015);
published version \textit{J.~Number Theory} \textbf{182},
36--90 (2018).

\bibitem{perez_sanchez2024}
C.~I.~Perez-Sanchez,
``Bratteli networks and the Spectral Action on quivers,''
arXiv:2401.03705 (2024).

\bibitem{perez_sanchez2025}
C.~I.~Perez-Sanchez,
``Comment on `Gauge networks in noncommutative geometry','',
arXiv:2508.17338 (2025).

\bibitem{connes_vs2021}
A.~Connes and W.~D.~van~Suijlekom,
``Spectral truncations in noncommutative geometry and operator
systems,''
\textit{Commun.~Math.~Phys.}\ \textbf{383}, 2021--2067 (2021).

\bibitem{hekkelman2022}
E.~Hekkelman,
``Truncated geometries from spectral triples,''
PhD thesis, Radboud University (2022).

\bibitem{hekkelman_mcdonald2024}
E.-M.~Hekkelman and E.~A.~McDonald,
``A noncommutative integral on spectrally truncated spectral triples,
and a link with quantum ergodicity,''
\textit{J.~Funct.~Anal.}\ \textbf{289} (2025), accepted for
publication; arXiv:2412.00628.

\bibitem{vansuijlekom2024_ksystems}
W.~D.~van~Suijlekom,
``A generalization of $K$-theory to operator systems,''
arXiv:2409.02773 (2024).

\bibitem{latremoliere2026}
F.~Latr\'emoli\`ere,
``Spectral continuity of almost commutative manifolds for the
$C^1$ topology on Riemannian metrics,''
arXiv:2603.19128 (2026).

\bibitem{krajewski1998}
T.~Krajewski,
``Classification of finite spectral triples,''
\textit{J.~Geom.~Phys.}\ \textbf{28}, 1--30 (1998).

\bibitem{paschke_sitarz2000}
M.~Paschke and A.~Sitarz,
``Discrete spectral triples and their symmetries,''
\textit{J.~Math.~Phys.}\ \textbf{39}, 6191--6205 (1998).

\bibitem{cacic2009}
B.~\'Ca\'ci\'c,
``Real structures on almost-commutative spectral triples,''
\textit{Lett.~Math.~Phys.}\ \textbf{103}, 793--816 (2013), arXiv:1209.4832.

\bibitem{camporesi_higuchi1996}
R.~Camporesi and A.~Higuchi,
``On the eigenfunctions of the Dirac operator on spheres and real
hyperbolic spaces,''
\textit{J.~Geom.~Phys.}\ \textbf{20}, 1--18 (1996).

\bibitem{gilkey1995}
P.~B.~Gilkey,
\textit{Invariance Theory, the Heat Equation, and the Atiyah--Singer
Index Theorem}, 2nd ed., CRC Press, Boca Raton, 1995.

\bibitem{vassilevich2003}
D.~V.~Vassilevich,
``Heat kernel expansion: user's manual,''
\textit{Phys.~Rep.}\ \textbf{388}, 279--360 (2003), arXiv:hep-th/0306138.

\bibitem{karamata1962}
J.~Karamata,
\textit{Theory and Applications of Stieltjes Integrals},
SANU Posebna izdanja CCCXLIII, Belgrade, 1962.

\bibitem{korevaar2004}
J.~Korevaar,
\textit{Tauberian Theory:\ A Century of Developments},
Grundlehren der mathematischen Wissenschaften \textbf{329},
Springer, Berlin, 2004.

\bibitem{tenenbaum2015}
G.~Tenenbaum,
\textit{Introduction to Analytic and Probabilistic Number Theory},
3rd ed., Graduate Studies in Mathematics \textbf{163},
American Mathematical Society, Providence, RI, 2015.

\bibitem{paper0}
GeoVac Project, ``Geometric packing and the origin of quantum
numbers,'' Paper~0 (2025).

\bibitem{paper2}
GeoVac Project, ``The fine structure constant from Hopf bundle
spectral invariants,'' Paper~2 (2025; observation status, May~2026).

\bibitem{paper7}
GeoVac Project, ``The Dimensionless Vacuum: $S^3$ Conformal
Equivalence and 18 Symbolic Proofs,'' Paper~7 (2025).

\bibitem{paper14}
GeoVac Project, ``Structurally Sparse Qubit Hamiltonians from Natural
Geometry,'' Paper~14 (2025).

\bibitem{paper18}
GeoVac Project, ``Spectral-Geometric Exchange Constants,'' Paper~18
(2026).

\bibitem{paper22}
GeoVac Project, ``Potential-Independent Angular Sparsity Theorem,''
Paper~22 (2026).

\bibitem{paper23}
GeoVac Project, ``Nuclear Shell Model Qubit Hamiltonians and the
Fock Projection Rigidity Theorem,'' Paper~23 (2026).

\bibitem{paper24}
GeoVac Project, ``The Bargmann--Segal Lattice: $\pi$-Free
Discretization of the 3D Harmonic Oscillator on $S^5$,'' Paper~24
(2026).

\bibitem{paper25}
GeoVac Project, ``The Hopf Graph as a Lattice Gauge Structure: A
Framework Observation,'' Paper~25 (2026).

\bibitem{paper28}
GeoVac Project, ``QED on $S^3$: Spectral Geometry of Perturbative
Transcendentals,'' Paper~28 (2026).

\bibitem{paper30}
GeoVac Project, ``SU(2) Wilson Lattice Gauge Structure on the GeoVac
Hopf Graph: The Non-Abelian Sibling of Paper~25,'' Paper~30 (2026).

\bibitem{paper31}
GeoVac Project, ``Universal/Coulomb Partition: The $\mathcal{A}/D$
Split of the GeoVac Spectral Triple,'' Paper~31 (2026).

\bibitem{paper34}
GeoVac Project, ``Projection Taxonomy and Empirical Matches
Catalogue,'' Paper~34 (2026).

\bibitem{paper35}
GeoVac Project, ``Time as Projection:\ Klein--Gordon Spectrum on
$S^3 \times \mathbb{R}$ and the Temporal-Window Mechanism for $\pi$,''
Paper~35 (2026).

\bibitem{paper36}
GeoVac Project, ``Bound-State QED on the GeoVac Spectral Triple:\
Hydrogen Lamb Shift at One Loop,'' Paper~36 (2026; observation status).

\bibitem{paper39}
J.~Loutey, ``Latr\'emoli\`ere propinquity convergence of truncated
tensor-product Camporesi--Higuchi spectral triples on $S^3 \times S^3$,''
Paper~39 (2026; companion to Paper~38), DOI via Zenodo.

\bibitem{bisognano_wichmann1975}
J.~J.~Bisognano and E.~H.~Wichmann,
``On the duality condition for a Hermitian scalar field,''
\textit{J.~Math.~Phys.}\ \textbf{16}, 985--1007 (1975).

\bibitem{bisognano_wichmann1976}
J.~J.~Bisognano and E.~H.~Wichmann,
``On the duality condition for quantum fields,''
\textit{J.~Math.~Phys.}\ \textbf{17}, 303--321 (1976).

\bibitem{sewell1982}
G.~L.~Sewell,
``Quantum fields on manifolds: PCT and gravitationally induced thermal
states,''
\textit{Annals of Physics}\ \textbf{141}, 201--224 (1982).

\bibitem{hartle_hawking1976}
J.~B.~Hartle and S.~W.~Hawking,
``Path-integral derivation of black-hole radiance,''
\textit{Phys.~Rev.~D}\ \textbf{13}, 2188--2203 (1976).

\bibitem{unruh1976}
W.~G.~Unruh,
``Notes on black-hole evaporation,''
\textit{Phys.~Rev.~D}\ \textbf{14}, 870--892 (1976).

\bibitem{witten2018_rmp}
E.~Witten,
``APS Medal for Exceptional Achievement in Research: Invited article on
entanglement properties of quantum field theory,''
\textit{Rev.~Mod.~Phys.}\ \textbf{90}, 045003 (2018), arXiv:1803.04993.

\bibitem{bizi_brouder_besnard2018}
N.~Bizi, C.~Brouder, and F.~Besnard,
``Space and time dimensions of algebras with applications to
Lorentzian noncommutative geometry and quantum electrodynamics,''
\textit{J.~Math.~Phys.}\ \textbf{59}, 062303 (2018), arXiv:1611.07062.

\bibitem{vandungen2016}
K.~van~den~Dungen,
``Krein spectral triples and the fermionic action,''
\textit{Math.\ Phys.\ Anal.\ Geom.}\ \textbf{19}, 4 (2016),
arXiv:1505.01939.

\bibitem{nieuviarts2025a}
G.~Nieuviarts,
``Emergence of Lorentz symmetry from an almost-commutative twisted
spectral triple,''
arXiv:2502.18105 (2025).

\bibitem{paper42}
J.~Loutey, \emph{Tomita--Takesaki modular structure on truncated SU(2)
spectral triples: a four-witness Wick-rotation theorem at finite
cutoff} (Paper~42 of the GeoVac series, Zenodo, 2026).

\bibitem{paper43}
J.~Loutey, \emph{Lorentzian extension of the four-witness
Wick-rotation theorem at finite cutoff} (Paper~43 of the GeoVac
series, Zenodo, 2026).

\bibitem{latremoliere2021dual}
F.~Latr\'emoli\`ere, ``The Dual Modular Propinquity and Completeness,''
\textit{J.\ Noncommut.\ Geom.}\ \textbf{15}, 347 (2021).
arXiv:1811.04534.

\end{thebibliography}

\end{document}
